\pdfminorversion=7
\documentclass[11pt,a4paper]{article}
\usepackage[T1]{fontenc}
\usepackage[utf8]{inputenc}
\usepackage{lmodern,amsmath,amssymb,amsthm,mathtools}
\usepackage[margin=27mm,headheight=14pt]{geometry}
\usepackage{microtype,booktabs,longtable,tabularx,array,enumitem,xurl,seqsplit}
\usepackage{fancyhdr}
\usepackage[hidelinks]{hyperref}
\usepackage[nameinlink,noabbrev]{cleveref}
\newtheorem{theorem}{Theorem}[section]
\newtheorem{lemma}[theorem]{Lemma}
\newtheorem{proposition}[theorem]{Proposition}
\newtheorem{corollary}[theorem]{Corollary}
\theoremstyle{definition}
\newtheorem{definition}[theorem]{Definition}
\theoremstyle{remark}
\newtheorem{remark}[theorem]{Remark}
\newcommand{\Ht}{H_t}
\newcommand{\Lam}{\Lambda}
\newcommand{\Ree}{\operatorname{Re}}
\newcommand{\Imm}{\operatorname{Im}}
\newcommand{\dd}{\mathrm d}
\newcommand{\C}{\mathbb C}
\newcommand{\R}{\mathbb R}
\newcommand{\N}{\mathbb N}
\newcommand{\Q}{\mathbb Q}
\newcommand{\e}{\mathrm e}
\newcommand{\ii}{\mathrm i}
\DeclareMathOperator{\supp}{supp}
\DeclareMathOperator{\ord}{ord}
\DeclareMathOperator{\Log}{Log}
\setlength{\parindent}{1.2em}
\setlength{\parskip}{2pt plus 1pt}
\setlength{\emergencystretch}{3em}
\setlist{itemsep=2pt,topsep=4pt,leftmargin=1.7em}
\setcounter{tocdepth}{2}
\setcounter{secnumdepth}{3}
\pagestyle{fancy}
\fancyhf{}
\fancyhead[L]{\small A computer-assisted bound for $\Lambda$}
\fancyhead[R]{\small D. J. Raymond}
\fancyfoot[C]{\small\thepage}
\renewcommand{\headrulewidth}{0.3pt}
\hypersetup{pdftitle={A computer-assisted bound Lambda <= 0.107964910079 for the de Bruijn-Newman constant},pdfauthor={Dale James Raymond},pdfsubject={Version 1.0 with direct source reconstruction, 29 October 2026},pdfkeywords={de Bruijn-Newman constant, heat flow, translated probes, exact arithmetic, interval arithmetic}}
\title{A computer-assisted bound\\
 $\Lambda\le 0.107964910079$\\
 for the de Bruijn--Newman constant}
\author{Dale James Raymond\\[3pt]
 \small ORCID: \href{https://orcid.org/0009-0001-1682-8529}{0009-0001-1682-8529}}
\date{29 October 2026\quad\textbar\quad Version 1.0\\[3pt]
\small DOI: \href{https://doi.org/10.5281/zenodo.23752103}{10.5281/zenodo.23752103}}

\begin{document}
\maketitle
\begin{abstract}
In the Polymath normalisation, a computer-assisted comparison gives
\[
 \Lambda\le
 \frac{427693072077942473661578098491}
 {3961408125713216879677197516800}
 <0.107964910079.
\]
The initial squared zero height is bounded by $9/16$, using OpenAI's
ordinary-zeta $7/8$ nonvanishing theorem formalised in Lean at the
identified source revision. Finite signed translated-probe
functionals, matched physical source estimates, and continuous-domain
kernel inequalities supply a piecewise comparison for that height.
A newly implemented exact rational consumer checks every one of the
$221\,010$ arcs, their outward rounding and joins, and the complete
final de Bruijn contraction. The fifteen-stage finite replay includes a direct reconstruction of all
$67$ source chords from their primitive payments, with explicit binding
of the main terms and the complete derivative functionals. Separate
consumers verify their continuous kernel certificates and the actual
physical and arithmetic charges. The analytic argument includes the
corrected $11/12$ Gaussian remainder, finite-head localisation,
attainment of a positive-time highest contact, and repeated complex
contacts. The accompanying packet contains the mathematical sources,
certificate data, checker programs, exact source identities, and replay
instructions. The initial theorem is used at cited-theorem trust with
authenticated external Lean-build evidence; a fresh local Lean build
and an end-to-end formalisation of the present comparison are not claimed.
\end{abstract}
\noindent\textbf{2020 Mathematics Subject Classification.}
Primary 11M26; secondary 11Y35, 65G30.\\
\textbf{Keywords.} de Bruijn--Newman constant, zero motion, translated
probes, finite logarithmic source, rigorous computation.

\clearpage
\tableofcontents
\clearpage

% BEGIN INCLUDED FILE: sections/introduction.tex
\section{Result and proof structure}\label{intro:result}
Let
\begin{align}
 H_t(z)&=\int_0^\infty e^{tu^2}\Phi(u)\cos(zu)\,\dd u,
 \label{intro:H}\\
 \Phi(u)&=\sum_{n=1}^{\infty}
 (2\pi^2n^4e^{9u}-3\pi n^2e^{5u})e^{-\pi n^2e^{4u}}.
 \nonumber
\end{align}
The decay of $\Phi$ makes $H_t$ entire for every real $t$. These functions
are even and real on the real axis, and satisfy
\begin{equation}\label{intro:normalisation}
 H_0(z)=\tfrac18\xi((1+iz)/2),\qquad
 \partial_t H_t=-\partial_z^2H_t,\qquad
 \xi(s)=\tfrac12s(s-1)\pi^{-s/2}\Gamma(s/2)\zeta(s).
\end{equation}
The de Bruijn--Newman constant $\Lambda$ is the real threshold for which
all zeros of $H_t$ are real exactly when $t\ge\Lambda$
\cite{deBruijn,Newman,Polymath}.

\begin{theorem}\label{intro:theorem}
For the normalisation \eqref{intro:H}--\eqref{intro:normalisation},
\begin{equation}\label{intro:bound}
 0\le\Lambda\le B_*:=
 \frac{427693072077942473661578098491}
 {3961408125713216879677197516800}
 <\frac{107964910079}{10^{12}}.
\end{equation}
The upper bound is computer assisted and uses the ordinary-zeta
$7/8$ theorem \cite{OAI:The-Quasi-Riemann-Hypothesis-September-30-2026} and the finite-height zero theorem of
Platt and Trudgian \cite{PlattTrudgian}. The exact finite certificates
and analytic estimates used to obtain the upper bound are specified
below and supplied in the accompanying packet.
\end{theorem}

The lower inequality is the theorem of Rodgers and Tao
\cite{RodgersTao}. The task here is to establish the upper inequality.
The comparison does not use an assumption of the full Riemann hypothesis.
The ordinary $7/8$ input, including its source and external verification
boundary, is stated separately so that a referee can check the imported
theorem independently of the new comparison.

\subsection{Initial data and terminal data}
By \cite{OAI:The-Quasi-Riemann-Hypothesis-September-30-2026}, a nontrivial zeta zero $\rho$ satisfies
$\Ree\rho\le7/8$. The functional equation also gives
$\Ree\rho\ge1/8$. Under \eqref{intro:normalisation} its corresponding
$H_0$ zero therefore has $|\Imm z|\le3/4$. The strict squared-height
barrier starts at
\begin{equation}\label{intro:q0}
 Q_0=\frac9{16}+10^{-12}
     =\frac{562500000001}{10^{12}}.
\end{equation}
The strict cushion is retained in every rational computation. At the
terminal comparison time the certified data are
\begin{equation}\label{intro:terminal}
 T=\frac{22101}{204800}=0.1079150390625,
 \qquad
 Q_T=\frac{7902377963830536747158643}
 {79228162514264337593543950336}.
\end{equation}
Here $Q_T>0$. Consequently a further time $Q_T/2$ is required for
de Bruijn's strip contraction. Exact addition gives
$T+Q_T/2=B_*$; the displayed decimal in
\eqref{intro:bound} is rounded upward. Stopping at $T$ alone would
omit part of the proof.

\subsection{From a finite functional to a global height barrier}
At a simple nonreal highest zero $x+iy$, set $q=y^2$.
A finite signed combination of logarithmic-derivative probes supplies
a physical lower bound $J$; its entire distinguished conjugate-pair
contribution is the explicitly defined function $O_{\mathrm{pair}}(q)$. A kernel inequality then gives
\begin{equation}\label{intro:current}
 q'\le -2-2q\bigl(J-O_{\mathrm{pair}}(q)\bigr).
\end{equation}
The term $-2$ is the attraction to the distinguished conjugate zero.
The subtraction $O_{\mathrm{pair}}(q)$ belongs to the whole signed functional and is
paid exactly once. In particular, optimising a source floor alone
does not optimise the contraction rate.

The proof has three distinct analytic obligations before one can
integrate \eqref{intro:current}. First, the sources must apply to
the actual function $H_t$ at every declared probe and real coordinate,
including arbitrarily large exterior coordinates. Second, the signed
kernel inequalities must cover all real spacings and the full permitted
heights of every other zero, including nonreal zeros. Third, the
first-contact argument must handle compactness, time zero and repeated
complex zeros. These obligations are proved in the following sections;
the finite numerical histories check the subsequent scalar comparison.

\begin{center}
\small
\begin{tabularx}{\textwidth}{@{}lX@{}}\toprule
Proof component & Content and verification boundary\\\midrule
Initial strip & Ordinary-zeta $7/8$ theorem; exact strip map
and strict seed \eqref{intro:q0}.\\
Finite head & Published finite RH height, full continuous wall,
and analytic zero continuation.\\
Exterior source & Corrected height-six approximation and
signed finite, matched-source inequalities.\\
Kernel & Rational continuous-domain certificates, whole own pair,
and harmonic extension to other-zero heights.\\
Global comparison & Simple and repeated contacts, remote compactness,
all closed time cells and upward numerical resets.\\
Exact flow & $221\,010$ rational arcs, independent checker,
and full terminal $Q_T/2$ cleanup.\\\bottomrule
\end{tabularx}
\end{center}

\subsection{Edition and reproducibility}
This edition supersedes referee edition 3 of 28 October 2026; the numerical plan, source table and comparison history are unchanged; Appendix B adds the direct source reconstruction.

\subsection{Relation to earlier versions}
The pinned public repository \cite{RaymondPrevious} states the ceiling
$0.116397726408291$. The present displayed ceiling is smaller by
$0.008432816329291$. This compares the stated numerical bounds; independent
recognition of a world record is not asserted.
An earlier research comparison gave $0.107936602217$ under the narrower
initial assertion $\Ree\rho\le0.8748867$. Xu and Sang
\cite[Theorem~1.1]{XuSang2026} obtain that boundary under the analytic
inputs stated there. Liu's related refinements give
$0.874957019420098946\ldots$ \cite{LiuRefinement2026} and
$0.874886647039964\ldots$ \cite[Theorem~17.1]{LiuHecke2026}, with their
stated analytic inputs. Since $0.8748867<7/8$, the narrower assertion
does not follow from the ordinary $7/8$ theorem alone. The present
argument therefore checks a new history from the literal seed
\eqref{intro:q0}; its input remains the ordinary $7/8$ theorem.

The local native and confluent calculations retained in the research
archive are not used in Theorem~\ref{intro:theorem}. Nor is an unproved
uniform-current premise used. In particular, neither $\Lambda<0.1$ nor
the Riemann or Lindelof hypothesis is a conclusion of this manuscript.


\subsection{Related work and the external strip input}
The heat-flow and zero-motion framework follows de Bruijn, Newman and
Polymath~\cite{deBruijn,Newman,Polymath}. Ki, Kim and Lee
\cite{KiKimLee2009} proved that, at every positive heat parameter,
all but finitely many zeros are real and simple, and obtained
$\Lambda<1/2$. The finite real-zero input used here is the verification
of Platt and Trudgian~\cite{PlattTrudgian}; the nonnegative lower bound
is due to Rodgers and Tao~\cite{RodgersTao}.

Raymond's earlier time-zero construction and translated geometry are
set out in the analytic companion~\cite{RaymondAnalytic}. The subsequent
phase-correlated, translated-probe and retuned releases are
\cite{RaymondPhase,RaymondCorrelated,RaymondRetuned}. The physical
estimates and direct source reconstruction used here are stated in
the present manuscript.

Planat~\cite[Section~8]{PlanatCollective2026} reports a complete numerical
reproduction of Raymond's frozen correlated-probe release, recovering
$0.149395073827688$, and a head-pair extension to
$0.149344833997286$. His later classical-strip comparison
\cite[Sections~4--5]{PlanatClassical2026} reports a stored-source replay
of the retuning archive, including all eight supplied comparison
histories, and independently checks the $57$-kernel catalogue and
$19\,023$ steps of its refined comparison. These replications concern
the earlier releases and retain the analytic premises stated in the
cited theorems.

The 2026 entries in Tao's maintained catalogue~\cite{TaoCatalogue2026}
are the Mosaic Intelligence packages for $0.1965$, $0.1875$, $0.16$ and
$0.1595$~\cite{MI2026,MI2026b,MI2026c,MI2026d}, and Kornaus's
$0.172422$ package~\cite{Kornaus2026}. The catalogue marks these as
self-archived certificate packages awaiting external review. Further
computer-assisted preprints include Gomila's $0.1787854$
release~\cite{Gomila2026} and Gordon's bound below
$0.158$~\cite{Gordon2026}. These references provide comparison context.

Recent comparisons also require attention to the initial theorem.
Hakimoglu-Brown~\cite{HakimogluBrown2026} states the $0.13963$ ceiling
under the seven-eighths hypothesis. The height-dependent region of
Bellotti, Trudgian and Yang~\cite{BellottiTrudgianYang2026} is a different
input from a fixed global seven-eighths half-plane. The OpenAI preprint
\cite{OAI:The-Quasi-Riemann-Hypothesis-September-30-2026} supplies the
ordinary $7/8$ input used here; its pinned Lean declaration is identified
in Section~\ref{ver:formal}.

% END INCLUDED FILE


% BEGIN INCLUDED FILE: sections/foundations.tex
\section{Analytic foundations and the passage from source bounds to a global strip}
\label{fnd:section}

This section fixes the normalisation and proves the analytic implications
used to consume the finite certificates. The published inputs are the
strip-contraction theorem of de Bruijn, the local zero-dynamics and
positive-time tail theorems in \cite{Polymath}, and the finite Riemann-hypothesis
verification in \cite{PlattTrudgian}. The explicit wall at the numerical value
of $X$ below is a computation in this work. In particular, the published
finite-height theorem is not being asserted at every later heat time.

\subsection{Normalisation, symmetry, and the initial strip}
\label{fnd:normalisation}

Put
\begin{align}
 \xi(s)&=\frac{s(s-1)}2\pi^{-s/2}\Gamma(s/2)\zeta(s),\label{fnd:xi}\\
 \Phi(u)&=\sum_{n=1}^{\infty}
 \bigl(2\pi^2n^4e^{9u}-3\pi n^2e^{5u}\bigr)e^{-\pi n^2e^{4u}},\label{fnd:phi}\\
 H_t(z)&=\int_0^{\infty}e^{tu^2}\Phi(u)\cos(zu)\,\dd u.
 \label{fnd:heatdefinition}
\end{align}
The apparent singularities in \eqref{fnd:xi} are removed. Thus $\xi$ is
entire, and the standard Fourier identity gives
\begin{equation}
 H_0(z)=\frac18\xi\left(\frac{1+iz}{2}\right),
 \qquad \partial_tH_t=-\partial_z^2H_t.
 \label{fnd:heat}
\end{equation}
The super-exponential decay of $\Phi$ justifies differentiation under the
integral on compact subsets of the $(t,z)$ domain. In particular $H_t$ is
jointly analytic locally in $t$ and $z$, and, for real $t$, is an entire
function of order one. It is real on the real axis and even:
\begin{equation}
 H_t(\overline z)=\overline{H_t(z)},\qquad H_t(-z)=H_t(z).
 \label{fnd:symmetry}
\end{equation}
All zero counts below include multiplicity. Every summand of
\eqref{fnd:phi} is positive when $u\geq0$, since
$2\pi n^2e^{4u}-3>0$. Hence
\begin{equation}
 H_t(iy)=\int_0^{\infty}e^{tu^2}\Phi(u)\cosh(yu)\,\dd u>0
 \quad(t,y\in\mathbb R).
 \label{fnd:imaginarypositive}
\end{equation}
In particular $H_t$ is never the zero function and has no zeros on the
imaginary axis.

The transformation between zeta zeros and $H_0$ zeros is worth keeping
explicit. If $z=x+iy$, then
\begin{equation}
 \frac{1+iz}{2}=\frac{1-y}{2}+i\frac{x}{2}.
 \label{fnd:coordinates}
\end{equation}
Thus the classical critical strip gives $|y|\leq1$ at $t=0$.
The ordinary seven-eighths theorem used here is the nonvanishing assertion
\begin{equation}
 \zeta(s)\ne0\qquad\left(\Ree s>\frac78,\ s\ne1\right)
 \label{fnd:seveneighths}
\end{equation}
for the ordinary Riemann zeta function; the finite extension assigned to
the pole by the formal library is irrelevant to this nontrivial-zero
statement. Its precise source version and verification provenance are
specified with the external theorem input
\cite{OAI:The-Quasi-Riemann-Hypothesis-September-30-2026}. By the functional equation, a nontrivial zero with
real part below $1/8$ would give one with real part above $7/8$.
Consequently all nontrivial zeta zeros have real parts in $[1/8,7/8]$.
The factors defining $\xi$ add no zeros at $s=0$ or $s=1$, and its zeros
are precisely the nontrivial zeta zeros. Equation \eqref{fnd:coordinates}
therefore gives the global initial strip
\begin{equation}
 H_0(x+iy)=0\quad\Longrightarrow\quad |y|\leq\frac34,
 \qquad y^2\leq\frac9{16}.
 \label{fnd:initialstrip}
\end{equation}
The weaker seven-eighths input is used literally; a strip with boundary
$0.8748867$ is not inferred from it.

The following classical contraction theorem is used twice in the proof. It is \cite[Theorem~3.2]{Polymath}, quoting
\cite[Theorem~13]{deBruijn}.
\begin{theorem}[de Bruijn strip contraction]\label{fnd:debruijn}
Suppose $t_0\in\mathbb R$ and $a>0$, and every zero of $H_{t_0}$ has
imaginary part in $[-a,a]$. Then for $t\geq t_0$ every zero of $H_t$ has
imaginary part of modulus at most
\[
 \sqrt{\max\{a^2-2(t-t_0),0\}}.
\]
In particular all zeros are real at every time $t\geq t_0+a^2/2$.
\end{theorem}
The initial use of this theorem and \eqref{fnd:initialstrip} yields
\begin{equation}
 H_t(x+iy)=0\quad\Longrightarrow\quad
 y^2\leq\max\left\{\frac9{16}-2t,0\right\},\qquad t\geq0.
 \label{fnd:initialcontraction}
\end{equation}
Only a very short initial segment of this elementary comparison is needed
for strict startup. The quantitative improvement comes from the source
and geometry inequalities, not from extrapolating the contraction theorem.

\subsection{The stationary wall and finite-head theorem}
\label{fnd:finitewall}

The finite wall is at
\begin{equation}
 X=6000000185827,\qquad N=690988,\qquad T_{\mathrm w}=\frac{41}{250}.
 \label{fnd:wallparameters}
\end{equation}
It is stationary in the $z$ coordinate. Define
\[
 L=\log\frac{X}{4\pi},\qquad v_n=\log n,
\]
and use the leading gamma factor
\begin{equation}
 M_0(s)=\frac{s(s-1)}{16}\pi^{-s/2}\sqrt{2\pi}
            (s/2)^{s/2-1/2}e^{-s/2}.
 \label{fnd:Mzero}
\end{equation}
The powers use the logarithm analytic on the relevant upper or lower
half-plane. Set
\[
 \theta_0=\arg\frac{M_0((1+iX)/2)}{M_0((1-iX)/2)}.
\]
The certified branch satisfies $-11/10<\theta_0<-1$. The wall computation
also certifies
\begin{equation}
 \sqrt{\frac{X}{4\pi}}\geq N,\qquad
 \sqrt{\frac{X}{4\pi}+\frac{T_{\mathrm w}}{16}}<N+1,
 \qquad \log N<\frac L2.
 \label{fnd:cutoff}
\end{equation}
In particular the literal Riemann--Siegel cutoff is the same $N$ throughout
the wall interval.

\begin{lemma}[Continuous prefix-to-wall transfer]\label{fnd:abel}
For $0\leq j\leq3$, let
\[
 t_j=\frac{41(2j+1)}{2000},\qquad r=\frac{41}{2000}.
\]
Suppose that for every $1\leq J\leq N$ the following directed inequalities
hold:
\begin{multline}
 \sum_{n\leq J}\frac1{\sqrt n}
 \cos\left(\frac{Xv_n}{2}-\frac{\theta_0}{2}
            -\frac{\pi t_j(L-2v_n)}{16}\right)\\
 -r\sum_{n\leq J}\frac{\pi(L-2v_n)}{16\sqrt n}
 >\frac{79}{100}.
 \label{fnd:prefix}
\end{multline}
Then the real part of the unit-dephased ideal paired model at every
$(t,y)\in[0,T_{\mathrm w}]\times[0,1]$ is greater than $79/100$.
\end{lemma}
\begin{proof}
The four closed cells $[t_j-r,t_j+r]$ cover $[0,T_{\mathrm w}]$.
By $|\cos a-\cos b|\leq|a-b|$, \eqref{fnd:prefix} implies that every
partial sum of
\[
 c_n(t)=\frac1{\sqrt n}\cos\psi_n(t),\qquad
 \psi_n(t)=\frac{Xv_n}{2}-\frac{\theta_0}{2}
                  -\frac{\pi t(L-2v_n)}{16}
\]
is greater than $79/100$ throughout the corresponding closed cell.
The ideal dephased paired expression is
\begin{equation}
 F_{\mathrm{id}}(t,y)=\sum_{n=1}^N\frac{e^{-tv_n(L-v_n)/4}}{\sqrt n}
 \left(e^{-yv_n/2}e^{i\psi_n(t)}
       +e^{-y(L-v_n)/2}e^{-i\psi_n(t)}\right).
 \label{fnd:idealpair}
\end{equation}
Its real part is
\begin{equation}
 \Ree F_{\mathrm{id}}(t,y)
 =2e^{-yL/4}\sum_{n=1}^N c_n(t)\,w_{t,y}(v_n),
 \quad
 w_{t,y}(v)=e^{-tv(L-v)/4}\cosh\left(\frac{y(L-2v)}4\right).
 \label{fnd:idealreal}
\end{equation}
Both factors of $w_{t,y}$ are positive and nonincreasing on $[0,\log N]$.
For the first this follows from $L-2v>0$; for the second it follows from
$y\geq0$ and $L-2v\geq0$. Write $C_J=\sum_{n\leq J}c_n(t)$ and
$w_n=w_{t,y}(v_n)$. Finite Abel summation gives
\[
 \sum_{n=1}^N c_n(t)w_n
 =C_Nw_N+\sum_{n=1}^{N-1}C_n(w_n-w_{n+1})
 >\frac{79}{100}w_1.
\]
Since $v_1=0$, multiplication by $2e^{-yL/4}$ gives
\[
 \Ree F_{\mathrm{id}}(t,y)
 >\frac{79}{100}\,2e^{-yL/4}\cosh(yL/4)
 =\frac{79}{100}(1+e^{-yL/2})\geq\frac{79}{100}.
\]
No sampling in $t$ or $y$ appears in this deduction.
\end{proof}

The inequalities \eqref{fnd:prefix} comprise $4N=2763952$ prefix tests.
Their direct 512-bit replay checks every prefix and both cutoff guards,
with a total of $3454987$ gates including the ancillary wall checks.
A ball-arithmetic lower endpoint must exceed the displayed rational
threshold; overlap or a floating approximation to the threshold is not
an acceptance criterion. The physical transfer uses the same nonzero
normaliser $B_t$ as the Riemann--Siegel model in
Section~\ref{src:section}. More explicitly, unit dephasing by
$\exp[-i(\theta_0/2+\pi tL/16)]$ takes its leading pair to
\eqref{fnd:idealpair}; the difference from the physical normalised function
is charged as follows.

Write $v_N=\log N$, and bound the source-coordinate and gamma changes by
\begin{align}
 d_\alpha&=5/X,& A_\alpha&=L/2+1,\\
 \epsilon_s&=T_{\mathrm w}d_\alpha v_N/2,&
 \epsilon_\gamma&=T_{\mathrm w}d_\alpha
                   (A_\alpha+d_\alpha/2)
                     +(5+T_{\mathrm w}A_\alpha)/X.
 \label{fnd:walldrift}
\end{align}
The full displacement of the ideal paired sum is bounded by
\begin{equation}
 E_{\mathrm{id}}=2\sqrt N
       \left[(e^{\epsilon_s}-1)
                   +(e^{\epsilon_s+\epsilon_\gamma}-1)\right]
 <\frac1{100}.
 \label{fnd:idealerror}
\end{equation}
The inequality uses the complete sum of both branches, not a termwise
comparison of selected low frequencies. The termwise heat/gamma error is
separately bounded by $E_{\mathrm{AB}}<10^{-6}$. It remains important to
use a remainder valid on the entire wall time interval. The corrected
broad-height transfer in the source supplement gives
\begin{equation}
 E_{\mathrm C}
 \leq4\exp\left(-\frac L4-\frac{tL^2}{16}+\frac1{10}\right)
 \leq4e^{1/10}\left(\frac{X}{4\pi}\right)^{-1/4}
 <\frac2{369}.
 \label{fnd:wallremainder}
\end{equation}
For reference, the last strict rational bounds are
\[
 \frac{X}{4\pi}>\frac{7X}{88}>820^4,
 \qquad e^{1/10}<\frac{10}{9}.
\]
Both complete Riemann--Siegel branches and the unrefined remainder are
included in this fee.

The wider-time normalisation in \eqref{fnd:wallremainder} is verified at
$x_*=X-4$ and $t=T_{\mathrm w}$ with the bound
\begin{multline}
 \frac{T_{\mathrm w}\pi^2}{32}
 +\frac{40}{x_*}
 +\frac{T_{\mathrm w}}4
       \left[\frac{20(26.90/2+1)}{x_*}+\frac{100}{x_*^2}\right]\\
 <\frac{4961}{98000}
 +\frac{40}{x_*}
 +\frac{T_{\mathrm w}}4
       \left[\frac{20(26.90/2+1)}{x_*}+\frac{100}{x_*^2}\right]
 <\frac1{10}.
 \label{fnd:wallnormalisation}
\end{multline}
Here $\pi<22/7$ gives
$T_{\mathrm w}(22/7)^2/32=4961/98000$ exactly.
The broad remainder is established on $\sigma\in[-1,2]$, $t\leq1/5$;
it therefore covers the closed interval in \eqref{fnd:wallparameters}.
The Gaussian in that derivation has constant $11/12$, as proved in the
source section. Specifically its full coefficient bracket is bounded by
\[
 \frac12(1+10^{-10})
       +\frac{(229/1000)110}{N-1}+10^{-100}<1.
\]
This displays the margin by which the corrected remainder, including
its far Gaussian tail, retains coefficient one. Neither an old narrower
remainder bound nor its old final wall-floor field is used.

Combining \eqref{fnd:idealerror}--\eqref{fnd:wallremainder} gives
\[
 E_{\mathrm{AB}}+E_{\mathrm C}
       <10^{-6}+\frac2{369}<\frac1{180}.
\]
Consequently Lemma~\ref{fnd:abel} implies the following explicit physical
wall:
\begin{proposition}[Full stationary wall]\label{fnd:wall}
For every $0\leq t\leq41/250$ and $0\leq y\leq1$,
\begin{equation}
 \left|\frac{H_t(X+iy)}{B_t(X+iy)}\right|
 >\frac{79}{100}-\frac1{100}-\frac1{180}
 =\frac{697}{900}>\frac34.
 \label{fnd:physicalwall}
\end{equation}
In particular the entire stated vertical wall is zero-free, including
all its time and height endpoints.
\end{proposition}

\begin{proposition}[Reality and simplicity in the finite head]
\label{fnd:finite-head}
Every zero of $H_t$ with $|\Ree z|\leq X$ is real for
$0\leq t\leq41/250$. Such zeros are simple for $0<t\leq41/250$.
\end{proposition}
\begin{proof}
Theorem~1 of \cite{PlattTrudgian} verifies the Riemann hypothesis through
zeta ordinate $3000175332800$. The conservative normalisation clearance
is
\begin{equation}
 3000175332800-\frac{X+1}{2}=175239886>0.
 \label{fnd:clearance}
\end{equation}
Equations \eqref{fnd:coordinates} and \eqref{fnd:clearance} therefore show
that all time-zero zeros with $|\Ree z|\leq X$ are real. This uses the
exact theorem endpoint, not only the rounded number $3\cdot10^{12}$ in
the paper's title.

Take the compact rectangle $|\Ree z|\leq X$, $|\Imm z|\leq2$.
Classical strip contraction keeps all zeros in $|\Imm z|\leq1$.
Thus the horizontal edges and the vertical pieces with $1<|\Imm z|\leq2$
are zero-free. Proposition~\ref{fnd:wall} and symmetry exclude zeros on
the remaining vertical pieces. The boundary is zero-free for the whole
closed time interval.

It remains to justify preservation of reality inside this rectangle.
At any fixed time there are finitely many zeros inside. Surround them by
disjoint small disks contained strictly inside the rectangle. Rouch\'e's
theorem preserves their total local multiplicities at nearby times,
and a compact zero-free complement excludes new zeros elsewhere. A simple
real zero persists as a real zero, by uniqueness and conjugation.
At a real zero of multiplicity $m\geq2$, Proposition~3.1(ii) of
\cite{Polymath} gives the forward expansions
\[
 z_k(t_0+\epsilon)
   =z_0+\sqrt{2\epsilon}\,\lambda_k+O(\epsilon),
 \qquad \lambda_1<\cdots<\lambda_m\in\mathbb R,
\]
where the $\lambda_k$ are the distinct roots of the probabilists'
Hermite polynomial. Each individual branch is localized in a disjoint
disk centered on the real axis. Conjugation and the one-zero count in
that disk force the zero to be real. Thus every all-real slice has an
all-real right neighborhood; the compact zero-free boundary and local
zero continuity also make the all-real property closed. Starting at
zero and taking the supremum of times through which reality holds proves
reality throughout the interval.

Finally, if a multiple real zero occurred at a positive time $t_0$,
the backward version of the same expansion would have a nonreal branch
$z_0+i\sqrt{2\epsilon}\,\lambda_k+O(\epsilon)$ at $t_0-\epsilon>0$.
That would contradict the reality just proved. No simplicity assumption
about the verified time-zero zeros was used.
\end{proof}

\subsection{Paired kernels and the complete signed probe functional}
\label{fnd:kernelinterface}

For $h>b\geq0$ and $d\in\mathbb R$, put
\begin{equation}
 K_h(d,b)=\frac{h-b}{d^2+(h-b)^2}
                 +\frac{h+b}{d^2+(h+b)^2}.
 \label{fnd:K}
\end{equation}
This is the contribution of the conjugate pair $a\pm ib$, with $d=x-a$,
to $-\Imm(H_t'/H_t)(x+ih)$. For a real zero the same formula has to be
given half weight: $K_h(d,0)/2=h/(d^2+h^2)$.

Fix one of the source rows. Its signed linear functional $\mathcal L$
is a finite real linear combination of symmetric translations in the
real variable and derivatives with respect to the probe height.
For example, a node $(h,a,w)$ contributes
\[
 \frac w2\left[S_t(x+a,h)+S_t(x-a,h)\right],\qquad
 S_t(x,h)=-\Imm\frac{H_t'}{H_t}(x+ih).
\]
Finite joint-jet combinations are included by applying the same linear
operator to the $h$ variable. The actual coefficients, factors of two,
heights and shifts are defined in Section~\ref{src:section}. Write
\begin{equation}
 J(t,x)=\mathcal L[S_t](x),\qquad
 R(d,b)=\mathcal L[K_{(\cdot)}(d+(\cdot),b)].
 \label{fnd:functional}
\end{equation}
The notation for $R$ means that exactly the same height, translation,
and height-derivative operator is applied to the pair kernel. The operator
is fixed within a source row; it is not differentiated along a moving
zero. Its full own-pair function at a contact square $q=y^2$ is
\begin{equation}
 O_{\mathrm{pair}}(q)=R(0,\sqrt q).
 \label{fnd:own}
\end{equation}
Here $O_{\mathrm{pair}}$ denotes the exact conjugate-pair contribution of
the fixed finite functional, not a Landau order term. The row-indexed
function is written $O_{\mathrm{pair},i}$ below.
For an ordinary symmetric node it is explicitly
\begin{equation}
 O_{\mathrm{pair};h,a,w}(q)
 =\frac{2wh(h^2+a^2-q)}{(h^2+a^2-q)^2+4a^2q}.
 \label{fnd:ownnode}
\end{equation}
This is the contribution of the \emph{entire} conjugate pair at $x\pm iy$.
Subtracting only the lower member from the probe would use incorrect
units. The lower member alone is removed separately in the zero-motion
identity below.

At a first highest contact with height $y$, all other zero heights have
modulus at most $y$. Every probe pole lies strictly outside this strip;
indeed the complete contact caps satisfy $h^2>q_{\max}$ for all original
heights and all joint-jet centers. Therefore all quantities in
\eqref{fnd:functional} are well-defined at such a contact.

\begin{lemma}[Convergence of paired imaginary logarithmic derivatives]
\label{fnd:convergence}
Suppose the zeros of the real entire order-one function $H_t$ have
imaginary parts bounded in modulus. At every point away from its zeros, the imaginary
part of its logarithmic derivative is the sum of the conjugately paired
imaginary contributions, with real zeros assigned half-pair weight.
The resulting sum is absolutely convergent. Finite signed combinations
and finite height derivatives at probe points outside the zero strip may
be interchanged with this sum.
\end{lemma}
\begin{proof}
The order-one growth of $H_t$ gives, by Jensen's formula,
$N_t(r)=O(r^{1+\eta})$ for every fixed $\eta>0$. In a strip of bounded
height, this implies convergence of the zero sum of $(1+|\Ree\rho|)^{-2}$,
by a dyadic decomposition with, say, $\eta=1/2$.
At a fixed point $z=x+ih$, a conjugate pair at $a\pm ib$ contributes
\[
 -\Imm\left(\frac1{z-a-ib}+\frac1{z-a+ib}\right)
       =K_h(x-a,b).
\]
For large $|a|$, with $b$ bounded, this is $O((1+|a|)^{-2})$.
The same estimate, with a changed constant, holds after any fixed finite
number of height derivatives on compact probe sets strictly outside the
zero strip. The genus-one canonical product has a linear exponential
factor whose logarithmic derivative is real: this follows from reality
on the real axis, or directly from the symmetric product representation.
The conjugate pairing also makes the imaginary parts of the canonical
$1/\rho$ correction terms cancel. Thus there is no additional imaginary
constant. Absolute convergence then gives the asserted exchanges.
\end{proof}

The boundary geometry can be extended to all other-zero heights by an
explicit positive Poisson argument. This is needed because other zeros
are not assumed real or to lie on the highest-contact boundary.

\begin{lemma}[Boundary kernel inequality implies a full strip inequality]
\label{fnd:poisson}
Fix $y>0$. Suppose $R(d,b)$ is the signed finite probe kernel just defined,
with every probe height strictly greater than $y$. If
\begin{equation}
 \frac12R(d,y)\leq\frac{2}{d^2+4y^2}
 \quad\hbox{for every }d\in\mathbb R,
 \label{fnd:boundarykernel}
\end{equation}
then
\begin{equation}
 K_y(d,b)\geq\frac y2R(d,b),\qquad 0\leq b\leq y,
 \label{fnd:fullkernel}
\end{equation}
apart from the distinguished singular location $d=0$, $b=y$.
\end{lemma}
\begin{proof}
As a function of $(d,b)$, each $K_h(d+a,b)$ and each of its height
derivatives is harmonic on $|b|<y$, is continuous and bounded on the
closed strip, is even in $b$, and tends to zero as $|d|\to\infty$.
So does $g=R/2$. It is the unique bounded harmonic extension of its equal
boundary values $g(d,y)=g(d,-y)$.

Let $\mathcal P_y$ denote the positive Dirichlet extension operator for
these equal boundary data. With Fourier convention
$\widehat f(\xi)=\int_{\mathbb R}e^{-id\xi}f(d)\,\dd d$, its multiplier
at height $b$ is
\[
 \frac{\cosh(b\xi)}{\cosh(y\xi)}.
\]
In the sense of boundary measures the positive harmonic function
$K_y(d,b)/y$ has, at \emph{each} edge of the strip, the boundary data
\begin{equation}
 \frac{2}{d^2+4y^2}+\frac\pi y\delta_0.
 \label{fnd:boundaryatoms}
\end{equation}
Indeed
$\epsilon/(d^2+\epsilon^2)\to\pi\delta_0$ as $\epsilon\downarrow0$.
Equivalently the Fourier transform of the density in
\eqref{fnd:boundaryatoms} is
$(\pi/y)(e^{-2y|\xi|}+1)$, whose Poisson extension is
\[
 \frac\pi y(1+e^{-2y|\xi|})
       \frac{\cosh(b\xi)}{\cosh(y\xi)}
 =\frac{2\pi}{y}e^{-y|\xi|}\cosh(b\xi),
\]
which is exactly the Fourier transform of $K_y(d,b)/y$.

By \eqref{fnd:boundarykernel} and positivity of the Poisson operator,
\[
 g=\mathcal P_y[g(\cdot,y)]
 \leq\mathcal P_y\left[\frac{2}{d^2+4y^2}\right]
 \leq\mathcal P_y\left[\frac{2}{d^2+4y^2}+\frac\pi y\delta_0\right]
 =\frac{K_y(d,b)}{y}
\]
in the interior. Continuity gives the nonsingular boundary points.
The atoms may be discarded for the inequality but must be retained in
any asserted identity for the target Poisson extension.
\end{proof}

The polynomial certificates in Theorem~\ref{src:kernel} establish
\eqref{fnd:boundarykernel} for every real spacing and for the entire closed
$q$ interval of each row, with the same signed coefficients as the
source. Their rational compactification and Bernstein subdivisions are
continuum sign proofs. Lemma~\ref{fnd:poisson} explains the additional
analytic step from these two-variable boundary certificates to every
possible other-zero height. Negative individual node coefficients cause
no problem, because the boundary inequality is proved for their full
signed sum before the positive extension is taken.

\begin{proposition}[Current inequality at a simple highest zero]
\label{fnd:current}
Let $z(t)=x(t)+iy(t)$ be a simple zero of $H_t$, with $y(t)>0$, at a time
when no zero has imaginary part exceeding $y(t)$. Let $q=y^2$. If a row
functional satisfies \eqref{fnd:fullkernel}, then
\begin{equation}
 q'\leq-2-2q\bigl[J(t,x)-O_{\mathrm{pair}}(q)\bigr].
 \label{fnd:currentbound}
\end{equation}
\end{proposition}
\begin{proof}
Proposition~3.1(i) of \cite{Polymath}, valid when the other zeros may be
complex or repeated, gives the local branch equation
\[
 z'(t)=2\sum_{\rho\ne z(t)}\frac1{z(t)-\rho},
\]
with the symmetric summation prescribed there. Taking imaginary parts
and pairing the other zeros with their conjugates is legitimate by
Lemma~\ref{fnd:convergence}. The own conjugate $x-iy$ contributes
$2/(2iy)=-i/y$. All other conjugate pairs contribute $-2K_y(d,b)$,
with real zeros assigned half weight. Hence
\begin{equation}
 y'=-\frac1y-2\sum_{\mathrm{other\ pairs}}K_y(d,b),\qquad
 q'=-2-4y\sum_{\mathrm{other\ pairs}}K_y(d,b).
 \label{fnd:motion}
\end{equation}
By Lemma~\ref{fnd:convergence}, applying the signed probe functional gives
\[
 J(t,x)=O_{\mathrm{pair}}(q)+\sum_{\mathrm{other\ pairs}}R(d,b).
\]
The same other pairs, including their multiplicities and the half weight
for real zeros, appear in both expressions. Apply \eqref{fnd:fullkernel}
to each pair and sum absolutely to obtain
\[
 q'\leq-2-2q\sum_{\mathrm{other\ pairs}}R(d,b),
\]
which is \eqref{fnd:currentbound}. The $-2$ is the own conjugate's exact
motion contribution and is not counted again in the sum.
\end{proof}

\subsection{Compact attainment and repeated nonreal contacts}
\label{fnd:globalcontact}

A derivative formula for a single zero does not by itself control the
supremum over an unbounded zero set. The compact-attainment and
multiple-contact arguments are therefore supplied separately.

\begin{lemma}[Uniform compactness after positive startup]
\label{fnd:compact}
For $0<\tau<T\leq1/2$, all nonreal zeros of all $H_t$ with
$t\in[\tau,T]$ lie in a common compact rectangle.
\end{lemma}
\begin{proof}
Theorem~1.5(i) of \cite{Polymath} has one sufficiently large absolute
constant $C$, independent of $0<t\leq1/2$, such that all zeros with
$\Ree z\geq\exp(C/t)$ are real. Symmetry handles the negative real tail.
Thus possible nonreal zeros for $t\in[\tau,T]$ have
$|\Ree z|\leq\exp(C/\tau)$. Classical strip contraction puts them in
$|\Imm z|\leq1$. The value of $C$ need not be effective at the practical
wall size: only existence of a common compact rectangle is used here.
The theorem is not used to replace the exterior source proof by a
numerical tail cutoff at $X$.
\end{proof}

The repeated-zero expansion is particularly transparent directly from
the heat equation. If $H_{t_0}(z_0+w)=a w^m+O(w^{m+1})$, $a\ne0$,
then at the backward time $t_0-\epsilon$ the heat scaling yields,
uniformly for $w$ in compact sets,
\begin{equation}
 \frac{H_{t_0-\epsilon}(z_0+\sqrt\epsilon\,w)}{a\epsilon^{m/2}}
 \longrightarrow
 p_m(w)=m!\sum_{k=0}^{\lfloor m/2\rfloor}
                   \frac{w^{m-2k}}{k!(m-2k)!}.
 \label{fnd:backwardpolynomial}
\end{equation}
This is also a special case of \cite[Proposition~3.1(ii)]{Polymath}.
Writing $\operatorname{He}_m$ for the probabilists' Hermite polynomial,
\begin{equation}
 p_m(w)=(i\sqrt2)^m\operatorname{He}_m\left(\frac{w}{i\sqrt2}\right).
 \label{fnd:hermite}
\end{equation}
For $m\geq2$, it has a simple root $i\sqrt2\,\lambda$ with
$\lambda>0$.

\begin{lemma}[No repeated nonreal first contact]\label{fnd:multiple}
Let $Y(t)>0$ be locally Lipschitz from the left at $t_0>0$. Suppose
all zeros at earlier nearby times have imaginary part less than $Y(t)$,
and $H_{t_0}(x+iY(t_0))=0$. Then this zero is simple.
\end{lemma}
\begin{proof}
If its multiplicity were $m\geq2$, Rouch\'e's theorem applied to
\eqref{fnd:backwardpolynomial} near the positive imaginary root in
\eqref{fnd:hermite} would give a zero of $H_{t_0-\epsilon}$ of the form
\[
 x+iY(t_0)+i\sqrt{2\epsilon}\,\lambda+O(\epsilon).
\]
Its imaginary part is greater than $Y(t_0)+c\sqrt\epsilon$ for some
$c>0$ and all sufficiently small $\epsilon>0$. But
$Y(t_0-\epsilon)=Y(t_0)+O(\epsilon)$. This contradicts the assumed
earlier strict upper barrier. The argument is local at a complex zero;
it does not require that this repeated zero be real.
\end{proof}

\begin{theorem}[Global comparison from the row certificates]
\label{fnd:comparison}
Let $0<\tau<T<\min\{41/250,1/2\}$. Suppose a positive upper barrier
$Q$ on $[\tau,T]$ consists of finitely many smooth closed incoming arcs,
with only upward resets between successive arcs. On an arc beginning
at $a$ with initial ceiling $q_a$, suppose
\begin{equation}
 Q'=-2-rQ,\qquad r>0,
 \qquad 0<Q\leq q_a\leq q_{\max}.
 \label{fnd:comparisonode}
\end{equation}
Suppose every actual highest contact in that arc has $|x|>X$, that the
row source and kernel are valid on its complete closed time and
$q$ domains, and that
\begin{align}
 J(t,x)&\geq B_a,\label{fnd:comparison-source}\\
 O_{\mathrm{pair}}(q)&\leq O_{\mathrm{pair},a}\qquad(0\leq q\leq q_a),\label{fnd:comparison-own}\\
 r&<2(B_a-O_{\mathrm{pair},a}).\label{fnd:comparison-rate}
\end{align}
Assume also a strict initial separation at $\tau$, and exclusion of a
violation during the preceding startup interval. Then every zero of
$H_t$ satisfies $(\Imm z)^2<Q(t)$ on the incoming arcs for
$\tau\leq t\leq T$, and satisfies the corresponding outward ceiling at
every reset.
\end{theorem}
\begin{proof}
Upward resets preserve the upper-barrier property, so it suffices to
exclude a first contact on a closed incoming arc. The barrier is positive
on every such arc. Hence its square root is Lipschitz there.
By Lemma~\ref{fnd:compact}, all zeros capable of reaching this positive
height lie in a common compact region. The zero sets vary locally
continuously with time by Rouch\'e's theorem. More explicitly, for any
sequence of violation times decreasing to an infimum, their zeros have
a convergent subsequence; continuity supplies a zero at the limiting
time, and local persistence rules out a strict overshoot there. Strict
startup separation and upward resets therefore force an attained first
contact on the upper boundary. By symmetry it may be taken at positive
real part. Proposition~\ref{fnd:finite-head} puts it at $x>X$.
At this first contact every other zero is below the same upper height.

Lemma~\ref{fnd:multiple} excludes a repeated contact. At a simple contact
Proposition~\ref{fnd:current} and
\eqref{fnd:comparison-source}--\eqref{fnd:comparison-rate} give
\[
 q'\leq-2-2q(B_a-O_{\mathrm{pair},a})<-2-rq=Q'.
\]
For a differentiable zero branch reaching $q-Q=0$ for the first time
from below, the incoming derivative of $q-Q$ is nonnegative. This
contradicts the strict inequality. If the contact is at an arc endpoint
or a row knot, its closed incoming arc supplies the same contradiction
from the left. No differentiation of an unattained global supremum and
no global simplicity assumption has been made.
\end{proof}

The exact numerical construction checks the hypotheses
of Theorem~\ref{fnd:comparison} at every one of its $221010$ arcs.
The actual common-source theorem holds at every center $x\geq X$; its
admitted translations obey $|a|\leq39/20<2$. Thus all translated probe
coordinates at a contact $x>X$ are greater than $X-2$ and stay within
the theorem's physical domain. All probe circles
and their complex padding are separately covered by the source bounds;
compactness does not authorize any enlargement of them. On ordinary
nodes and the three conic rows, the full own-pair rational function is
proved nondecreasing on its whole contact interval, so the upper bound
$O_{\mathrm{pair},a}$ may be evaluated at the outward $q_a$. Source lower chords are
nondecreasing in time, so their value at the beginning of each small
arc is a valid lower bound throughout that arc. These facts are used
with the actual signed full functional, never with a sum of independently
optimized flow rates.

\subsection{The terminal global strip and the complete cleanup}
\label{fnd:cleanupsection}

\begin{lemma}[Complete terminal cleanup]\label{fnd:cleanup}
Suppose at a time $T$ all zeros of $H_T$ have imaginary part of modulus
at most $\sqrt{q_e}$, where $q_e>0$. Then
\begin{equation}
 \Lambda\leq T+\frac{q_e}{2}.
 \label{fnd:cleanupbound}
\end{equation}
No finite-head wall or exterior source estimate beyond $T$ is required.
\end{lemma}
\begin{proof}
Apply Theorem~\ref{fnd:debruijn} to the global strip at time $T$.
At $T+q_e/2$ every zero is real. The defining characterization of the
de Bruijn--Newman constant \cite{deBruijn,Newman} gives
\eqref{fnd:cleanupbound}. The conclusion uses forward strip contraction
only and makes no backward inference from a later all-real slice.
\end{proof}

For the certificate consumed here the endpoint data are the exact rationals
\begin{equation}
 T=\frac{22101}{204800},\qquad
 q_e=\frac{7902377963830536747158643}
           {79228162514264337593543950336}.
 \label{fnd:terminaldata}
\end{equation}
Both the source horizon $19/125$ and the wall horizon $41/250$ are greater
than $T$. The complete cost, including the positive terminal strip, is
\begin{align}
 T+\frac{q_e}{2}
 &=\frac{427693072077942473661578098491}
          {3961408125713216879677197516800}\label{fnd:exactendpoint}\\
 &<\frac{107964910079}{10^{12}}=0.107964910079.
 \label{fnd:decimalendpoint}
\end{align}
The exact-flow section proves the finite arithmetic leading to
\eqref{fnd:terminaldata}; this section establishes why those data control
all zeros. Together with $\Lambda\geq0$ from \cite{RodgersTao}, the resulting
interval is $0\leq\Lambda\leq0.107964910079$.

The analytic arguments above are ordinary mathematical proofs assembled
with explicit interval and exact certificates. They are not claimed to
be an end-to-end Lean formalisation of the Lambda bound. In particular,
using a pinned external formal theorem for \eqref{fnd:seveneighths}
does not turn the custom wall, source estimates, or contact argument into
machine-checked Lean declarations. Their proofs and their finite
certificate consumers are supplied separately for scrutiny.

% END INCLUDED FILE


% BEGIN INCLUDED FILE: sections/source.tex
\section{The physical source and its exact finite certificates}
\label{src:section}

The comparison uses a finite signed linear functional of logarithmic
derivatives.  Negative coefficients in that functional are essential: they
must be combined before estimating the arithmetic head, whereas errors
without a proved common representation must be charged with absolute
coefficients.  This section gives the analytic reductions behind the source
file and the continuous kernel certificates.  The finite rational data in
the accompanying supplement are part of the proof.
Theorem~\ref{thm:direct67-reconstruction} directly reconstructs the operational
source for every row. Appendices~\ref{app:physical-debit-ledger}
and~\ref{sec:arithmetic-credit-audit} give the debit and credit
derivations used in that reconstruction.

\subsection{The functional, the certificate interface, and the data}
\label{src:functional}

Write
\[
 S_t(x,h)=-\operatorname{Im}\frac{H_t'}{H_t}(x+ih)
       =\partial_h\log|H_t(x+ih)|.
\]
Here the prime on $H_t$ denotes differentiation in its complex argument.
For a node $(h,a,w)$ with $a\geq0$, the associated functional is
\[
 w\mathcal A_{h,a}S_t(x)
   :=\frac w2\{S_t(x+a,h)+S_t(x-a,h)\}.
\]
Thus a zero shift is one evaluation, not two evaluations.  An ordinary row
has
\begin{equation}
 J_i(t,x)=\mathcal L_i S_t(x)
       :=\sum_j w_{ij}\mathcal A_{h_{ij},a_{ij}}S_t(x).
 \label{src:functional-def}
\end{equation}
The exact rational node lists, their closed time intervals, and their
permitted squared contact-height intervals are the $67$ records of
\path{M8_LATE_PLAN.json}.  These lists are not floating point proposals.
The plan has SHA-256
\begin{center}\footnotesize\ttfamily
 ec346b7c1136eda38021eb6323033b3b48ed21e6f7458eeeac7cc9579db2497a.
\end{center}
Rows $0,16,44$ additionally contain a joint derivative functional, defined
in Section~\ref{src:jets}; every other row is of the form
\eqref{src:functional-def}.

The source file \path{MAIN_CLOCK_SOURCE.json} provides four operational
rational numbers for each row:
\[
 M_i=\texttt{main\_lo},\quad N_i=\texttt{native\_debit},\quad
 A_i=\texttt{first\_hi},\quad B_i=\texttt{last\_hi}.
\]
If its time interval is $[t_i,t_{i+1}]$, put
\begin{equation}
 \underline J_i(t)=M_i-N_i-(1-\theta)A_i-\theta B_i,
 \qquad \theta=\frac{t-t_i}{t_{i+1}-t_i}.
 \label{src:operational}
\end{equation}
Historical fields in the source file are not alternative decompositions
of this expression.  In particular an earlier lower endpoint for an
arithmetic head is not paired with a subsequently reduced upper debit.

\begin{theorem}[Uniform physical source]\label{src:source}
For each declared row, throughout its entire closed time interval and
for every actual central coordinate
\[
 x\geq X:=6000000185827,
\]
the physical functional specified by that row obeys
\[
 J_i(t,x)\geq\underline J_i(t).
\]
All values and derivative evaluations in this assertion are of the same
function $H_t$, with its actual physical phase, native cutoff, and
normaliser.  The assertion uses no hypothesis about an unbounded set of
zeta zeros and no additional initial-strip assumption.
\end{theorem}

The proof is distributed below into the physical approximation,
coefficient identities, complete analytic and physical charges, and
finite phase/resource certificates.  The final source file is bound by
SHA-256
\begin{center}\footnotesize\ttfamily
 b59905d7b994d3424ba14f9347be75a2429ae4d4b0de2c2404971ffda1d4d536.
\end{center}
Its relation to the pinned plan is checked as an identity of the entire
JSON objects, including each unchanged row, rather than inferred from
matching displayed decimal values.

\subsection{One finite analytic model at all the translations}
\label{src:common}

Fix $x\geq X$ and $t\geq0$.  The notation in this subsection is local to
this fixed physical center.  Put
\[
 L=\log\frac{x}{4\pi},\qquad
 \tau=\frac x2+\frac{\pi t}{8},\qquad
 N=\left\lfloor\sqrt{e^L+t/16}\right\rfloor,\qquad
 b_t(n)=\exp\left(\frac t4\log^2n\right).
\]
For complex $z$ near $x+i[0,6]$, set
\[
 a(z)=\frac{1-iz}{2},\qquad b(z)=1-a(z),
\]
and, on the appropriate half-plane logarithm branches, define
\begin{align*}
 M_0(s)&=\frac{s(s-1)}{16}\pi^{-s/2}\sqrt{2\pi}
       (s/2)^{s/2-1/2}e^{-s/2},\\
 \alpha(s)&=\frac1{2s}+\frac1{s-1}+\frac12\Log\frac{s}{2\pi},\\
 M_t(s)&=M_0(s)\exp(t\alpha(s)^2/4),
       &B_t(z)&=M_t(a(z)).
\end{align*}
These functions are analytic and nonzero on every source disk used here.
The logarithmic derivative of $M_t$ is
\[
 \beta_t(s)=\alpha(s)+\frac t2\alpha(s)\alpha'(s).
\]
Let
\[
 \Gamma=\frac{M_t((1+ix)/2)}{M_t((1-ix)/2)}.
\]
In particular $|\Gamma|=1$; it is determined by $x,t$ and is not an
independent phase.  For $\zeta=h-iv$ the common finite model is
\begin{align}
 f_x(\zeta)
 &=\sum_{n\leq N} b_t(n)n^{-1/2-tL/4+i\tau}
                   e^{-\zeta\log n/2}\notag\\
 &\quad+\Gamma\sum_{n\leq N} b_t(n)n^{-1/2-tL/4-i\tau}
                   e^{-\zeta(L-\log n)/2}.
 \label{src:common-model}
\end{align}
All of $N,L,\tau,\Gamma$ in this formula are frozen at the central
coordinate.  The physical normaliser is still evaluated at the actual
point $x+i\zeta=x+v+ih$.  This distinction makes it possible to apply one
holomorphic error bound or one spectral measure to multiple translations.

Write
\[
 \ell_*=\frac{2689}{100},\qquad c_*=\frac{2689}{400},\qquad
 X_*=X-4.
\]
The elementary anchor inequalities are
\[
 26.89<\log\frac{X_*}{4\pi}<26.90,
 \qquad \log\frac{X_*}{4\pi}>2\log 690988.
\]
The latter inequality permits all subsequent comparisons with
$L_0=2\log N_0$, $N_0=690988$.

\begin{proposition}[Complete common-model physical errors]
\label{src:physical}
Suppose a closed thermal row $[t_-,t_+]\subset[0,19/125]$ and a native
height $h$ satisfy
\begin{equation}
 |v|\leq\frac{39}{20},\quad h-\frac1{100}\geq0,\quad
 h+\frac1{100}\leq\frac{601}{100},\quad
 \frac{1+h-1/100}{2}+\frac{t_-\ell_*}{8}>\frac54.
 \label{src:physical-guard}
\end{equation}
For every time in the row and all physical centers $x\geq X$,
\begin{align}
 \left|\frac{H_t(x+v+ih)}{B_t(x+v+ih)}-f_x(h-iv)\right|
       &\leq E_0(h,t_-),\label{src:E0}\\
 \left|\partial_h\left(\frac{H_t}{B_t}\right)(x+v+ih)
                      -f_x'(h-iv)\right|
       &\leq E_1(h,t_-),\label{src:E1}
\end{align}
where
\begin{equation}
 \begin{split}
 \mathcal S&=(1+h-1/100)c_*+t_-c_*^2,\\
 E_0(h,t_-)&=8e^{-\mathcal S}+5\cdot10^{-10},\\
 E_1(h,t_-)&=800e^{-\mathcal S}+7\cdot10^{-9}.
 \end{split}\label{src:matched-errors}
\end{equation}
The baseline bounds $e_0=5\cdot10^{-7}$ and $e_1=5\cdot10^{-5}$ also
hold on the same row domains.  Every moving-cutoff term and both complete
Riemann--Siegel remainders are included.
\end{proposition}

\begin{proof}
The following estimates are especially material to the larger-height
extension; the complete finite scalar inequalities and their exact
endpoints are in the physical-source supplement.  The individual heated
term estimate is \cite[Proposition 6.1]{Polymath}.  Its use is valid for
all real parts, before the narrower-height simplification in the printed
numerical theorem.  The two real parts needed here lie in $[-3,4]$.

The full Gaussian calculation must use the constant $11/12$.  Indeed,
with $D=T-3$ and real $u$,
\[
 \frac{u^2+6|u|+2/3}{4D}\leq\frac{u^2+11/12}{D},
\]
because, after multiplication by $4D$, the difference is
$3(|u|-1)^2$.  Completing the square therefore gives the exact majorant
\begin{equation}
 G_\sigma(c)=
 \frac{\exp\left\{\frac{11}{12D}
    +\frac{\sigma^2/D+\sigma c+tc^2/4}{1-t/D}\right\}}
      {\sqrt{1-t/D}}.
 \label{src:gaussian}
\end{equation}
Convexity in $\sigma$ and endpoint comparison on $[-3,4]$ give
\[
 G_4(\log9)<8104,\qquad G_{-3}(-\log2)<8.5
\]
throughout the stated time and coordinate domain.  In particular, with
the deliberately smaller native lower bound $690000$, the complete
remainder bracket is bounded by
\[
 \frac12G_4(0)+\frac{0.2}{690000-.865}G_4(\log9)
 +\frac{.029}{690000-.353}G_{-3}(-\log2)+10^{-100}<1.
\]
The omitted negative Gaussian range in the intermediate expansion is
not discarded: for its variable $r\geq100$ its complete integrand is
bounded by
\[
 e^{-4r^2+40r+20+(r^2+1)/D}\leq e^{-3r^2},
 \qquad
 \int_{100}^{\infty}e^{-3r^2}\,\dd r
           \leq e^{-30000}/600.
\]
The extra $1/12$ relative to the erroneous $5/6$ value in an earlier
source note is paid in \eqref{src:gaussian}; that earlier value is not
used.  Consequently each complete remainder is bounded by
$e^{t\pi^2/64}|M_0(iT')|$ at the corresponding shifted ordinate.

For $|\operatorname{Re}s|\leq4$ and $|\operatorname{Im}s|=x'/2$,
direct estimates in the displayed formula for $\alpha$ give
\[
 |\alpha(s)-(L'/2\pm i\pi/4)|<10/x',\quad
 |\alpha'(s)|<2/x',\quad |\alpha(s)|<L'/2+1.
\]
After division by the actual $B_t$, both remainders together are thus
bounded by
\begin{equation}
 e_C<4\exp\left(-\frac{(1+y)L'}4-\frac{t(L')^2}{16}
                            +\frac1{10}\right).
 \label{src:heat-physical}
\end{equation}
The $1/10$ includes the positive correction $t\pi^2/32$ and all
coordinate and $\alpha$ corrections on this source domain.

The full first and reflected native sums have, per branch, zeroth,
first-logarithmic, and second-logarithmic moments at most $6,22,180$;
the deliberately enlarged total value mass bound is $18$.
For example the uninflated majorants are supplied by
$\sum n^{-5/4}$ and its first two logarithmic moments, followed by the
row guard \eqref{src:physical-guard}.  The complete individual term
errors from Proposition 6.1 of \cite{Polymath} total less than
$e_{AB}=5.5\cdot10^{-12}$.

Across a translated disk of radius $1/100$, every real displacement
from the central coordinate is less than $2$.  Consecutive activation
coordinates are $4\pi(m^2-t/16)$ and have spacing $4\pi(2m+1)>4$.
At most one native pair can change.  For a changed index the exact
inequality
\[
 \left|\log m-\frac{L'}2\right|<10^{-12}
\]
holds on the whole disk.  Expanding the terminal thermal exponent
before bounding its linear term gives the complete crossing-pair
cost $3e^{10^{-5}}e^{-\mathcal S}$.  This covers both directions of
crossing and both native branches.  The normaliser and fixed-index
comparison are holomorphic on the disk, so Cauchy's formula gives a
factor $100$ for its height derivative.  In view of
\[
 4e^{.1}+3e^{10^{-5}}<8,
\]
this accounts for the first terms in \eqref{src:matched-errors}.

For completeness, the passage from exact local $\alpha$ exponents to
the local leading model has value error at most
\[
 \frac{440t+12yC_\gamma}{x'},\qquad
 C_\gamma=10+t(L'/2+1),
\]
and derivative error at most
\[
 \frac{yC_\gamma(6L'+22)+5t(22L'+360)+12C_\gamma+44t}{x'}.
\]
These are respectively less than $2\cdot10^{-10}$ and
$2.5\cdot10^{-9}$ at the worst anchor, and decrease at larger $x'$.
Finally freezing the common gamma, clock, and frequency gives, with
\[
 B_{\rm fr}=\frac{30+2t(L+3)}{x-2},
\]
the separate charges
\[
 36B_{\rm fr}<2.35\cdot10^{-10},\qquad
 18B_{\rm fr}(L+1)+\frac{36}{x-2}<3.3\cdot10^{-9}.
\]
They include the change of reflected frequency as well as the common
phase.  Adding these costs and $e_{AB}$ proves
\eqref{src:E0}--\eqref{src:matched-errors}.  Applying the same complete
calculation without the height/time refinement gives value error less
than $4.167459\cdot10^{-7}$ and derivative error less than
$4.637616\cdot10^{-5}$, so the stated baseline allowances also hold.
All endpoint inequalities used here are directly outward-verified by
the supplied corrected physical-source checkers.
\end{proof}

\subsection{Finite thermal coefficients and exact matching}
\label{src:thermal}

Let $D$ denote the formal Dirichlet derivation $(Da)(n)=(\log n)a(n)$.
Formal multiplication, logarithms, and inverses are coefficientwise:
the coefficient at any fixed integer uses finitely many divisors.  Define
\[
 d_t=D\log_*b_t,\qquad
 E_t(s)=\prod_{p\in\{2,3,5\}}(1-b_t(p)p^{-s}),
 \qquad c_t=m_t*b_t,
\]
where $m_t$ is the eight-term coefficient sequence of $E_t$, and put
$\widetilde d_t=D\log_*c_t$.  No positive-time infinite numerical
Dirichlet series with coefficients $b_t(n)$ is used.

\begin{lemma}[Coefficient positivity and recurrences]\label{src:coeff}
For $n\geq2$ and $t\geq0$, the coefficients
$d_t(n),\widetilde d_t(n),c_t(n)$ are nonnegative and nondecreasing in
$t$.  Their exact triangular recurrences are
\begin{align}
 d_t(n)&=(\log n)b_t(n)
       -\sum_{\substack{r\mid n\\2\leq r<n}}d_t(r)b_t(n/r),
       \label{src:triangular}\\
 (\log n)c_t(n)&=\sum_{\substack{r\mid n\\r\geq2}}
                         \widetilde d_t(r)c_t(n/r),\label{src:coeff-match}\\
 \widetilde d_t(n)&=d_t(n)
 -\sum_{\substack{p\in\{2,3,5\},\ j\geq1\\n=p^j}}
           (\log p)e^{jt\log^2p/4}.
 \label{src:deflation}
\end{align}
\end{lemma}
\begin{proof}
Formal heat differentiation and the derivation rule give
\[
 \partial_t d_t(n)=\frac{\log^2n}{4}d_t(n)
 +\frac{\log n}{4}\sum_{ab=n,\ a,b\geq2}d_t(a)d_t(b).
\]
At time zero these are von Mangoldt coefficients.  Induction on $n$
and the integrating-factor formula prove nonnegativity and
monotonicity.  In \eqref{src:deflation}, at $n=p^j$ the subtracted term
has derivative $j(\log^2p)/4$ times itself, whereas the first linear
coefficient in the equation for $d_t(p^j)$ is
$j^2(\log^2p)/4$.  The extra forcing is therefore nonnegative,
proportional to $j(j-1)$.  The same induction applies to
$\widetilde d_t$, whose initial coefficients are nonnegative.
Equation \eqref{src:coeff-match} then proves positivity and monotonicity
of $c_t$.  The triangular identities follow from
$Da=(D\log_*a)*a$.
\end{proof}

At the fixed physical center introduce the finite Euler multiplier
\[
 E_x(\zeta)=\prod_{p\in\{2,3,5\}}
 \left(1-b_t(p)p^{-1/2-tL/4+i\tau}e^{-\zeta\log p/2}\right),
 \qquad Q=E_xf_x=1+U,\quad Q'=-V/2.
\]
The first branch is collected only through $N_0$.  Every product label
above $N_0$, reflected label, and newly activated index is retained
separately, including its sign.  Their absolute amplitudes and their
frequency-weighted amplitudes form finite positive majorants.
All these frequencies are nonnegative.  Reflected labels have frequency
$L-\log n+\log d$ for $d\mid30$.

The matched logarithmic-derivative polynomial is
\[
 T_K(\zeta)=\sum_{2\leq k\leq K}\widetilde d_t(k)
 k^{-1/2-tL/4+i\tau}e^{-\zeta\log k/2},\qquad K=4096.
\]
For the collected first head, the coefficient remaining after matching
$QT_K$ against $V$ is exactly
\begin{equation}
 (\log n)c_t(n)
 -\sum_{\substack{k\mid n\\2\leq k\leq K}}
           \widetilde d_t(k)c_t(n/k)
 =\sum_{\substack{k\mid n\\k>K}}
           \widetilde d_t(k)c_t(n/k)\geq0.
 \label{src:positive-residual}
\end{equation}
The physical integer phases multiply correctly in this identity; no
independently selectable phase is assigned to a divisor product.

At a native height, let $A,B$ denote the complete nonconstant amplitude
and frequency moment of the fixed head and its retained exterior labels.
Let $D_K$ be the positive norm of $T_K$ and let $Z_K$ be its overlap
with the nonconstant collected head:
\[
 Z_K=\sum_{2\leq k\leq K}\widetilde d_t(k)k^{-\sigma}
       \sum_{2\leq m\leq\lfloor N_0/k\rfloor}c_t(m)m^{-\sigma},
 \qquad \sigma=\frac{1+h}{2}+\frac{tL}{4}.
\]
The full positive residual decomposition gives the bound
\begin{equation}
 |V-QT_K|\leq R_*=B-(1-A)D_K-2Z_K
 \label{src:residual}
\end{equation}
before activation is added.  This difference is nonnegative because it
equals a sum of the residual coefficients
\eqref{src:positive-residual}, overflow products, and absolute exterior
and reflected labels.  It is not the difference of unrelated monotone
estimates.  The complete prefixes through $\lfloor N_0/k\rfloor$ are
used in $Z_K$; they are not truncated to $K$.

For a block, coefficient monotonicity permits the artificial common
comparison point with thermal parameter $t_+$ and abscissa
$(1+h)/2+t_-L_0/4$.  All matching masses in
\eqref{src:residual} are evaluated at that same point.  Upper bounds for
$A,B$ and lower bounds for $D_K,Z_K$ then give an upper bound for
$R_*$, provided $A<1$.  Beyond $L_0$, the actual positive residual
decreases term by term.  For example a reflected weighted label obeys
\[
 \frac{\dd}{\dd L}(\lambda w)
 =w\left[1-\lambda\left(\frac h2+\frac t4\log(dn)\right)\right]
 \leq0,
\]
since $\lambda\geq\log N_0$ and $h\log N_0/2>1$ on every row.
This proves frequency uniformity after exact matching, without applying
an anchor lower mass at an arbitrary larger frequency.

\subsection{All future cutoffs and the retained main slope}
\label{src:activation}

Put $v_0=\log N_0$ and write
$L=2(v_0+s)$, $\log n=v_0+w$ for a newly active index $N_0<n\leq N$.
Then
\[
 0<w\leq s+\delta,\qquad\delta=10^{-12}.
\]
The small extension is required by $N=\lfloor\sqrt{e^L+t/16}\rfloor$.
After inserting counting measure, the two positive densities are bounded
by
\[
 C_0e^{-a_0s}e^{-dw},\qquad
 C_0e^{-(a_0+h)s}e^{pw},
\]
where
\[
 a_0=t_-v_0/2,\quad d=(h-1)/2>0,\quad p=(h+1)/2,
 \quad C_0=e^{-t_-v_0^2/4-dv_0}.
\]
The dropped thermal factor is at most $e^{t\delta^2/4}$.
The full adjacent-integer and reflected-growth inflation is less than
$10^{-4}$ in logarithm; the saved prefactors pay it before any integral
comparison.  In particular the frequency moment includes the derivative
of the Euler factor, not just the derivative of a native term.

On a finite near interval it is convenient to use positive rational
endpoints $v_-<v_0<v_+$ and an upper constant
$C\geq e^{-t_-v_-^2/4-dv_-+1/5000}$.  For $u=s+\delta$, integration
over $0\leq w\leq u$ gives the simple complete bounds
\begin{equation}
 A_{\rm act}(s)\leq2Cu,\qquad
 D_{\rm act}(s)\leq C(v_++2u)u.
 \label{src:activation-near}
\end{equation}
If $P,P_1$ are the positive Euler norm and log moment, set
\[
 a_a(s)=2PC(s+\delta),\qquad
 b_a(s)=2PC(v_++2(s+\delta))(s+\delta)+2P_1C(s+\delta).
\]
These are polynomials with nonnegative coefficients.  The full analytic
and physical charges at that height are bounded by
\begin{align}
 E_{\rm an}(s)&=\frac12
 \left(\frac{R_*+b_a(s)+a_a(s)D^+}{g(s)}+T\right),
 \label{src:analytic-cap}\\
 E_{\rm ph}(s)&=\frac{r_1}{g(s)-r_0}
 +\frac{(B+b_a(s))r_0}{2g(s)(g(s)-r_0)}+2\cdot10^{-12},
 \label{src:physical-cap}\\
 g(s)&=1-A-a_a(s),\qquad r_0=e_0P,
 \quad r_1=e_1P+(e_0/2)P_1.\notag
\end{align}
Here $T$ is the complete omitted geometric Euler-power tail.  The strict
condition $g>r_0$ precedes division.  Formula
\eqref{src:physical-cap} is precisely the difference-of-quotients
bound, with $|Q|\geq g$ and $|Q'|\leq(B+b_a)/2$.

Both charges are increasing and convex as long as that gap remains
positive: expand their reciprocal denominators in geometric power
series in $s$, whose coefficients are nonnegative.  For a row choose
one fixed analytic spectral recombination, add absolute-weighted physical
charges and the fixed coordinate allowance, and call the result $N(s)$.
The same fixed choice is used throughout the near interval.  It is not
asserted that a pointwise minimum of convex caps is convex.

For the complementary half-line one needs global activation constants
that remain valid when the near denominator is no longer positive.
Appendix~\ref{verification:sharp-activation} proves the constants actually
consumed by the source certificates. With $d=(h-1)/2$, $p=(h+1)/2$,
$a=t_-V_-/2$, $k=a+d$, $c=a+h$, $U=1/d-1/p$, and the complete
integer/end-point factor $C$ defined there, the combined branch integrals are
\begin{align}
 A(s)&=\frac{e^{-as}}d-Ue^{-ks}-\frac{e^{-cs}}p,
     \label{src:activation-global-value}\\
 D(s)&=A_a e^{-as}-(A_k+B_ks)e^{-ks}-(A_c+B_cs)e^{-cs}.
     \label{src:activation-global-derivative}
\end{align}
The rational coefficients are specified in the appendix. Factored
derivative monotonicity proves a unique maximum on $s\geq0$ when
$a>0$ and gives the exact limiting supremum when $a=0$.
Outward evaluation on rational critical-point enclosures proves, at every
actual row and height,
\[
 C\sup_{s\geq0}A(s)\leq\overline A_{\rm act},\qquad
 C\sup_{s\geq0}D(s)\leq\overline D_{\rm act},
\]
where the right sides are the actual stored rational bounds.
In particular the two native branches are combined before maximizing.
At $t_-=0,h=2$ this gives $\sup A=2$, rather than the weaker
sum of two separate suprema. These inequalities cover every future
finite activation set on the entire unbounded clock domain.

After multiplying by the Euler factor put
\[
 \bar a=P\overline A_{\rm act},\qquad
 \bar b=2P\overline D_{\rm act}+P_1\overline A_{\rm act},\qquad
 \bar g=1-A-\bar a.
\]
The complete global analytic and physical caps at that height are
\begin{align}
 \overline E_{\rm an}
  &=\frac12\left(\frac{R_*+\bar b+\bar aD^+}{\bar g}+T\right),
      \label{src:global-analytic-cap}\\
 \overline E_{\rm ph}
  &=\frac{r_1}{\bar g-r_0}
   +\frac{(B+\bar b)r_0}{2\bar g(\bar g-r_0)}
   +2\cdot10^{-12},\qquad \bar g>r_0.
      \label{src:global-physical-cap}
\end{align}
The gap is explicitly checked; it is not inferred from a favorable
near-$s$ gap.  In particular the Euler derivative term
$P_1\overline A_{\rm act}$ is present.  At the scalar-parent row level,
one valid global debit is the absolute-weight sum of these analytic
and physical caps plus the fixed coordinate debit.  A valid shared
spectral envelope from Section~\ref{src:measure} can instead replace
that analytic sum by its certified nonnegative combination of the
global analytic moments, while retaining the same absolute physical
sum and coordinate charge.  Either fixed branch proves a constant
debit for all $s\geq0$; taking the smaller of two proved constants is
also safe. The directed consumers apply these formulas to the sharp
activation constants just proved, including the full Euler derivative
moment. Conic rows inherit the complete scalar
parent bound with their stated positive multiplier and pay their jet
component separately, as in Section~\ref{src:jets}.

Let $M>0$ be the signed node mass.  The finite source certificates give a rational
$0<\beta<1$ and $R>0$ such that
\[
 N(R)\leq N(0)+\frac{\beta MR}{2},\qquad
 G\leq N(0)+\frac{\beta MR}{2}.
\]
Convexity for $s\leq R$ and the global cap for $s\geq R$ prove
\begin{equation}
 N_{\rm actual}(L)\leq N(0)+\frac{\beta M}{4}(L-L_0),
 \qquad L\geq L_0.
 \label{src:affine-native}
\end{equation}
Subtracting from $ML/4$ leaves the unspent main slope
$(1-\beta)M(L-L_0)/4$.  This is how infinitely many possible future
activation states are covered: each is still a finite native sum, and
its entire added mass is dominated by the positive integral estimates.
No infinite heat-weighted Dirichlet series has been introduced.

The refined errors of Proposition~\ref{src:physical} may safely refund
part of the physical cap.  If $E_0\leq\kappa e_0$ and
$E_1\leq\kappa e_1$ with $0\leq\kappa\leq1$, then
\[
 \frac{\kappa r_1}{g-\kappa r_0}
 +\frac{W\kappa r_0}{2g(g-\kappa r_0)}
 \leq\kappa\left\{\frac{r_1}{g-r_0}
             +\frac{Wr_0}{2g(g-r_0)}\right\}.
\]
The fixed $2\cdot10^{-12}$ and coordinate charges are retained.
For an affine row this refund is taken from the anchor cap and is valid
on both its near and far branch.  It changes neither $\beta$ nor its
paid slope.  The supplied consumer rebuilds all $469$ native-height
caps and all $67$ row refunds from their rational primitives.

\subsection{One common positive-spectrum analytic remainder}
\label{src:measure}

The preceding scalar analytic caps are moments of one measure, which
allows the signed translations to be combined more efficiently.  At the
fixed actual $x,t,N$, let $\nu$ assign absolute $U$-label coefficients
to their frequencies.  Let $\rho$ assign the nonnegative matched head
residual coefficients of \eqref{src:positive-residual}, and absolute
unmatched exterior, reflected, and activated residual coefficients, to
their frequencies.  Both are positive atomic measures attached to the
same finite analytic model.

At the lowest licensed native height $h_0$,
$\int e^{-h_0\lambda/2}\,\dd\nu<1$.  The reciprocal of $Q=1+U$
therefore has an absolutely convergent geometric expansion in the
weighted total-variation norm.  Including the positive measure
$\epsilon$ of omitted Euler prime powers gives the complete majorant
\begin{equation}
 \mu_{x,t}=\frac12\rho*\sum_{m\geq0}\nu^{*m}
                         +\frac12\epsilon.
 \label{src:positive-measure}
\end{equation}
The analytic logarithmic-derivative remainder is the Laplace transform
of a complex measure dominated by $\mu_{x,t}$.  Every residual frequency
is at least $\log4097$, so the same is true after convolution.  All seven
native moment bounds for the row thus concern this one measure:
\[
 \int e^{-h_l\lambda/2}\,\dd\mu_{x,t}(\lambda)
 \leq r_l=\frac12\left(\frac{R_l}{g_l}+T_l\right).
\]
Only the reciprocal of a finite function with a proved gap and convergent
finite-prime geometric tails are summed infinitely here.

For the signed scalar functional define
\[
 F_i(\lambda)=\sum_jw_{ij}e^{-h_{ij}\lambda/2}
                                     \cos(a_{ij}\lambda/2).
\]
Any nonnegative rational numbers $c_l$ satisfying
\begin{equation}
 |F_i(\lambda)|\leq\sum_lc_le^{-h_l\lambda/2}
 \quad\hbox{for every }\lambda\geq\log4096
 \label{src:spectral-envelope}
\end{equation}
give combined analytic charge at most $\sum_lc_lr_l$.
The slightly larger half-line in \eqref{src:spectral-envelope} is safe.
Physical approximation errors are not put into this measure; their
quotient charges continue to have absolute node weights.

For clarity, the continuous proof of the finite envelope is as follows.
Factor out its lowest exponential and form, for both signs,
\[
 D_\pm(\lambda)=\sum_lc_le^{-\beta_l\lambda}
       \pm\sum_jw_{ij}e^{-\alpha_j\lambda}\cos(\delta_j\lambda).
\]
On a rational interval of midpoint $m$ and radius $r$, compute outward
enclosures for $D_\pm(m),D_\pm'(m)$ and an upper bound $B_2$ for
$|D_\pm''|$ throughout it.  The sufficient check is
\[
 D_\pm(m)_{\rm lo}-r|D_\pm'(m)|_{\rm hi}-r^2B_2/2\geq0.
\]
Exact endpoint chaining proves a complete interval cover.  Beyond its
last endpoint, retain the nondecaying positive envelope coefficient and
bound every decaying numerator exponential at that endpoint.  This
proves the entire tail.  Thus the envelope certificates are continuous
inequalities in the spectral variable, not samples of a physical orbit.

The finite thermal head is handled separately.  Its exact signed filter
is
\[
 F_i(n)=\sum_jw_{ij}n^{-h_{ij}/2}
                              \cos(a_{ij}\log n/2),
\]
and its absolute charge is
\begin{equation}
 \frac12\sum_{2\leq n\leq4096}
       d_t(n)n^{-1/2-tL/4}|F_i(n)|.
 \label{src:head-filter}
\end{equation}
The signed node sum is formed before its modulus at each integer.
For the basic chord one first replaces $d_t(n)$ by
$d_{t_+}(n)$ and $L/4$ by $c_0=L_0/4$.  Convexity of the positive
exponential gives, for $\theta=(t-t_-)/(t_+-t_-)$,
\[
 d_t(n)n^{-tc_0}\leq d_{t_+}(n)
 \{(1-\theta)n^{-t_-c_0}+\theta n^{-t_+c_0}\}.
\]
After multiplication by $n^{-1/2}|F_i(n)|/2$ and summation, the two
complete endpoint costs therefore bound the whole thermal cell.
The endpoint certificates round the first of these costs upwards to
its first head field and the second upwards to its last head field.
Any further retained time majorant is checked on its complete declared
cell.  Thus \eqref{src:operational} is a proved chord, not interpolation
of sampled lower bounds.  Later credits described next reduce this
same complete charge, with an explicit disjoint resource ledger.

\subsection{Arithmetic resource allocation and one common physical clock}
\label{src:clock}

This subsection describes the $39$ retained-source improvements entering
the final source file.  It is also useful for distinguishing a genuine
all-center certificate from a numerical calculation at one center.
Let
\[
 c=L/4,\quad c_0=\tfrac12\log N_0,\quad\Delta=c-c_0\geq0.
\]
After the native slope in \eqref{src:affine-native} has been paid, put
$A=(1-\beta)M>0$.  Exact finite coefficient allocation gives
simultaneously, in its declared near resource domain,
\begin{equation}
 J_i\geq B_i^{\rm comp}(t)+A\Delta
                   +\sum_p\alpha_p(1-\cos\phi_p),
 \qquad \alpha_p\geq0,
 \label{src:resource}
\end{equation}
and globally $J_i\geq B_i^{\rm comp}(t)+A\Delta$.
Here the composite credit is already in $B_i^{\rm comp}$.
For $n=\prod p^{e_p}$ an allocated coefficient $r_n$ expends
$r_n\Omega(n)e_p$ of prime $p$'s deficit capacity.  Subtracting every
such expenditure from the original prime capacity defines the remaining
$\alpha_p$.  This finite accounting proves that the composite credit
and the displayed prime deficits coexist; the capacities are never
reset after a credit is taken.  Their lower bounds use the lower thermal
endpoint and attenuation at $L=27$, so they remain valid throughout
the licensed near domain.  The composite allocation spends none of
the retained main slope $A$.

For each integer exponent relation $R_j=\prod_pp^{k_{jp}}$, let
\[
 a_j=2N_0^2\log R_j,\qquad
 v=e^{4\Delta}-1,\qquad q_{\rm clk}=\frac{t}{16N_0^2}.
\]
The gauged phases and an authenticated integer $h_j$ satisfy the exact
identity
\begin{equation}
 \frac{k_j\cdot\phi}{\pi}=2m_j+1+\delta_j(v,q_{\rm clk}),
 \qquad
 \delta_j=a_j-h_j+a_j(v+q_{\rm clk}).
 \label{src:clock-relation}
\end{equation}
The odd parity includes the actual gauges of all primes in the relation.
It follows from the physical clock
$\tau=\pi(2e^L+t/8)$, rather than being imposed on a relaxed independent
phase vector.  If $|\delta_j|<1$, the true sign
$\sigma_j=\operatorname{sgn}(k_j\cdot\phi)$ obeys
\[
 \sigma_j k_j\cdot\phi\geq\pi(1+\sigma_j\delta_j).
\]

For fixed nonnegative rational multipliers $\lambda_j$, set
$b_p=\sum_j\lambda_j\sigma_jk_{jp}$ and use the exact principal-angle
conjugate
\[
 C^*(\alpha,b)=\max_{-\pi\leq u\leq\pi}
                           \{bu-\alpha(1-\cos u)\}.
\]
Then the sum of prime deficits is bounded below by
\[
 \pi\sum_j\lambda_j(1+\sigma_j\delta_j)
                                -\sum_pC^*(\alpha_p,b_p).
\]
The maximum defining each conjugate includes both principal endpoints
and every interior stationary maximum.  In particular it is not replaced
by a local quadratic near a preferred angle.

Choose a total desired angular/main credit $d>0$ and put
\[
 \omega=4d/A<1,\qquad V=\frac{\omega}{1-\omega}.
\]
For every one of the full sign branches certify $|\delta_j|<1$ on the
rectangle
\[
 0\leq v\leq V,\qquad
 \frac{t_-}{16N_0^2}\leq q_{\rm clk}\leq\frac{t_+}{16N_0^2},
\]
and verify $L_0+\log(1+V)<27$.  The whole lower payment is
\begin{equation}
 \mathcal F(v,q_{\rm clk})=
 \pi\sum_j\lambda_j(1+\sigma_j\delta_j(v,q_{\rm clk}))
 -\sum_pC^*(\alpha_p,b_p)+\frac A4\log(1+v).
 \label{src:clock-payment}
\end{equation}
For fixed multipliers it is jointly concave on this rectangle: its
phase part is affine and its only nonlinear term is the displayed
positive multiple of $\log(1+v)$.  Consequently its minimum is at least
the minimum of its four corner values.  The same physical $v,q_{\rm clk}$
is used in every relation before evaluation; separate worst clocks for
individual relations are not substituted.

If all four values are at least $d$ in every sign branch, then
$J_i\geq B_i^{\rm comp}+d$ in the near interval $0\leq\Delta\leq d/A$,
since $e^{4\Delta}-1\leq\omega/(1-\omega)$.  For the complementary
far interval $\Delta\geq d/A$, the global phase-free main bound gives
$A\Delta\geq d$.  The two closed domains cover every physical center
and spend main growth once.

For the seven-relation rows all $128$ signs are checked.  Row $8$ uses
an eighth relation and all $256$ signs.  Its new exact relation is
\[
 \frac{3^4\cdot5\cdot11\cdot19^2\cdot23\cdot47\cdot71}
      {2^4\cdot7^5\cdot13\cdot17\cdot31\cdot67}
       =\frac{123435179505}{123435179504},\qquad h_8=7.
\]
Its total new credit is $11043047769/10^{12}$, replacing the original
$8143047769/10^{12}$; the source improvement is therefore $29/10000$.
The exact remaining $38$ operational decrements are in the supplied
delta table.  They are applied relative to their pinned common parent,
not added to intermediate improvements for the same prime resources.
In particular row $55$'s earlier $5\cdot10^{-6}$ native-zero credit is
removed when its pure-source credit replaces it.  Mixed rows $0,16$
multiply their old-component gain by their positive conic coefficient;
their separately charged derivative component is untouched.

\subsection{Joint derivative rows with one physical logarithm}
\label{src:jets}

For a real center $p>0$ and rational coefficients $b_0,\ldots,b_m$, define
\begin{equation}
 T_bS_t(x,p)=\frac1p\sum_{r=0}^m\frac{b_rp^r}{r!}
                              \partial_p^r S_t(x,p),
 \qquad P_b(z)=\sum_{r=0}^m\frac{b_rz^r}{r!}.
 \label{src:jet-def}
\end{equation}
The three derivative rows have the form
\[
 J_i=\alpha_i\mathcal L_i^{\rm old}S_t+2T_{b_i}S_t(x,p_i),
 \qquad \alpha_i\geq0,\quad b_{i,0}\geq0.
\]
The exact coefficient vectors are included in the plan.  Their orders
are $4,4,8$ for rows $0,16,44$ respectively; $p_0=p_{16}=61/20$ and
$p_{44}=4$.  A single summed kernel, not an addition of separately
proved rates, allocates the available zero-motion force.

Let $G_x(\zeta)=H_t(x+i\zeta)/B_t(x+i\zeta)$, and choose a licensed
circle $|\zeta-p|\leq r$ above a native height $h_0$.  The same finite
Euler gap supplies $|f_x|\geq g/P$ throughout it.  Hence
\[
 W=(G_x-f_x)/f_x,\qquad |W|\leq\delta<1
\]
defines one holomorphic logarithm $\log(1+W)$ on the full disk, bounded
by $\eta=\delta/(1-\delta)$.  It is this logarithm that supplies all
physical derivative errors.  Cauchy's formula gives the safe joint
charge
\begin{equation}
 |T_b(\log(1+W))'(p)|
 \leq\frac{\eta}{pr}\sum_{r'=0}^m
          (r'+1)|b_{r'}|(p/r)^{r'}.
 \label{src:jet-physical}
\end{equation}
The earlier Parseval version with the square root of the corresponding
sum of squares is also valid where its certificate uses it.  Neither
formula assigns independent errors to the derivatives.

The baseline quotient ledger implies
$e_0P/g\leq E_{\rm ph,old}(h_0)/100$.  Proposition~\ref{src:physical}
therefore permits
\begin{equation}
 \delta\leq
 \min\{1,E_0(h_0,t_-)/e_0\}\cdot
                         E_{\rm ph,old}(h_0)/100.
 \label{src:jet-refined-ratio}
\end{equation}
Indeed the direct new bound is
$|(G_x-f_x)/f_x|\leq E_0P/g$; no extra factor of $P$ is inserted.
No $E_1$ estimate is needed for this value-log route.  The normaliser is
still present: on all these circles
\[
 \left|\partial_\zeta\log B_t(x+i\zeta)
                       -(L/4-i\pi/8)\right|<10^{-10}.
\]
Its separately paid charge is
$(10^{-10}/p)\sum_{r'}|b_{r'}|(p/r)^{r'}$.

For the analytic remainder, use the same measure
\eqref{src:positive-measure}.  If $h_0<p$ and
$a=(p-h_0)/2$, its complete charge is bounded by
\begin{equation}
 \frac{E_{\rm an}(h_0)}p
 \sup_{\lambda\geq\log4097}
       |P_b(-p\lambda/2)|e^{-a\lambda}.
 \label{src:jet-analytic}
\end{equation}
A monomial-by-monomial bound follows from the maximum of
$\lambda^{r'}e^{-a\lambda}$ at $\max(\log4097,r'/a)$.
The selected row-$44$ bound instead keeps the whole polynomial before
absolute value.  It evaluates its endpoint and every zero on the
half-line of
\[
 \frac{\dd}{\dd\lambda}P_b(-p\lambda/2)
                                      -aP_b(-p\lambda/2).
\]
An exact rational Sturm sequence proves that the disjoint root
isolators exhaust the half-line; directed interval evaluation bounds
each critical value.  The endpoint is rounded down below $\log4097$,
which enlarges the domain, and the value tends to zero at infinity.
Zeros of the polynomial itself contribute zero.  This proves the full
semi-infinite envelope, with no sampled-maximum assumption.

The finite derivative head is charged as
\begin{equation}
 \frac1{2p}\sum_{2\leq n\leq4096}d_{t_+}(n)
 n^{-1/2-t_-c_0-p/2}|P_b(-p\log n/2)|,
 \label{src:jet-head}
\end{equation}
with the appropriate end-time attenuation for each endpoint chord.
All $4095$ coefficient enclosures are reconstructed by the
proper-divisor recurrence \eqref{src:triangular}.  The main term is
$b_0c_0/p$, frozen downwards only because $b_0\geq0$.

Combining the main term with
\eqref{src:jet-physical}, \eqref{src:jet-analytic},
\eqref{src:jet-head}, and the separate normaliser charge gives the
new jet source.  Every old source charge is multiplied by $\alpha_i$,
and every jet source field by $2$.  Thus the operational source fields
are exactly $\alpha_i$ times the complete old fields plus twice the new
fields.  The full old record is retained in each certificate; an old
diagnostic split is not used as a decomposition of the conic source.

The selected row-$16$ circle uses native floor $51/20$ and radius $1/2$;
row $44$ uses $51/25$ and radius $39/20$.  Each circle independently
passes $p-r\geq h_0$, $p+r\leq6$, $r\leq39/20$, $p+r<7$, and the
thermal guard.  The independent source consumers reconstruct $4128$
gates for mixed row $0$, $8363$ for the later conic extensions, and
$4215$ for the final order-eight row.  These counts describe finite
proof obligations; they do not replace the preceding analytic argument.

\subsection{Continuous signed kernel and the entire own pair}
\label{src:geometry}

For an upper/lower pair at horizontal spacing $d$ and height $b$, define
\[
 K_h(d,b)=\frac{h-b}{d^2+(h-b)^2}
                         +\frac{h+b}{d^2+(h+b)^2}.
\]
Apply the same row functional to this kernel, including all signed nodes
and all derivatives for a joint row, and denote the result by
$R_i(d,b)$.  Its own-pair contribution at squared height $q=y^2$ is
\begin{equation}
 O_{\mathrm{pair},i}(q)=R_i(0,y).
 \label{src:own-def}
\end{equation}
For an ordinary row this is the explicit rational function
\begin{equation}
 O_{\mathrm{pair},i}(q)=\sum_j
 \frac{2w_{ij}h_{ij}(h_{ij}^2+a_{ij}^2-q)}
      {(h_{ij}^2+a_{ij}^2-q)^2+4a_{ij}^2q}.
 \label{src:own}
\end{equation}
For a conic row add
\[
 2\sum_{r=0}^m\frac{b_rp^{r-1}}{r!}
             \partial_p^r\frac{2p}{p^2-q}
\]
to $\alpha_i$ times the old function.  The signs stay inside the entire
sum; neither individual own terms nor their derivatives are presumed
positive.

\begin{theorem}[Complete continuous kernel certificates]
\label{src:kernel}
For every row of the pinned plan and every $q$ in its closed declared
contact interval, $h_{ij}^2>q$ at every node and
\begin{equation}
 R_i(d,\sqrt q)\leq\frac4{d^2+4q}
             \qquad(d\in\mathbb R,\ (d,q)\ne(0,0)).
 \label{src:boundary-kernel}
\end{equation}
The corresponding inequality at every other-zero height
$0\leq b\leq y=\sqrt q$ is
\begin{equation}
 K_y(d,b)\geq\frac y2R_i(d,b),
 \label{src:whole-kernel}
\end{equation}
away from the distinguished singular pair.  Each complete $O_{\mathrm{pair},i}$ is
nondecreasing on its declared contact interval.
\end{theorem}

\begin{proof}
The finite part of the proof is an exact rational polynomial
certificate, regenerated directly from the displayed Poisson
fractions and the plan.  Symmetric shifts make the rational expression
even in $d$.  Put $D=d^2$ and, away from $(D,q)=(0,0)$, clear all positive denominators in
$4/(D+4q)-R_i(d,\sqrt q)$.  Their positivity follows from
$h_{ij}^2>q$; for a derivative row the corresponding powers have the
same sign.  Compactify $D=z/(1-z)$ and multiply by the required
positive power of $1-z$.  The endpoint $z=1$ represents the leading
far-field coefficient and is included.  The result is a rational
polynomial on the compact rectangle in $(z,q)$.

On every rational subdivision cell the checker converts this
polynomial to its tensor-product Bernstein basis.  If its local degree
is $(m,n)$, the basis functions are
\[
 \binom{m}{r}u^r(1-u)^{m-r}
 \binom{n}{s}v^s(1-v)^{n-s}\geq0
 \qquad(0\leq u,v\leq1).
\]
The stored rational coefficients are all checked to be nonnegative,
and the checker verifies the complete subdivision tree and its
boundary ownership.  Thus their sum proves the inequality at every
point of every cell.  It independently regenerates the cleared
polynomial, so an irrelevant nonnegative coefficient table would not
pass.  At $(d,q)=(0,0)$ the target fraction is not defined; the
denominator-cleared polynomial is still checked on that closed corner.
Actual nonreal contacts have $q>0$.  This proves
\eqref{src:boundary-kernel}, including $q$ and
spacing continua and the far-field endpoint.

The same procedure applied to the rational derivative $O_{\mathrm{pair},i}'(q)$
proves aggregate monotonicity.  For example the final row-$44$
boundary polynomial has degree $118$ in each of its two variables
and $29$ certified cells; its own-derivative polynomial has degree
$116$ and one cell.  Degree size does not add an analytic assumption:
all its arithmetic is rational.

Finally, as a function of $(d,b)$ each probe kernel is harmonic
throughout $|b|\leq y$, since its poles are outside the strip.
The same holds after real translations and differentiation in a probe
height.  The positive strip-Poisson formula extends the nonnegative
boundary defect to the whole interior.  The target kernel carries a
nonnegative boundary atom at the distinguished point, which only
increases the interior defect.  This proves
\eqref{src:whole-kernel}.  The explicit extension, convergence of its
zero sums, and treatment of the distinguished pair are given in the
analytic foundations section.
\end{proof}

\subsection{Closing the finite source assertion}
\label{src:close}

Here is the exact dependency order in Theorem~\ref{src:source}.
First, the corrected physical calculation proves the common finite
model and the absolute errors on every actual source disk.  Finite
thermal recurrences and complete native labels give a denominator gap,
the matched residual, and full future-activation charges.  The common
analytic measure permits only the explicitly certified spectral
recombinations; physical quotient charges remain separate.  The finite
head is summed with signs before moduli.  The affine native bound
retains an authenticated portion of the main growth.  Disjoint
composite capacities and the complete physical clock branches yield
the licensed finite head decrements.  Selected derivative rows are
charged by one analytic physical logarithm and a complete spectral
envelope.  Each conic row uses one combined source and kernel.

The final source-assembly consumer reconstructs the M8 parent from its
physical and analytic primitives and then applies the $39$ globally
proved pure-source gains to \texttt{first\_hi} and \texttt{last\_hi}.
Every other operational field is preserved. In row~$55$, an earlier
local credit is replaced, rather than retained in addition to the new
global credit. The direct M8-to-final reconstruction therefore has no
local native exclusion as a premise. Rows $0,16$ carry the exact positive
conic scaling of their scalar parent. Whole-object comparisons establish
that the resulting rational fields are precisely those in
\eqref{src:operational}. Together with the estimates above, this proves
Theorem~\ref{src:source}.

The finite supplement contains the raw rational primitives, branch
multipliers, phase factorizations, interval envelopes, conic coefficients,
kernel subdivision trees, source assembly and independent consumers.
Checks of file digests identify those objects but do not prove the
analytic inequalities.  Conversely, a report saying that a numerical
optimizer passed is not one of the source premises: the consumers
reconstruct the declared finite inequalities using exact rational
arithmetic or outward Arb enclosures.  The comparison section uses
only Theorems~\ref{src:source} and~\ref{src:kernel} for these pinned
versions.

% END INCLUDED FILE


% BEGIN INCLUDED FILE: sections/flow.tex
\section{An exact positive comparison and its complete terminal cleanup}
\label{flow:section}

This section specifies the finite calculation that turns the analytic
source inequalities into an upper bound for $\Lam$.  The inputs are the
rational plan, the rational source ledger, and the list of rational arcs.
The list is a certificate: the optimization and integration procedures
that found it are unnecessary for its verification.  In particular,
rounding toward a larger comparison height is permitted only after the
whole preceding ideal arc has been shown to remain in its licensed domain.
The endpoint below includes the additional time needed to remove the
\emph{entire} terminal strip.

\subsection{The interface with the analytic estimates}
\label{flow:interface}

Write $q=y^2$ for the square of the height of a highest nonreal zero.
There are $67$ rows, indexed by $i=0,\ldots,66$.  Row $i$ supplies a
closed time cell $[a_i,b_i]$, a contact cap $c_i>0$, a signed probe
functional, and rational constants
\[
 M_i,\qquad N_i,\qquad A_i\geq B_i>0.
\]
The cells meet at their endpoints.  They cover $[0,87/800]$; the
comparison will stop before the right endpoint of the last cell.  Put
\begin{equation}
 \theta_i(t)=\frac{t-a_i}{b_i-a_i},\qquad
 S_i(t)=M_i-N_i-(1-\theta_i(t))A_i-\theta_i(t)B_i.
 \label{flow:source-affine}
\end{equation}
The meanings of the four stored fields are respectively
\texttt{main\_lo}, \texttt{native\_debit}, \texttt{first\_hi}, and
\texttt{last\_hi}.  Thus $S_i$ is a lower bound, rather than a numerical
estimate of the actual source.  Since $A_i\geq B_i$, it is nondecreasing
on its entire cell.

For reference, the exact own-pair function for the retained translated-probe
part of a row is the following rational function.  If the row has nodes
$(h_j,v_j,w_j)$, with signed $w_j$, define
\begin{equation}
 O_{\mathrm{pair},i}^{\rm old}(q)=
 \sum_j\frac{2w_jh_j(h_j^2+v_j^2-q)}
 {(h_j^2+v_j^2-q)^2+4v_j^2q}.
 \label{flow:old-own}
\end{equation}
The signs in this sum are retained.  Replacing them separately by absolute
values would change the functional and invalidate the numerical rate.
All node poles lie outside the complete contact cap: $h_j^2>c_i$.

Rows $0$, $16$, and $44$ also have joint-jet terms.  A row of this form
has rational $\alpha\geq0$, centres $p$, and coefficient vectors
$(b_0,\ldots,b_m)$.  In the units of the stored source its whole own pair is
\begin{equation}
 O_{\mathrm{pair},i}(q)=\alpha O_{\mathrm{pair},i}^{\rm old}(q)+
 2\sum_{\text{centres }p}\sum_{m=0}^{d_p}
 b_m p^{m-1}\frac{1}{m!}
 \frac{\partial^m}{\partial p^m}\left(\frac{2p}{p^2-q}\right).
 \label{flow:jet-own}
\end{equation}
For the other rows $O_{\mathrm{pair},i}=O_{\mathrm{pair},i}^{\rm old}$.  The poles satisfy $p^2>c_i$.
The geometry certificate proves that the \emph{whole signed} $O_{\mathrm{pair},i}$ is
nondecreasing on $0\leq q\leq c_i$.  This is not inferred from the
monotonicity of individual signed summands.  The same coefficient vector
and the same scaling occur in the source and kernel certificates.

The source theorem (Theorem~\ref{src:source}) and kernel theorem
(Theorem~\ref{src:kernel}), together
with the all-complex first-contact argument in the analytic foundations,
give the following interface at a simple highest contact in the exterior
region:
\begin{equation}
 q'(t)\leq-2-2q(t)\bigl(S_i(t)-O_{\mathrm{pair},i}(q(t))\bigr),
 \quad t\in[a_i,b_i],\quad 0<q(t)\leq c_i.
 \label{flow:differential-interface}
\end{equation}
The finite head is covered by its separate wall theorem.  Compactness
and the treatment of multiple contacts are also supplied by the analytic
foundations.  Equation~\eqref{flow:differential-interface} is the only
analytic zero-motion input to the computation here; the finite consumer
cannot establish it by examining an arc file.

\subsection{A comparison arc with an upward reset}
\label{flow:arc}

An arc record consists of
\[
 (i,t_0,\delta,q_0,q_1,s,o,r,\ell),
\]
all rational except the row index $i$.  Here $q_0$ is the stored initial
square, $s$ a source lower bound, $o$ an own-pair upper bound, and $r>0$
the comparison rate.  The ideal arc is
\begin{equation}
 Q(u)=\left(q_0+\frac2r\right)e^{-ru}-\frac2r,
 \qquad 0\leq u\leq\delta.
 \label{flow:ideal}
\end{equation}
It solves $Q'=-2-rQ$ and $Q(0)=q_0$.  Its outgoing ideal value need
not be rational.  The record instead stores an outward rational endpoint
$q_1\geq Q(\delta)$, which becomes the next arc's initial value.
The following finite conditions are enough:
\begin{align}
 &a_i\leq t_0<t_0+\delta\leq b_i,
       &&0<q_0\leq c_i,\label{flow:arc-domain}\\
 &s\leq S_i(t_0),\qquad o\geq O_{\mathrm{pair},i}(q_0),
       &&0<r\leq2(s-o)-\eta,\label{flow:arc-rate}\\
 &\ell=q_0-(2+rq_0)\delta>0,
       &&Q(\delta)\leq q_1<q_0,
 \label{flow:arc-positive}
\end{align}
where the fixed strictness allowance is
\begin{equation}
 \eta=10^{-12}.
 \label{flow:eta}
\end{equation}
Each $q_1$ has denominator a power of two at most $2^{96}$.
The allowance $\eta$ has the units of the rate in
\eqref{flow:arc-rate}; it is not a lower bound for the source itself.

\begin{lemma}[A positive ideal arc]
\label{flow:positive-arc}
If $q_0>0$, $r>0$, and $\ell$ in
\eqref{flow:arc-positive} is positive, then on $0\leq u\leq\delta$
\[
 0<\ell\leq Q(u)\leq q_0,
 \qquad Q'(u)<0.
\]
\end{lemma}
\begin{proof}
The elementary inequality $e^{-ru}\geq1-ru$ gives
\[
 Q(u)\geq q_0-(2+rq_0)u\geq\ell>0.
\]
Consequently $Q'=-2-rQ<0$, and $Q(u)\leq Q(0)=q_0$.
This proof controls the entire closed ideal arc.  Positivity of the stored
outgoing value alone would not prove the assertion.
\end{proof}

\begin{lemma}[Strict exclusion of a simple contact on an arc]
\label{flow:strict-arc}
Assume \eqref{flow:differential-interface} and the whole-row own-pair
monotonicity.  A zero branch starting below $Q$ cannot first meet $Q$ at
a simple highest contact on the arc satisfying
\eqref{flow:arc-domain}--\eqref{flow:arc-positive}.
\end{lemma}
\begin{proof}
At such a contact, $q=Q(u)>0$ and $q\leq q_0\leq c_i$.  Since $S_i$
is nondecreasing and $O_{\mathrm{pair},i}$ is nondecreasing on the cap,
\[
 S_i(t_0+u)-O_{\mathrm{pair},i}(q)\geq s-o\geq\frac{r+\eta}{2}.
\]
Thus \eqref{flow:differential-interface} yields
\[
 q'\leq-2-(r+\eta)q
     < -2-rq=Q'(u),
\]
contrary to a first meeting from below.  The strict inequality is valid
at every positive contact height, however small it is.  Lemma
\ref{flow:positive-arc} guarantees positivity through the incoming end.
\end{proof}

At $t_0+\delta$ the replacement of $Q(\delta)$ by $q_1$ is an
\emph{upward} jump of the comparison, so it preserves a previously
established strict separation from all zeros.  A terminal contact at the
incoming ideal value is excluded by the same strict differential
inequality.  The next cell may have a different signed functional; that
change is legal because the next record checks its own initial cap and
its original absolute time cell.  No continuous interpolation of the
functionals is required.

\subsection{Exact exponential enclosures and rational own-pair evaluation}
\label{flow:exact-arithmetic}

For an exact rational $x\geq0$ put
\begin{equation}
 P(x)=1+\frac{x}{2}+\frac{5x^2}{44}+\frac{x^3}{66}
      +\frac{x^4}{792}+\frac{x^5}{15840}+\frac{x^6}{665280}.
 \label{flow:p20}
\end{equation}
All coefficients are positive. Repeated integration by parts gives
\[
 e^{-x}P(x)-P(-x)
 =-\frac{x^{13}}{12!}\int_0^1 e^{-tx}t^6(1-t)^6\,\dd t\leq0.
\]
Thus the even diagonal $[6/6]$ Pad\'e expression is an exact upper bound,
\begin{equation}
 e^{-x}\leq \frac{P(-x)}{P(x)}.
 \label{flow:reciprocal-exp}
\end{equation}
The consumer evaluates the two polynomials with exact rational Horner
arithmetic and checks
\begin{equation}
 \left(q_0+\frac2r\right)\frac{P(-r\delta)}{P(r\delta)}-\frac2r
 \leq q_1.
 \label{flow:exp-check}
\end{equation}
Together with \eqref{flow:reciprocal-exp}, this establishes the outgoing
inequality in \eqref{flow:arc-positive}. No floating approximation to an
exponential and no trust in the endpoint producer is needed.

The joint derivative terms also admit an independent rational evaluation.
Formally write $q=y^2$ and set
\[
 u=\frac1{p-y},\quad v=\frac1{p+y},\quad
 A=u+v=\frac{2p}{p^2-q},\quad B=uv=\frac1{p^2-q}.
\]
The symmetric power sums $S_j=u^j+v^j$ are rational in $p,q$ and obey
\[
 S_0=2,\qquad S_1=A,\qquad S_{j+1}=AS_j-BS_{j-1}.
\]
Consequently
\begin{equation}
 \frac1{m!}\frac{\partial^m}{\partial p^m}\frac{2p}{p^2-q}
 =(-1)^m S_{m+1}.
 \label{flow:jet-binomial}
\end{equation}
The fresh consumer uses this Newton recurrence, independently of the
power-series derivative calculation used to construct the candidate
record. The identity is exact for all rational inputs in the licensed
domain $p^2>q$. Every signed term of the whole own pair is added before
testing its proposed upper bound.

\subsection{Initialization at the ordinary seven-eighths strip}
\label{flow:startup}

The ordinary zeta theorem used in the analytic foundations excludes
nontrivial zeros with real part greater than $7/8$.  The functional
equation gives the reflected bound.  Under the coordinate conversion
between zeta and $H_0$, the initial imaginary halfwidth is therefore
$2(7/8)-1=3/4$, and its square is $9/16$.  The comparison is initialized
strictly outside that strip:
\begin{equation}
 q_{\rm init}=\frac9{16}+10^{-12}
             =\frac{562500000001}{10^{12}}.
 \label{flow:initial}
\end{equation}
The cushion is part of the exact computation.  It is not permission to
use a narrower unproved initial zeta plane.

To enter the positive-time compact first-contact setting, take
\begin{equation}
 \tau=10^{-14}.
 \label{flow:tau}
\end{equation}
Classical de Bruijn strip contraction, applied from $t=0$, bounds every
actual squared height at $t\leq\tau$ by $9/16-2t$.  On the first ideal
arc, Lemma~\ref{flow:positive-arc} gives
\[
 Q(t)\geq q_{\rm init}-(2+r_0q_{\rm init})t.
\]
The separation from the classical upper strip is consequently at least
\begin{equation}
 Q(t)-\left(\frac9{16}-2t\right)
 \geq10^{-12}-r_0q_{\rm init}t
 \geq10^{-12}-r_0q_{\rm init}\tau.
 \label{flow:startup-gap}
\end{equation}
Here the first rate in the exact certificate is
\[
 r_0=\frac{122615339047817575291909362361}
 {9903520314283042199192993792}.
\]
The last expression in \eqref{flow:startup-gap} is exactly
\begin{equation}
 \frac{921380903213784218478552545296645590637639}
 {990352031428304219919299379200000000000000000000000000}
 >\frac1{2\cdot10^{12}}.
 \label{flow:startup-exact}
\end{equation}
It is approximately $9.3035695790\cdot10^{-13}$.  Also
$0<\tau<1/2048000$, so it lies strictly inside the first arc.
This proves a strict starting separation at an actual positive time;
it does not assume positive-time simplicity of every zero at $t=0$.
Theorem~\ref{fnd:comparison} justifies passing from this separation to the
first-contact argument on $[\tau,T]$.

\subsection{The finite certificate and its coverage}
\label{flow:finite-certificate}

The exact computation consists of $221010$ arcs.  Every arc has
\begin{equation}
 \delta=\frac1{2048000}.
 \label{flow:step}
\end{equation}
The first $20$ time rows have length $1/400$ and therefore contain
$5120$ arcs each.  Rows $20$ through $65$ have length $1/800$ and
contain $2560$ arcs each.  The calculation uses $850$ arcs of row $66$.
In particular,
\[
 20\cdot5120+46\cdot2560+850=221010.
\]
The finer decomposition is data, rather than an inference from a plot
or a numerical time mesh: exact stored times and states are checked at
every join.

The following table gives several checkpoints.  The height columns are
decimal summaries for orientation; the certificate stores their exact
rational values, and none of the displayed rounded entries is used as
an inequality in the proof.
\begin{center}
\small
\begin{tabular}{rccrrr}
\hline
Row & Time cell & Arcs & Entry square & Exit square & Number of nodes\\
\hline
0  & $[0,1/400]$ &5120& .562500000001& .540389791362&9\\
1  & $[1/400,1/200]$ &5120& .540389791362& .518757332785&9\\
16 & $[1/25,17/400]$ &5120& .263984448758& .249128464055&9\\
44 & $[2/25,13/160]$ &2560& .075504237976& .071222671787&25\\
55 & $[3/32,19/200]$ &2560& .033261081999& .029931516439&10\\
65 & $[17/160,43/400]$ &2560& .003496950134& .000934680690&10\\
66 & $[43/400,87/800]$ &850& .000934680690& .000099742033&9\\
\hline
\end{tabular}
\end{center}
For row $66$, the exit is at the stopping time, before the right
endpoint of the displayed licensed cell.  In all rows the argument stays
below the cap supplied by the plan; it is not necessary that the cap
equal the comparison's entry value.

The minimum exact entry-cap clearance, over the complete list, is
\begin{equation}
 \min_{\text{arcs}}(c_i-q_0)=\frac{21059999}{10^{12}}>0.
 \label{flow:cap-clearance}
\end{equation}
The smallest comparison rate occurs at the first arc and is $r_0$
above, approximately $12.3809852615$.  The rates range up to approximately
$22.6946264088$.  These are the rates $r$ in $Q'=-2-rQ$, rather than
values of the undoubled source $S_i-O_{\mathrm{pair},i}$.

The checker additionally verifies that all $2560$ closed ideal arcs of
row $55$, including their outward jumps, remain inside
\[
 \frac{289}{10000}\leq Q\leq\frac{361}{10000}.
\]
This retained tube check is conservative.  The old native improvement
associated with that tube has already been replaced in the final source
ledger and is not credited a second time.  Likewise, every source
change is compared against its specified parent in its original
functional units; a flow-only verifier must not silently rescale a source
or widen a cap.

\subsection{A fully specified first record}
\label{flow:first-record}

The first arc makes the normalisation and directions of all roundings
concrete.  It belongs to row $0$, whose licensed cap is
\[
 c_0=\frac{28126053}{50000000}=0.56252106.
\]
Its exact inputs are $t_0=0$, $\delta=1/2048000$, and
$q_0=562500000001/10^{12}$.  The stored source and own bounds are
\begin{align*}
 s_0&=\frac{18915863536666823719831936483}
 {2475880078570760549798248448},\\
 o_0&=\frac{114846276982028443785727385843}
 {79228162514264337593543950336}.
\end{align*}
They are approximately $7.64005643908$ and $1.44956380834$,
respectively.  The difference between twice their difference and the
stored rate is exactly
\[
 2(s_0-o_0)-r_0=
 \frac{39614081257132169}{39614081257132168796771975168}
 >10^{-12}.
\]
Thus the strictness allowance is paid before the comparison is used.
The record's positive linear lower bound is
\[
 \ell_0=
 \frac{11408766623905503936988141385986459896876653639}
 {20282409603651670423947251286016000000000000000}>0,
\]
and its stored outgoing value is
\[
 q_1=\frac{5570686828209889305941576421}
 {9903520314283042199192993792}.
\]
These are approximately $0.5624956228993352$ and
$0.5624956229125658$.  The two inequalities are used for different
purposes: $\ell_0>0$ bounds the entire ideal arc from below, while
\eqref{flow:exp-check} bounds its final value from above by $q_1$.
The initial cap clearance is
$c_0-q_0=21059999/10^{12}$, which also equals the smallest clearance
in the entire certificate.  Every later record is verified by precisely
the same rational conditions; it need not have the same source or
functional as this example.

\subsection{What the complete consumer proves}
\label{flow:consumer}

A successful numerical verification establishes the following finite
statement, independently of the method used to construct the records.

\begin{proposition}[Finite comparison certificate]
\label{flow:certificate-proposition}
For the pinned plan and source ledger, the $221010$ records start at
\eqref{flow:initial}, satisfy
\eqref{flow:arc-domain}--\eqref{flow:arc-positive} with
$\eta=10^{-12}$, meet exactly in stored time and outward state, and end
at the rational pair $(T,q_T)$ in
\eqref{flow:terminal-pair} below.  The strict startup inequality
\eqref{flow:startup-exact} holds.  Every ideal arc lies within a licensed
source time cell and contact cap.
\end{proposition}

Here is a complete logical outline of the finite consumer.  The
pseudocode uses exact integer fractions throughout.
\begin{enumerate}
\item Check the input hashes, schemas, and the $67$ row identities.
For every row check its absolute time interval, common source/kernel cap,
source-functional scaling, and the separate source and geometry proof
bindings.  In particular, verify $A_i\geq B_i>0$ and poles outside the
full cap.  Check the strict far-field condition needed by the harmonic
extension in the analytic kernel proof.
\item Set $(t,q)=(0,q_{\rm init})$.  For each arc check that its incoming
$(t_0,q_0)$ equals the current exact pair, its row is the current or next
row, and \eqref{flow:arc-domain} holds.
\item Evaluate \eqref{flow:source-affine} at $t_0$ and
\eqref{flow:old-own}--\eqref{flow:jet-own} at $q_0$ by
\eqref{flow:jet-binomial}.  Verify the downward source and upward own
bounds and the strict rate inequality \eqref{flow:arc-rate}.
\item Check the exact positive lower bound $\ell$ and the outgoing
exponential enclosure \eqref{flow:exp-check}.  Check $q_1<q_0$ and the
$96$-bit dyadic denominator.  Where recorded, check the complete closed
tube as well.  Update $(t,q)=(t_0+\delta,q_1)$.
\item Verify that the resulting $(t,q)$ equals the stated terminal pair,
that $0<q\leq10^{-4}$, and that all stated time horizons hold.  Check
\eqref{flow:startup-exact}, compute $t+q/2$ exactly, and round that sum
outward to the claimed decimal.
\end{enumerate}
The consumer \path{verification/arcs/verify_arcs_fresh.py} evaluates
all of these inequalities directly with arbitrary-precision rational
arithmetic. Its execution performs $5\,797\,314$ finite checks on all
$221\,010$ arcs. The count records the scope of an execution; the proof
consists of the inequalities and their analytic implications.

For each arc, the supplied compact transcript gives six exact $128$-bit
dyadic lower bounds: source margin, own-pair margin, strict rate margin,
outgoing-state margin, positive linear floor, and cap clearance.
The first four are nonnegative and the final two strictly positive.
These actual witnesses are in
\path{verification/arcs/COMPACT_ALL_ARC_MARGIN_LOWER_BOUNDS.jsonl.gz}.
They permit direct inspection of the inequalities and are not a substitute
for recomputing the arc from its rational data.

The exact numerical inputs have the following SHA-256 identities.  Paths
are relative to the evidence workspace in the proof packet.
\begin{quote}\small\ttfamily\raggedright
lambda\_round14/material/m8\_late/M8\_LATE\_PLAN.json\\
\seqsplit{ec346b7c1136eda38021eb6323033b3b48ed21e6f7458eeeac7cc9579db2497a}\\[3pt]
lambda\_round15/flow\_main\_v2/MAIN\_CLOCK\_SOURCE.json\\
\seqsplit{b59905d7b994d3424ba14f9347be75a2429ae4d4b0de2c2404971ffda1d4d536}\\[3pt]
lambda\_round16/ordinary\_7\_8\_flow/full\_flow/FLOW.json.gz\\
\seqsplit{a664e829a3681d96e474bcf875d80cba7b16fe8f65e050c69f01b5ea07de4737}
\end{quote}
The compressed flow identity is checked before decompression.  The
producer identity is also bound, but correctness does not depend on
trusting the producer: the output inequalities are re-evaluated.

\subsection{Terminal values and the upper bound}
\label{flow:terminal}

The exact terminal pair is
\begin{equation}
 T=\frac{22101}{204800}=0.1079150390625,
 \qquad
 q_T=\frac{7902377963830536747158643}{79228162514264337593543950336}.
 \label{flow:terminal-pair}
\end{equation}
Here $q_T$ is approximately $0.00009974203254313493$.  Both source and
wall time restrictions are respected: $T<19/125$, and even the eventual
cleanup endpoint is below $41/250$.  Only the source/wall bounds through
$T$ are actually needed; after $T$, the classical contraction theorem is
applied to the global strip already obtained.

\begin{proposition}[Assembly of the global terminal strip]
\label{flow:global-terminal}
Assume the ordinary seven-eighths theorem and the analytic source,
whole-kernel, finite-head, compactness, and multiple-contact results
proved in the preceding sections.  Then every zero of $H_T$ has
imaginary part at most $\sqrt{q_T}$ in absolute value.
\end{proposition}
\begin{proof}
Classical contraction and \eqref{flow:startup-gap} give strict separation
at $\tau$.  Suppose that a zero first reaches the piecewise ideal
comparison before its terminal time.  Positive-time remote-tail
compactness produces an attained highest contact; the finite-head wall
puts a nonreal contact in the exterior region covered by the source
estimates.  A multiple contact is excluded by the backward complex
Hermite splitting argument in the analytic foundations.  A simple one
contradicts Lemma~\ref{flow:strict-arc}.  All ideal arcs are positive and
within their licensed caps by Proposition
\ref{flow:certificate-proposition}.  Upward resets preserve the
separation, including time-row joins.  Hence there is no first contact,
and at time $T$ all zeros lie within the outward terminal strip.
\end{proof}

\begin{corollary}[Exact comparison-plus-cleanup endpoint]
\label{flow:endpoint}
The de Bruijn--Newman constant satisfies
\begin{equation}
 \Lam\leq T+\frac{q_T}{2}
 =\frac{427693072077942473661578098491}
 {3961408125713216879677197516800}
 <\frac{107964910079}{10^{12}}.
 \label{flow:headline}
\end{equation}
\end{corollary}
\begin{proof}
De Bruijn contraction applied at time $T$ to the complete global strip
$|\Imm z|\leq\sqrt{q_T}$ makes every zero real by time $T+q_T/2$.
The rational addition in \eqref{flow:headline} is exact.  Finally,
\[
 T+q_T/2\approx0.10796491007877157,
\]
and the displayed fraction with denominator $10^{12}$ is its upward
decimal ceiling.  In particular the full cleanup has been paid, rather
than omitting $q_T/2$ or treating a small positive terminal height as a
real-zero conclusion.
\end{proof}

The computation is initialized at exactly $7/8$.  An earlier, numerically
smaller endpoint used a narrower initial zeta strip and cannot be
relabelled as this theorem.  Similarly, the local fixed-head confluent
certificates and hypothetical uniform-current comparisons play no role
in Proposition~\ref{flow:global-terminal} or Corollary
\ref{flow:endpoint}.  The bound here uses the proved row source ledger
and its complete finite comparison.

% END INCLUDED FILE


% BEGIN INCLUDED FILE: sections/verification.tex
\section{Certificate interfaces and reproducibility}\label{ver:section}

The proof has three distinct verification layers.  First, it uses ordinary
mathematical theorems, including a particular externally verified formal
nonvanishing theorem.  Second, it uses the analytic estimates proved in this
manuscript, with their domains and normalisations stated explicitly.  Third,
it uses finite rational or outward-rounded interval certificates.  A successful
numerical producer is not an additional mathematical premise: its output is
used only through the inequalities checked by the corresponding consumer.
This section specifies these interfaces and the evidence retained for them.

\subsection{The imported ordinary nonvanishing theorem}\label{ver:formal}

The only new initial-strip input is the ordinary statement
\begin{equation}\label{ver:zeta-input}
 \forall s\in\mathbb C\setminus\{1\},\qquad
 \Ree s>\frac78\quad\Longrightarrow\quad\zeta(s)\ne0.
\end{equation}
The formal source \cite{SevenEighths} accompanying the preprint
\cite{OAI:The-Quasi-Riemann-Hypothesis-September-30-2026} exports the Lean declaration
\begin{quote}\small\ttfamily
OAI.riemannZeta\_ne\_zero\_of\_seven\_eighths\_lt\_re.
\end{quote}
Its type has the variable $s:\mathrm{Complex}$ and the single assumption
$(7/8:\mathrm{Real})<s.\mathrm{re}$; the conclusion concerns Mathlib's
\texttt{riemannZeta}.  In particular, it does not accept a source certificate,
an additional zero-free strip, a moment estimate, or a detector inequality as
an argument.  The certificate used by its final assembly is constructed in
the imported development.  The usual completed-zeta functional equation,
conjugation and zero correspondence then give precisely
$|\Imm z|\le3/4$ for zeros of $H_0$; no strictly smaller halfwidth is imported.
The convention assigning a finite value to the meromorphic zeta function at
its pole does not affect any zero of the completed function.

The retained source and build identities are as follows.  Long hexadecimal
identifiers below are Git commit or tree identifiers, rather than numerical
error bounds.
\begin{center}\small
\begin{tabular}{p{.29\linewidth}p{.64\linewidth}}
\hline
Object & Pinned identity\\\hline
Official source commit &
\texttt{\seqsplit{fd4aeeb2ee4fc729c18d98444fed42fd0529eeeb}}\\
External rebuild source commit &
\texttt{\seqsplit{adc7f1241b42e322a6451854ab7e4b4c146bf78a}}\\
Shared proof subtree &
\texttt{\seqsplit{1e340710a6e6220842ad878bb89a1de4b1989af4}}\\
Mathlib commit &
\texttt{\seqsplit{d13f23b723b8a846827a245b89c10fc7d3f11612}}\\
External verification report commit &
\texttt{\seqsplit{b83a73c6cb626a391e551baafdb7d0cceafaa55a}}\\\hline
\end{tabular}
\end{center}
The proof subtree, toolchain, dependency manifest and relevant compatibility
patches agree at the two official-source commits.  The downloaded ordinary
proof closure consists of $2924$ imported OAI modules.  There are $2926$
downloaded source files in the saved subtree, two of which are outside this
ordinary-zeta dependency path.  Their saved SHA-256 values and Git blob
identities are recorded separately.  A static inspection of the closure
finds neither unproved declarations nor a replacement of the target zeta
function, real numbers or complex numbers.  Static inspection is useful for
identifying the claimed statement; it does not replace kernel checking.

The separately pinned report \cite{ExternalVerifier} records a complete
Lean~4.34.1 build of the $2924$ OAI modules and the independently retyped
ordinary-zeta challenge.  Its retained \texttt{\#print axioms} output lists
only
\[
 \texttt{propext},\qquad \texttt{Classical.choice},\qquad
 \texttt{Quot.sound}.
\]
The report used the pinned official Mathlib compiled cache; it did not
re-elaborate every Mathlib declaration.  Its additional \texttt{leanchecker}
replay completed for the final two modules.  The optional additional replay
of all $2924$ modules was stopped and is not reported as complete.

This manuscript imports \eqref{ver:zeta-input} as that identified external
theorem.  The work reported here did not perform a fresh full local Lean
build.  Neither the custom source estimates nor the zero-dynamics assembly
nor the final $\Lambda$ inequality has been formalised end to end in Lean.
The distinction is a verification boundary, not an undisclosed RH hypothesis
in \eqref{ver:zeta-input}.  A reader wishing to replace the external build
evidence by a new kernel run should build the pinned target
\texttt{OAI.NumberTheory.DirichletL.Nonvanishing}, then run the retained
root-qualified \texttt{RootZetaChallenge.lean} and save both its elaborated
type and axiom output.

\subsection{Finite mathematical obligations}\label{ver:obligations}

Every row of the source plan binds one functional, one full contact-height
interval, one thermal interval and one own-pair expression.  The certificate
interface requires all of the following.
\begin{enumerate}
\item The signed boundary-kernel defect is nonnegative for every real
spacing throughout the declared closed contact interval.  Its denominators
are positive, its far-field coefficient has the required sign, and the whole
signed own-pair function has the stated monotonicity.  Rational polynomial
identities and Bernstein/conic certificates certify a continuum, rather than
a set of sampled spacings.
\item The physical source lower bound applies to that same functional, at
every admitted exterior physical center and every time of its closed thermal
cell.  Its complete native labels, analytic tails, common normalisation,
translated cutoffs and physical-error charges are retained.  A source refund
is applied only to the reserved variable charge it actually bounds.
\item The wall certificate applies to the complete fixed wall in time and
height.  Its finite prefix computation is combined with Abel summation and
the physical-error transfer proved in the text.  In particular, the usable
wall floor is the corrected $697/900$, rather than the superseded narrower
historical error field.
\item The flow consumer checks each arc's source debit, full signed own-pair
upper bound, strict comparison rate, complete domain caps, positive ideal
arc, outward endpoint and adjoining state and time.  Startup and final
global strip contraction are checked separately.
\end{enumerate}
The words ``complete'' and ``same'' in these obligations are essential.
For example, a proof for a single physical phase, an individual unsigned
probe, or an initial state outside a certified kernel cap would not satisfy
this interface.  The local confluent-current experiments are not used in
any of the four obligations or in the theorem.

The three operational input identities are
\begin{center}\footnotesize
\begin{tabular}{p{.27\linewidth}p{.66\linewidth}}
\hline
Object & SHA-256\\\hline
Source ledger &
\texttt{\seqsplit{b59905d7b994d3424ba14f9347be75a2429ae4d4b0de2c2404971ffda1d4d536}}\\
Kernel plan &
\texttt{\seqsplit{ec346b7c1136eda38021eb6323033b3b48ed21e6f7458eeeac7cc9579db2497a}}\\
Complete flow &
\texttt{\seqsplit{a664e829a3681d96e474bcf875d80cba7b16fe8f65e050c69f01b5ea07de4737}}\\\hline
\end{tabular}
\end{center}
The complete bundle inventory supplies byte counts and SHA-256 digests
for the actual input files, checker programs and mathematical explanations.
Hashes identify the inputs; the consumers re-evaluate their inequalities.

\subsection{Independent finite consumers}\label{ver:executions}

The finite gates were reimplemented from the actual coefficients and
primitive data. Historical success fields and proposed critical roots
are not acceptance premises. All consumers use the same trusted
FLINT/Arb arithmetic implementation for their transcendental enclosures;
the distinct implementations concern the mathematical reductions and
finite algorithms. The principal gates are as follows.
\begin{center}\small
\begin{tabularx}{\textwidth}{@{}p{.24\textwidth}X@{}}\toprule
Gate & Continuous or finite claim checked\\\midrule
Native coefficients & All $469$ heights; full support $1\leq n\leq690988$;
$27\,484$ independently generated bins and $1658$ exact cutoffs.\\
Global activation & Both native branches on the complete half-line;
unique critical-point enclosures or exact limiting suprema.\\
Spectral envelopes & All $67$ rows on $\lambda\geq\log4096$, including
$14\,727$ closed interval pieces and the unbounded tails.\\
Arithmetic head & All $4095$ finite labels; $28$ shared-resource rows
and $39$ retained rows; all allocation, phase and Gram inequalities.\\
Physical source & All $469$ quotient caps and all $67$ row refunds; actual common
model; three complete derivative circles; all $67$ operational rows.\\
Kernel & All $67$ rational boundary functions and whole own-pair
derivatives; complete subdivision coverage and $2\,716\,496$
positive terminal integer coefficients.\\
Wall & Direct $512$-bit cosine sums for all $2\,763\,952$ prefixes,
followed by the corrected broad physical transfer.\\
Exact flow & Every source, signed-own, rate, positivity, outgoing and
cap inequality on all $221\,010$ arcs; strict startup and full cleanup.\\
Interfaces & Same actual functional, time/contact domains and source
fields at every source-to-kernel-to-flow join.\\
Direct reconstruction & All $67$ actual main bounds; the complete scalar and jet functional identities; every primitive debit and arithmetic payment; all $268$ final source-field identities.\\\bottomrule
\end{tabularx}
\end{center}

The actual source table is
\path{verification/source_overlays/ALL67_OPERATIONAL_SOURCE_TABLE.json}.
It includes every signed rational scalar node, conic derivative component,
absolute thermal/contact domain and the four operational source fields.
The polynomial data are the $67$ files
\path{verification/kernel/ROW_000_ACTUAL_CERTIFICATE.json.gz} through
\path{verification/kernel/ROW_066_ACTUAL_CERTIFICATE.json.gz}.
Each contains the full initial integer tensor and its certified closed
leaf boxes, rather than only a reported number of passed tests.
The exact arc data and the all-arc margin transcript accompany them.

Rational checks use arbitrary-precision integer fractions.
Transcendental comparisons accept only when the entire outward-rounded
ball lies on the required side of the rational threshold. Midpoint
agreement or an inconclusive overlapping interval is insufficient.
Consumers that use Python assertions reject optimized execution; they
must be run without \texttt{-O} or \texttt{-OO}.

\subsection{Replay and mathematical review}\label{ver:replay}

The bundle supplies editable and flattened manuscript sources, actual
primitive and operational inputs, continuous-domain certificates,
new consumers, source explanations and the imported ordinary-theorem
provenance. The custom certificate data and their consumers are included
in the accompanying archive, together with the retained ordinary-theorem
source and attributed verification evidence. The reference-location ledger
supplies exact URLs and hashes for contextual source articles and external
report files that are not redistributed. Its top-level replay guide specifies installation and the
commands for an extracted bundle. The raw input layout is preserved
under \path{inputs/workspace/}; portability wrappers are separate from the
immutable data. No consumer may silently read an original research
workspace that is absent from the bundle.

The final stage is \path{direct_source/verify_direct_source.py}.
It evaluates Theorem~\ref{thm:direct67-reconstruction} from the primitive
data, verifies all main-term and derivative-operator bindings, and checks
complete endpoint/critical-value coverage for the order-eight source.
The fifteen stages run in four dependency groups; their exact names and
execution results appear in the matching verification guide and receipt.

The authoritative full replay command, run from the extracted root, is
\begin{quote}\small\ttfamily
python3.12 replay/run.py --all --jobs 1 --work ../lambda-review-replay
\end{quote}
The custom finite proof can be replayed without network access on the
supported platform using the bundled runtime. The included
python-flint~0.8.0/FLINT~3.3.1 runtime supports
Linux x86-64 with CPython~3.12; the system interpreter and its platform
libraries remain software dependencies. Other platforms can install
the pinned package as specified in \path{replay/README.txt}.
Writable output is directed to the chosen external work directory;
the bundled scientific checker files and raw inputs remain unchanged.

A referee should proceed in dependency order:
\begin{enumerate}
\item Verify the inventory and the exact terminal arithmetic. Establish
that the declared source and kernel rows, their thermal/contact domains,
and the actual flow bytes agree with the inputs just identified.
\item Inspect the precise imported ordinary theorem in
\S\ref{ver:formal}. A fresh formal build is distinct from accepting the
retained external theorem and its authenticated build evidence.
\item Replay the primitive, activation, arithmetic, physical and spectral
source gates; then the $67$ continuous kernel certificates and full wall.
All source debits and credits must be derived in the same functional
units, including the near/far joins and every unused capacity reserve.
\item Traverse every exact comparison arc and require strict startup,
complete ideal-arc coverage, outward resets and the entire terminal
$q_T/2$ cleanup. A metadata-only or endpoint-only run does not establish
this step.
\item Review the analytic implications linking the finite inequalities
to all actual complex zeros: finite-head continuation, harmonic kernel
extension, compact positive-time contact and repeated-zero splitting.
The finite arithmetic does not replace these arguments.
\end{enumerate}
The ordinary formal theorem, analytic implications, arithmetic software
and finite input checks are separate trust boundaries. The packet is a
self-contained record of the present analytic and computational proof,
with its software requirements disclosed; it is not an offline Lean
kernel build of all imported libraries.

\subsection{Status of review}\label{ver:status}

The proof packet is organised for external mathematical review.  The
separate producer and consumer inspections described above were conducted
within the same research session; they are not journal peer review and do
not constitute independent public recognition of a record.  The theorem
claims the stated numerical bound at ordinary cited-theorem and
computer-assisted-proof trust.  It claims neither $\Lam<1/10$, the Riemann
hypothesis, the Lindelof hypothesis, nor a complete formalisation of the
assembled argument.

% END INCLUDED FILE


% BEGIN INCLUDED FILE: sections/conclusion.tex
\section{The assembled proof and its scope}\label{end:assembly}
The ordinary $7/8$ theorem and the functional equation initialize the
strict barrier \eqref{intro:q0}. The finite wall and finite-height zero
theorem localize every possible nonreal contact to the exterior source
region. The signed continuous-domain kernels and the corrected physical
sources imply the comparison inequality on each closed cell. The
positive-time compactness and repeated-zero arguments make the
first-contact comparison global. The rational consumer verifies every
arc and every reset, giving the global strip
$|\Imm z|\le\sqrt{Q_T}$ at the time $T$ in
\eqref{intro:terminal}. De Bruijn contraction then makes all zeros real
by $T+Q_T/2=B_*$. This proves the upper bound of
Theorem~\ref{intro:theorem}; Rodgers--Tao supplies its lower bound.

The proof is a conventional analytic proof with finite exact and
interval-arithmetic certificates, using an externally verified formal
theorem as an input. It is not an end-to-end Lean proof. The machine
checks certify the finite inequalities and joins described here; they
do not replace the analytic implications linking those inequalities
to all complex zeros. Source hashes identify the reviewed versions;
hash agreement alone is not mathematical verification.

\paragraph{Provenance and assistance.}
This manuscript assembles the author's retained research, public source
materials, and the exact-version proof records. Its preparation,
code inspection and internal cross-checking used AI assistants. Such
internal checking is distinguished from an independent external
referee's judgment. The full source packet is supplied so that the
analytic derivations and computational boundaries can be assessed
directly.

% END INCLUDED FILE


% BEGIN INCLUDED FILE: bibliography.tex


% END INCLUDED FILE


\clearpage\appendix
\section{Sharp all-frequency activation}\label{verification:sharp-activation}
The global source cap requires a joint bound for both activated native
branches, retaining their common activation length until after integration.
This appendix proves the $469$ actual activation bounds used by all $67$
thermal rows. It is the quantitative justification of the constants in
\texttt{ACTIVATION.json}.

\paragraph{Parameters and the original densities.}
Let $N_0=690988$, $v_0=\log N_0$, and let $[t_-,t_+]$ be one of the
actual closed time rows. Let $h>1$ be one of its seven native heights.
Write $L=2(v_0+s)$, $s\ge0$, and $\log n=v_0+w$ for an activated
integer. The literal native cutoff gives $0<w\le s+\delta$, where
$\delta=10^{-12}$. Put
\begin{gather*}
 V_-=\frac{2689}{200},\qquad V_+=\frac{6723}{500},\qquad
 a=\frac{t_-V_-}{2},\\
 d=\frac{h-1}{2},\quad p=\frac{h+1}{2},\quad k=a+d,\quad c=a+h.
\end{gather*}
Thus $p-d=1$, $p+d=h$, $V_-<v_0<V_+$, and $d,p,k,c>0$.

After the counting measure $\dd n=e^{v_0+w}\dd w$ is inserted, the
first-branch amplitude is exactly the exponential of
\[
 -\frac{t v_0^2}{4}+\frac{1-h}{2}v_0
 -\frac{t v_0s}{2}+\frac{1-h}{2}w
 +\frac t4(w^2-2sw).
\]
The reflected amplitude has the same first two terms and last term,
but its $s,w$ part is $-(t v_0/2+h)s+pw$. On the activated support,
$w^2-2sw\le\delta^2$. Their height-derivative frequencies are
$(v_0+w)/2$ and $(v_0+2s-w)/2$.

The adjacent-integer and endpoint correction is bounded by
\begin{equation}\label{eq:rho}
 \rho=\frac{t_+\delta^2}{4}
 +\left(\frac{t_+v_0}{2}+h\right)\delta
 +\frac{h-1}{2N_0}+\frac{t_+\delta}{2N_0}.
\end{equation}
These are the complete positive-density bounds for the finite activated
sum. They contain both branches, including the reflected derivative
frequency; no positive-time infinite Dirichlet series is used.

Let $u=s+\delta$. Enlarging the derivative frequency to
$(V_++2u-w)/2$ is safe. Replacing $s$ by $u$ in the exponent costs at
most $\exp[(t_+v_0/2+h)\delta]$. The directed consumer verifies
\begin{equation}\label{eq:reserve}
 \rho+\left(\frac{t_+v_0}{2}+h\right)\delta<\frac1{5000}
\end{equation}
for every actual row and height. Since the prefactor decreases with
$v_0$ and $t$, the single constant
\begin{equation}\label{eq:C}
 C=\exp\left(-\frac{t_-V_-^2}{4}-dV_-+\frac1{5000}\right)
\end{equation}
therefore pays all the corrections. In what follows the variable $u$
is renamed $s$ and is allowed to range over the larger half-line
$[0,\infty)$.

Define $I_j(z,s)=\int_0^s w^j e^{zw}\,\dd w$. The complete activation
amplitude and its absolute height derivative are at most $C A(s)$ and
$C D(s)$, where
\begin{align}
 A(s)&=e^{-as}I_0(-d,s)+e^{-cs}I_0(p,s),\label{eq:Aintegral}\\
 D(s)&=\frac12e^{-as}\{V_+I_0(-d,s)+I_1(-d,s)\}\notag\\
 &\quad+\frac12e^{-cs}\{(V_++2s)I_0(p,s)-I_1(p,s)\}.
 \label{eq:Dintegral}
\end{align}
Each integrand in these expressions is nonnegative on its declared
support. In particular the reflected derivative factor is retained.

\begin{lemma}[Three-exponential reduction]\label{lem:three}
Put $U=1/d-1/p=1/(dp)$ and
\[
 A_a=\frac{V_+}{2d}+\frac1{2d^2},\qquad
 A_k=\frac{V_+U}{2}+\frac12\left(\frac1{d^2}-\frac1{p^2}\right),
 \qquad B_k=\frac U2,
\]
\[
 A_c=\frac{V_+}{2p}+\frac1{2p^2},\qquad B_c=\frac1p.
\]
Then, exactly for $s\ge0$,
\begin{align}
 A(s)&=\frac1d e^{-as}-Ue^{-ks}-\frac1p e^{-cs},\label{eq:Aclosed}\\
 D(s)&=A_a e^{-as}-(A_k+B_k s)e^{-ks}
                         -(A_c+B_c s)e^{-cs}.
 \label{eq:Dclosed}
\end{align}
\end{lemma}
\begin{proof}
Use $I_0(z,s)=(e^{zs}-1)/z$ and
$I_1(z,s)=[e^{zs}(zs-1)+1]/z^2$, then collect the three exponentials.
Here $p-d=1$ and $p+d=h$, so all displayed coefficients follow by
rational algebra. No numerical approximation enters this identity.
\end{proof}

\begin{proposition}[The global amplitude maximum]\label{prop:A}
If $a=0$, then $A$ is strictly increasing and
$\sup_{s\ge0}A(s)=1/d$. If $a>0$, $A$ has exactly one positive
critical point and its global maximum occurs there.
\end{proposition}
\begin{proof}
By \eqref{eq:Aclosed},
\[
 e^{as}A'(s)=-\frac ad+kUe^{-ds}+\frac cp e^{-hs}.
\]
The last two terms are positive and strictly decreasing. At $s=0$
their sum with the first term is $2$, either by direct algebra or by
differentiating \eqref{eq:Aintegral}. For $a>0$ the limit is $-a/d<0$.
Thus there is exactly one sign change, from positive to negative.
For $a=0$ the derivative is positive everywhere, and the limit in
\eqref{eq:Aclosed} is $1/d$.
\end{proof}

\begin{proposition}[The global derivative maximum]\label{prop:D}
Let
\[
 H_k=kA_k-B_k,\quad J_k=kB_k,\qquad
 H_c=cA_c-B_c,\quad J_c=cB_c.
\]
For each actual native height the following rational inequalities hold:
\begin{equation}\label{eq:guards}
 H_k,J_k,H_c,J_c>0,\qquad
 J_k-dH_k<0,\qquad J_c-hH_c<0.
\end{equation}
Under these inequalities, if $a=0$ then $D$ is strictly increasing with
$\sup_{s\ge0}D(s)=A_a$. If $a>0$, it has exactly one positive critical
point, which is its global maximum.
\end{proposition}
\begin{proof}
Differentiating \eqref{eq:Dclosed} gives
\[
 e^{as}D'(s)=-aA_a+(H_k+J_k s)e^{-ds}
                         +(H_c+J_c s)e^{-hs}.
\]
The two nonconstant terms are positive. Their derivatives have signs
of $J_k-dH_k-dJ_k s$ and $J_c-hH_c-hJ_c s$, so
\eqref{eq:guards} makes both strictly decreasing throughout $s\ge0$.
At zero the expression equals $V_+>0$, directly from
\eqref{eq:Dintegral}. Its limit is $-aA_a$.
For $a>0$ this proves a unique maximum. For $a=0$ it proves strict
increase, with limiting value $A_a$ by \eqref{eq:Dclosed}.
\end{proof}

\paragraph{The actual directed inequalities.}
The rational parameters in \eqref{eq:guards} are checked exactly for
all $469$ heights. For each positive $a$, the new checker independently
starts with $[0,1]$, doubles its right endpoint until the factored
derivative is negative, and performs $110$ rational bisections. Both
endpoint signs are checked by outward Arb balls. The unique critical
point is therefore contained in the resulting rational interval.
Evaluation of \eqref{eq:Aclosed} or \eqref{eq:Dclosed} on that entire
interval bounds the global maximum. The $a=0$ case uses the exact
limiting constants proved above. This is a proof on the complete
unbounded activation domain, not a sampled search.

The final directed comparisons are, separately at every row and height,
\[
 C\le C_{\rm stored},\qquad
 C\sup_{s\ge0}A(s)\le A_{{\rm act},{\rm stored}},\qquad
 C\sup_{s\ge0}D(s)\le D_{{\rm act},{\rm stored}}.
\]
The exact lower and upper rational endpoints, stored constants,
critical-root brackets and source-time ownership are in
\path{verification/source_primitives/ACTIVATION_PRIMITIVES_FRESH_384.json}.
The consumer is \path{verification/source_primitives/check_activation_primitives.py}. All $67$ rows and $469$ heights
pass at $384$ Arb bits. The saved historical critical-root proposals and
success fields are not used to establish any inequality.

At the diagnostic corner $t_-=0,h=2$, $d=1/2$ and $p=3/2$, so
\[
 A(s)=2-\frac43e^{-s/2}-\frac23e^{-2s}<2,
 \qquad \sup_{s\ge0}D(s)=V_++2.
\]
This explicitly proves the actual $2C$ bound and explains why summing
separate pointwise suprema into $4C$ would lose the required constant.

\paragraph{How these bounds enter the global source.}
Let $P,P_1$ be the independently verified Euler absolute norm and its
logarithmic moment. The complete added value and frequency moment are
\[
 \bar a=P A_{{\rm act},{\rm stored}},\qquad
 \bar b=2P D_{{\rm act},{\rm stored}}
                              +P_1 A_{{\rm act},{\rm stored}}.
\]
The $P_1$ term is the Euler-factor derivative and must be included.
These are exactly the activation fields consumed by the affine near/far
source and physical-refund constructions. A separate independent source
composition proves their quotient gaps, rational cap and near/far joins.
The remaining source composition uses these actual constants; it is
checked separately from this activation gate.

\paragraph{Verification boundary.}
This is a mathematical derivation and directed computation
of the stated activation inequalities. It uses python-flint/Arb for
outward transcendental arithmetic, not a second implementation of that
arithmetic library. It does not by itself prove the heat approximation,
the finite source capacities, kernel inequality, ordinary $7/8$ theorem,
or final comparison; those have separate proof and verification gates.


\clearpage
% New reconstruction for the available 14-stage scientific source.
% Insert immediately before the existing physical-debit appendix so that
% the present theorem receives the number B.1; use labels for references.
\clearpage
\section{Direct reconstruction of the 67 physical source chords}
\label{app:direct-source-reconstruction}

This appendix assembles the operational source directly from the primitive
native moments, finite coefficient capacities, arithmetic certificates,
physical error refunds and derivative functionals.  A previously assembled
source table is not an additional source theorem.  In particular, the
validity of \path{WEIGHTED_SOURCE.json} is not assumed in this argument.
The four rational fields in the final table are the output of the formulas
below, whose ingredients have the analytic meanings established in
Sections~\ref{src:common}--\ref{src:jets} and
Appendix~\ref{sec:arithmetic-credit-audit}.

For row $i$, write its exact closed thermal interval as
$I_i=[t_i^-,t_i^+]$, and put
\[
 \theta_i(t)=\frac{t-t_i^-}{t_i^+-t_i^-},\qquad
 c_0=\frac12\log N_0,\quad N_0=690988,\qquad
 c=\frac14\log\frac{x}{4\pi}.
\]
Let $\mathcal L_i^{\rm sc}$ be the scalar functional specified by the
complete signed list $(h_{ij},a_{ij},w_{ij})$, with symmetric translations
interpreted as in \eqref{src:functional-def}.  Define
\[
 m_i^{\rm mass}=\sum_jw_{ij}>0,\qquad
 F_i(n)=\sum_jw_{ij}n^{-h_{ij}/2}
                         \cos(a_{ij}\log n/2).
\]
The two sets of scalar rows are
\[
 \mathcal R=\{0,\ldots,26\}\cup\{52,\ldots,61\}\cup\{65,66\},
 \qquad
 \mathcal S=\{27,\ldots,51\}\cup\{62,63,64\}.
\]
Thus $|\mathcal R|=39$ and $|\mathcal S|=28$; they are disjoint and
exhaust the row set.  The derivative rows are
$\mathcal J=\{0,16,44\}$.

The following quantities are defined from their primitive certificates.
Their definitions are mathematical inequalities and rational formulas,
rather than a reading of an upstream success flag.
\begin{enumerate}[label=(\roman*)]
\item
$m_i$ is the actual rational scalar main field, checked to satisfy
\begin{equation}
 m_i\le m_i^{\rm mass}c_0.
 \label{direct67:main}
\end{equation}
This comparison concerns the value used in the final reconstruction,
not a main field in a different source record.

\item
$H_{i,-}$ and $H_{i,+}$ are complete raw finite-head upper bounds:
\begin{align}
 H_{i,-}&\ge\frac12\sum_{n=2}^{4096}
 d_{t_i^+}(n)n^{-1/2-t_i^-c_0}|F_i(n)|,\notag\\
 H_{i,+}&\ge\frac12\sum_{n=2}^{4096}
 d_{t_i^+}(n)n^{-1/2-t_i^+c_0}|F_i(n)|.
 \label{direct67:raw-heads}
\end{align}
For $i\in\mathcal S$, they are the sums of the $4095$ respective
dyadic upper capacities, denoted by $U_{i,n,-}$ and $U_{i,n,+}$.
For $i\in\mathcal R$, they are the directed
upper sums checked from the literal scaled probes and the finite
Dirichlet-log recurrence.  They contain no arithmetic payment.

\item
$D_i$ is the scalar analytic, physical and coordinate debit before the
matched physical refund.  It is reconstructed from the native moment
formulas \eqref{direct67:native-caps} below.  For $i\in\mathcal R$,
it is the intercept of the verified native affine bound; for
$i\in\mathcal S$, it is the verified global debit.

\item
$\rho_i$ is the simultaneous physical refund in
\eqref{direct67:refund}.  It refunds the variable quotient error only.
The fixed allowance in every native physical cap and the coordinate
allowance remain charged.

\item
The complete scalar arithmetic payment is
\begin{equation}
 C_i=
 \begin{cases}
 C_i^{\rm comp}+d_i,&i\in\mathcal R,\\[2pt]
 \Delta_i^{\rm fix}+\mu_i\kappa_0+C_i^{\rm orbit},&i\in\mathcal S,
 \end{cases}
 \qquad \kappa_0=\frac1{5000000}.
 \label{direct67:credit}
\end{equation}
Here $d_i$ is the \emph{total replacement} common-clock payment.  If
the earlier angular payment is $d_i^{(0)}$ and the certified additional
payment is $\gamma_i$, then $d_i=d_i^{(0)}+\gamma_i$.
The three payments in the second line use the fixed lower-capacity
allocation, the coordinate reserve, and the remaining-prime weighted
Gram certificate, respectively.  Their simultaneous validity is proved
below.

\item
For $i\in\mathcal J$, let the actual derivative functional have the
coefficient vector and centre $(b_i,p_i)$ and nonnegative scalar
multiplier $\alpha_i$ recorded in the final plan.  Write its independently
charged source fields as
\[
 (m_i^{\rm j},D_i^{\rm j},H_{i,-}^{\rm j},H_{i,+}^{\rm j}).
\]
They bound $T_{b_i}S_t(x,p_i)$ in the units of
\eqref{src:jet-def}, before multiplication by $2$.  The orders are
$4,4,8$ for rows $0,16,44$; the order-eight row-$44$ certificate is used
once.  For $i\notin\mathcal J$, set $\alpha_i=1$ and set all four jet
fields equal to zero.
\end{enumerate}

\begin{theorem}[Direct reconstruction of all physical source chords]
\label{thm:direct67-reconstruction}
For every row $i$, define four rational numbers by
\begin{align}
 M_i&=\alpha_i m_i+2m_i^{\rm j},\label{direct67:M}\\
 N_i&=\alpha_i(D_i-\rho_i)+2D_i^{\rm j},\label{direct67:N}\\
 A_i&=\alpha_i(H_{i,-}-C_i)+2H_{i,-}^{\rm j},\label{direct67:A}\\
 B_i&=\alpha_i(H_{i,+}-C_i)+2H_{i,+}^{\rm j}.\label{direct67:B}
\end{align}
With the actual signed functional
\[
 J_i(t,x)=\alpha_i\mathcal L_i^{\rm sc}S_t(x)
                      +2T_{b_i}S_t(x,p_i),
\]
where the second term is absent outside $\mathcal J$, one has
\begin{equation}
 J_i(t,x)\ge
 M_i-N_i-(1-\theta_i(t))A_i-\theta_i(t)B_i
 \quad(t\in I_i,\ x\ge6000000185827).
 \label{direct67:chord}
\end{equation}
The numbers \eqref{direct67:M}--\eqref{direct67:B} are exactly the
\texttt{main\_lo}, \texttt{native\_debit}, \texttt{first\_hi} and
\texttt{last\_hi} fields, respectively, in the supplied final
\path{MAIN_CLOCK_SOURCE.json}.  The complete signed nodes, derivative
operators, thermal intervals and contact intervals used for this
identification agree with the final kernel and arc plan.
\end{theorem}

\begin{proof}
\emph{Step 1: identify the finite source and reconstruct its debit.}
Fix an actual $x,t$ in the stated domain and freeze the common finite
model \eqref{src:common-model} at that centre.  The exact coefficient
identity \eqref{src:positive-residual}, the finite Euler multiplier and
its restored prime-power terms identify the finite head with
\[
 P_{i,t,c}(\chi)=\frac12\sum_{n=2}^{4096}
 d_t(n)n^{-1/2-tc}F_i(n)\chi_n,
 \qquad |\chi_p|=1,\quad \chi_{mn}=\chi_m\chi_n.
\]
The physical tuple is used only after the inequalities for this finite
polynomial have been established.  The leading normaliser contributes
$m_i^{\rm mass}c$; its discrepancy is charged in the coordinate allowance.
The finite polynomial and the error measure come from the same common
model, so an independently phased or differently truncated model is
not substituted.

To give the debit explicitly, at a native height write the verified
primitive bounds as $A_h,B_h,P_h,P_{1,h},D_h^+,D_h^-,Z_h^-,T_h$.
For the licensed activation amplitude and half-frequency moment
$A_h^{\rm act},D_h^{\rm act}$, put
\[
 a_h=P_hA_h^{\rm act},\quad
 b_h=2P_hD_h^{\rm act}+P_{1,h}A_h^{\rm act},\quad
 g_h=1-A_h-a_h,
\]
\[
 R_h=B_h-(1-A_h)D_h^--2Z_h^-,\qquad
 r_{0,h}=\frac{P_h}{2000000},\quad
 r_{1,h}=\frac{P_h}{20000}+\frac{P_{1,h}}{4000000}.
\]
The primitive consumers check $R_h\ge0$ and $g_h>r_{0,h}>0$ before
any division.  The exact charged caps are
\begin{align}
 E_{{\rm an},h}
 &=\frac{R_h+b_h+a_hD_h^+}{2g_h}+\frac{T_h}{2},\notag\\
 E_{{\rm ph},h}
 &=\frac{r_{1,h}}{g_h-r_{0,h}}
   +\frac{(B_h+b_h)r_{0,h}}{2g_h(g_h-r_{0,h})}
   +2\cdot10^{-12}.
 \label{direct67:native-caps}
\end{align}
The shared spectral envelope is applicable to the analytic caps by
\eqref{src:positive-measure}--\eqref{src:spectral-envelope}.
Physical caps are combined with absolute node weights.  Thus the
shared-row global debit is exactly
\[
 D_i=\sum_l c_{i,l}E_{{\rm an},h_l}^{\rm global}
      +\sum_j|w_{ij}|E_{{\rm ph},h_{ij}}^{\rm global}
      +\frac{\sum_j|w_{ij}|}{125000000000}
       \quad(i\in\mathcal S).
\]
The $c_{i,l}\ge0$ are the coefficients of the verified full-half-line
envelope of the actual scalar filter.

For a retained row, the prescribed fixed near branch combines either
the absolute analytic caps or its licensed common spectral envelope,
and adds the absolute physical caps and the coordinate allowance
$10^{-9}$.  Denote this combined expression by $\mathcal N_i(s)$,
where $s=(L-L_0)/2=2(c-c_0)$, and set $D_i=\mathcal N_i(0)$.
The verified near/far inequalities and convexity in
\eqref{src:affine-native} give
\[
 \mathcal N_{i,\rm actual}(x,t)
 \le D_i+\beta_i m_i^{\rm mass}(c-c_0),\qquad 0<\beta_i<1.
\]
The remaining main slope is therefore
$A_i^{\rm rem}=(1-\beta_i)m_i^{\rm mass}>0$.
This reconstructs the actual intercept and slope from native moments;
it does not infer them from an assembled source record.

The separate coordinate fees have an explicit uniform interpretation.
Write $\zeta=h-iv$, $x'=x+v$, $L'=\log(x'/(4\pi))$ and
$s_*=(1+h)/2-ix'/2$.  The full padded native disks satisfy
$0\le h\le6.01$ and $|v|<2$; the jet disks satisfy the smaller
$0\le h\le6$ and $|v|\le39/20$.  Hence $|\Re s_*|\le4$ on
all of these complex disks.  For $s_*=u-ix'/2$, the chosen logarithm
has the exact decomposition
\[
 \Log\frac{s_*}{2\pi}
 =L'-\frac{i\pi}{2}+\Log\left(1+\frac{2iu}{x'}\right).
\]
The inequality $|\Log(1+w)|\le |w|/(1-|w|)$ for $|w|<1$,
and the displayed formula for $\alpha$, give
\[
 |\alpha(s_*)-(L'/2-i\pi/4)|
 \le\frac3{x'}+\frac4{x'-8}<\frac{10}{x'},\qquad
 |\alpha'(s_*)|\le\frac1{x'}+\frac6{(x')^2}<\frac2{x'},
\]
and $|\alpha(s_*)|<L'/2+1$.  Since
$\partial_\zeta\log B_t=\alpha(s_*)/2+t\alpha(s_*)\alpha'(s_*)/4$,
it follows that
\begin{align}
 \left|\partial_\zeta\log B_t(x+i\zeta)
                    -(L/4-i\pi/8)\right|
 &<\frac{5+t(L'/2+1)/2}{x'}+\frac{|L'-L|}{4}\notag\\
 &\le\frac{32991}{5000(X-2)}<2\cdot10^{-12}.
 \label{direct67:normaliser}
\end{align}
For the second line, the first term decreases for $x'\ge X-2$,
$t\le19/125$, and $\log((X-2)/(4\pi))<26.90$; also
$|L'-L|\le2/(x-2)$.  This proves the bound on the entire complex
disks, not just their real-height diameters.

Assign the scalar normaliser discrepancy entirely to the separate
coordinate allowance.  Its bound is
$2\cdot10^{-12}\sum_j|w_{ij}|$.  For a shared row this is at most
the already charged $8\cdot10^{-12}\sum_j|w_{ij}|$.
For every actual retained row, exact summation gives
$\sum_j|w_{ij}|<100$, so it is less than the charged $10^{-9}$.
The direct consumer checks the required domination for each actual
functional.  The additional fixed $2\cdot10^{-12}$ in every native
physical cap remains intact; it is not spent again in this argument or
included in a refund.  For a jet, the separately charged $10^{-10}$
bound is valid on its whole disk by \eqref{direct67:normaliser}; its
height derivatives are paid by the independent jet Cauchy allowance.

\emph{Step 2: prove both complete raw head chords.}
Lemma~\ref{src:coeff} gives $0\le d_t(n)\le d_{t_i^+}(n)$ on $I_i$.
Convexity of the exponential then gives, term by term,
\[
 d_t(n)n^{-tc_0}
 \le d_{t_i^+}(n)
 \bigl((1-\theta_i)n^{-t_i^-c_0}
                      +\theta_i n^{-t_i^+c_0}\bigr).
\]
After forming the signed filter and taking its modulus at each integer,
\eqref{direct67:raw-heads} bounds the full triangle head by
$H_i(t):=(1-\theta_i)H_{i,-}+\theta_iH_{i,+}$.
No interpolation of sampled source values is involved.

\emph{Step 3: pay the scalar arithmetic improvement globally.}
For $i\in\mathcal S$, every fixed triad allocation spends at most the
verified lower coefficient capacity, with repeated labels charged twice.
Summing \eqref{eq:credit-triad-deficit} therefore gives the simultaneous
payment $\Delta_i^{\rm fix}$.  The unused prime capacities, rather than
their pre-allocation values, define the weighted Gram matrix.  Its exact
Loewner/LDL certificate and \eqref{eq:credit-cosine-envelope} give
$C_i^{\rm near}$ on the complete physical near-coordinate band.
At a fixed phase tuple, finite differentiation gives the remaining
positive coordinate margin
\[
 \mu_i=m_i^{\rm mass}
       -t_i^+\sum_{n=2}^{4096}(\log n)^+U_{i,n,-}>0.
\]
The first coordinate payment is $\mu_i\kappa_0$.  Beyond the near
band the remaining coordinate reserve gives $C_i^{\rm far}$ after
subtracting that first payment.  The verified inequality
$C_i^{\rm orbit}\le\min(C_i^{\rm near},C_i^{\rm far})$ covers the
two alternative coordinate domains.  It follows that
\[
 m_i^{\rm mass}c-\Re P_{i,t,c}(\chi_{x,t})
 \ge m_i-H_i(t)+\Delta_i^{\rm fix}
                      +\mu_i\kappa_0+C_i^{\rm orbit}.
\]
The derivative argument is first uniform in the phase tuple and only
then evaluated at $\chi_{x,t}$; it does not differentiate the oscillating
physical tuple with respect to $x$.

For $i\in\mathcal R$, the composite allocations charge prime $p$ the
full amount $r_n\Omega(n)e_p$ for every allocated composite $n$.
They give $C_i^{\rm comp}$ and leave the verified residual capacities
$\alpha_{i,p}$.  The same finite prime polynomial extends the composite
bound beyond $L=27$ by the polydisc maximum principle, as proved in
Appendix~\ref{sec:arithmetic-credit-audit}.  This spends none of
$A_i^{\rm rem}$.

For completeness, the total replacement $d_i$ is paid as follows.
Put $\Delta=c-c_0$, $\omega=4d_i/A_i^{\rm rem}<1$,
$V=\omega/(1-\omega)$, $v=e^{4\Delta}-1$, and
$q_{\rm clk}=t/(16N_0^2)$.  For each actual relation and every complete
sign branch, use its nonnegative multipliers to form
\[
 \mathcal F_i(v,q_{\rm clk})=
 \pi\sum_j\lambda_j(1+\sigma_j\delta_j(v,q_{\rm clk}))
 -\sum_p C^*(\alpha_{i,p},b_p)
 +\frac{A_i^{\rm rem}}4\log(1+v).
\]
The checked odd aliases and the principal-angle conjugate include both
endpoints and all stationary maxima.  Joint concavity makes the four
checked corner values a lower bound on the entire common-clock
rectangle.  All of them are at least $d_i$.  If $\Delta\le d_i/A_i^{\rm rem}$,
then $v\le e^\omega-1\le V$ and the near bound applies.  If
$\Delta\ge d_i/A_i^{\rm rem}$, then the global unused main term
$A_i^{\rm rem}\Delta$ itself pays $d_i$.  These are alternative
domains, so the same main reserve is spent once.  After the native
affine debit has been subtracted, the result is
\[
 \mathcal L_i^{\rm sc}S_t(x)
 \ge m_i-D_i-H_i(t)+C_i^{\rm comp}+d_i
       \quad(i\in\mathcal R),
\]
before the refund.  This proves the first line of
\eqref{direct67:credit} from the total payment.  Merely adding a new
payment to the old source inequality would not justify this step.

\emph{Step 4: make the physical refund without changing the paid slope.}
At every scalar probe height the matched common-model theorem supplies
one value/derivative error pair bounded by
$\kappa_{i,h}(e_0,e_1)$, with $0\le\kappa_{i,h}\le1$.
The quotient identity \eqref{debit:quotient} shows that its variable
physical cap scales by at most $\kappa_{i,h}$.  Accordingly
\begin{equation}
 \rho_i=\sum_j|w_{ij}|(1-\kappa_{i,h_{ij}})
           \bigl(E_{{\rm ph},h_{ij}}^{\rm paid}-2\cdot10^{-12}\bigr).
 \label{direct67:refund}
\end{equation}
For a shared row the paid cap is global.  For a retained row it is the
anchor cap.  The near physical cap increases with $s$, and the global
physical cap dominates the anchor, so this anchor refund is valid on
both branches of the already proved affine inequality.  In particular,
neither $\beta_i$ nor its remaining main slope changes.  Combining
Steps 1--4 gives the direct scalar bound
\begin{equation}
 \mathcal L_i^{\rm sc}S_t(x)
 \ge m_i-(D_i-\rho_i)-H_i(t)+C_i
       \quad(i=0,\ldots,66).
 \label{direct67:scalar}
\end{equation}

\emph{Step 5: add the actual derivative functional.}
For each of the three derivative rows, the complete physical disk
has one analytic common model, one Euler gap and one holomorphic
logarithm of the physical/model quotient.  Equations
\eqref{src:jet-physical} and \eqref{src:jet-refined-ratio} charge all
derivatives of that one logarithm.  The positive-spectrum envelope
\eqref{src:jet-analytic}, its full endpoint and critical-point proof,
the complete finite head \eqref{src:jet-head}, and the separately paid
normaliser give
\[
 T_{b_i}S_t(x,p_i)\ge
 m_i^{\rm j}-D_i^{\rm j}
       -(1-\theta_i)H_{i,-}^{\rm j}-\theta_iH_{i,+}^{\rm j}.
\]
The downward main bound uses $b_{i,0}\ge0$.  Exact equality of the
final plan's $\alpha_i,p_i,b_i$, old scalar probes, time interval and
contact interval with the corresponding checked proposal is required
here.  Equality of a source table and a kernel table alone would not
establish that the source proof concerned the same derivative operator.

Multiply \eqref{direct67:scalar} by $\alpha_i\ge0$ and the jet bound
by $2$, then add.  The coefficients are exactly
\eqref{direct67:M}--\eqref{direct67:B}, proving
\eqref{direct67:chord}.  For row $44$, the order-eight jet is the jet
in this formula; its order-four predecessor is not added.

\emph{Step 6: identify the operational data.}
The direct rational consumer
\path{verification/direct_source/verify_direct_source.py} evaluates
\eqref{direct67:M}--\eqref{direct67:B} from the primitive inputs and
checks every one of the $4\cdot67$ identities against the actual final
source fields.  It also verifies \eqref{direct67:main} for the actual
scalar mains and checks the functional and domain identities described
in Step 5.  The primitive native, activation, spectral, head, phase and
jet consumers establish the inequalities used above; the direct
consumer establishes their assembly and data identity.  In particular,
the row-$44$ critical-value list is checked to contain its endpoint and
every isolated critical interval, so a shorter list cannot silently
omit an extremum during paired iteration.

As a useful cross-check, on a retained row the old scalar head was
$H_{i,\pm}-C_i^{\rm comp}-d_i^{(0)}$.  The direct formula therefore
changes that complete head by exactly
$-\alpha_i(d_i-d_i^{(0)})$, and leaves the main and native fields
unchanged after their own prescribed refund/conic assembly.  At row
$8$ the scalar difference is $29/10000$.  The former local native
credit in row $55$ occurs nowhere in
\eqref{direct67:M}--\eqref{direct67:B}.  These identities confirm that
no successive historical credit is being counted twice.
\end{proof}

\paragraph{Verification boundary.}
The reconstruction has three different kinds of input: the written
analytic identification of the physical model and its continuous-domain
extensions; exact or directed primitive inequalities checked by the
native, activation, spectral, affine, head, common-clock and jet
consumers; and the exact final assembly and functional identities.
A digest identifies the data, and a successful assembly check establishes
the stated identities.  Neither substitutes for the first two inputs.
The present theorem concerns the source chords.  The boundary kernels,
global contact argument, initial ordinary-zeta theorem and exact flow
remain the separate ingredients of the main theorem.

\section{The complete physical debit ledger}
\label{app:physical-debit-ledger}

This appendix specifies how the physical approximation is consumed. A
check of the rational ledger proves the advertised debit only after the
common-model error theorem has been established. In particular, neither a
saved success field nor an enclosure of a final debit proves that analytic
theorem. Throughout, all functions, cutoffs, carrier phases, and errors
refer to the same actual central pair $(x,t)$.

\subsection{From the common model to a logarithmic derivative}

Let $G(\zeta)=H_t(x+i\zeta)/B_t(x+i\zeta)$ and let $f_x$ be the
finite common model in \eqref{src:common-model}. At a native height let
$E$ denote the finite Euler multiplier, $Q=Ef_x$, and let
$\Delta=E(G-f_x)$. Use the positive Euler norm $P$ and its logarithmic
moment $P_1$. Since the Euler frequencies are $\log d/2$,
\[
 |E|\le P,\qquad |E'|\le P_1/2.
\]
Consequently a single physical error satisfying
$|G-f_x|\le e_0$ and $|G'-f_x'|\le e_1$ gives
\begin{equation}
 |\Delta|\le r_0:=Pe_0,\qquad
 |\Delta'|\le r_1:=Pe_1+P_1e_0/2.
 \label{debit:error-product}
\end{equation}
The term $P_1e_0/2$ is compulsory. It differentiates the same Euler
multiplier used in the value comparison.

Let $A,B$ be the nonconstant amplitude and frequency moment of the fixed
native labels, and let $a,b$ be their additional activation amplitude
and frequency moment after multiplication by $E$. Thus
\[
 |Q|\ge g:=1-A-a,\qquad |Q'|\le (B+b)/2.
\]
Every denominator below is used only after proving $g>r_0$. Direct
subtraction, before any estimate, gives
\[
 \frac{Q'+\Delta'}{Q+\Delta}-\frac{Q'}Q
 =\frac{\Delta'Q-Q'\Delta}{Q(Q+\Delta)}.
\]
It follows that the logarithmic-derivative discrepancy is bounded by
\begin{equation}
 \mathcal P(g,B+b;r_0,r_1)
 :=\frac{r_1}{g-r_0}
 +\frac{(B+b)r_0}{2g(g-r_0)}.
 \label{debit:quotient}
\end{equation}
The ledger stores $E_{\rm ph}=\mathcal P+2\cdot10^{-12}$; the fixed
last allowance is not part of the refundable error.
The normaliser identity
\[
 \partial_h\log H_t=\partial_h\log B_t+G'/G
\]
accounts separately for the leading main term and its coordinate
allowance. It is not permissible to apply \eqref{debit:quotient} to an
unnormalised $H_t$ and keep the normalised main term.

\subsection{The value and derivative errors have one origin}

For the licensed row $[t_-,t_+]$, let
\[
 S=(1+h-1/100)\frac{2689}{400}
       +t_-\left(\frac{2689}{400}\right)^2.
\]
The uniform physical theorem supplies
\[
 E_0=8e^{-S}+5\cdot10^{-10},\qquad
 E_1=800e^{-S}+7\cdot10^{-9}.
\]
For clarity, its four contributions to the value and derivative
estimates are as follows. These are deterministic bounds, not independent
random errors.
\begin{center}\small
\begin{tabular}{p{.41\textwidth}p{.22\textwidth}p{.23\textwidth}}
\toprule
Contribution & Value & Height derivative\\\midrule
Both complete heat remainders and the possible crossing pair
 & $<8e^{-S}$ & $<800e^{-S}$\\
All individual heated-term errors
 & $<5.5\cdot10^{-12}$ & $<5.5\cdot10^{-10}$\\
Exact local exponents to the local leading model
 & $<2\cdot10^{-10}$ & $<2.5\cdot10^{-9}$\\
Local leading model to the central common model
 & $<2.35\cdot10^{-10}$ & $<3.3\cdot10^{-9}$\\\bottomrule
\end{tabular}
\end{center}
The factor $100$ in the first two derivative entries follows from
Cauchy's formula on the radius-$1/100$ disk, applied to the discrepancy
with the \emph{fixed central cutoff}. No discontinuous moving truncation
is differentiated. The remaining derivative entries are obtained by the
product rule on their explicit finite models.

The complete heat remainder uses the corrected square-completion
constant $11/12$. In particular,
\[
 4\left(u^2+\frac{11}{12}\right)
 -(u^2+6|u|+2/3)=3(|u|-1)^2\ge0.
\]
The Gaussian is integrated on the full real line. Its separate negative
tail is bounded by $e^{-30000}/600$; neither Riemann--Siegel branch is
removed. The estimate is the all-real-abscissa calculation of the cited
Polymath expansion on $[-3,4]$, with the explicit constants and domain
guards in Proposition~\ref{src:physical}. The narrower printed numerical
simplification alone would not license the present height-six probes.

To make the two leading-model transfers auditable, put
$L'=\log(x'/(4\pi))$ and $C_\gamma=10+t(L'/2+1)$. The exact finite
exponents and the local leading exponents have discrepancies at most
\[
 \delta_s=5t/x',\quad \delta_s'=t/(2x'),\quad
 \delta_\gamma=yC_\gamma/x',\quad
 \delta_\gamma'=C_\gamma/x'.
\]
The logarithm in $\delta_\gamma$ is integrated from height zero on one
branch. Using the per-branch moments $6,22,180$, the product rule gives
\begin{align*}
 \text{value}&\le4\delta_s\,22+2\delta_\gamma\,6,
 \\
 \text{derivative}&\le
 \delta_\gamma(6L'+22)+\delta_s(22L'+360)
       +12\delta_\gamma'+88\delta_s'.
\end{align*}
For example, the first-branch derivative error is at most
$\delta_s M_2+2\delta_s'M_1$; the second branch supplies the other
terms. This explains the second moment and prevents a value estimate
from being differentiated without justification.

Finally, with
\[
 B_{\rm fr}=\frac{30+2t(L+3)}{x-2},
\]
the common-clock transfer pays $36B_{\rm fr}$ in value and
$18B_{\rm fr}(L+1)+36/(x-2)$ in derivative. The last term pays the
change in the reflected frequency. All these quantities are bounded at
the anchor and then decrease: they are positive logarithmic polynomials
of degree at most two divided by $x'$ or $x-2$, and $L'>26.89$.
This proves all-center coverage, rather than extrapolating from a
numerical evaluation at $x=X$.

\subsection{Exact primitive assembly and all future cutoffs}

Let $D^+$, $D^-$ and $Z^-$ be the matched finite-head upper norm,
lower norm and overlap lower bound, respectively. The positive residual
identity gives
\[
 R_*=B-(1-A)D^--2Z^-\ge0.
\]
For any of the licensed near or global activation bounds,
\begin{equation}
 E_{\rm an}=\frac{R_*+b+aD^+}{2g}+\frac T2,
 \quad E_{\rm ph}=\mathcal P(g,B+b;r_0,r_1)
                         +2\cdot10^{-12}.
 \label{debit:cap-assembly}
\end{equation}
Here $T$ is the complete omitted Euler-power tail. The frequency moment
must use
\[
 a=P A_{\rm act},\qquad
 b=2P D_{\rm act}+P_1A_{\rm act}.
\]
Thus the derivative of the Euler factor is present in the activation
ledger as well as in \eqref{debit:error-product}.

Near the anchor, $A_{\rm act}$ and $D_{\rm act}$ are positive
polynomials in the activation length. Their composition with positive
denominators gives increasing convex caps. Beyond the near split, the
sharp combined-branch activation maxima in the activation appendix
supply a bound for every $s\ge0$. The old termwise $4C$ estimate does
not prove the actual $2C$ table entry at $t_-=0,h=2$; the sharp
combined-branch proof is required. With that correction, the same
stored $469$ activation entries are valid on the whole half-line.

Only analytic residual errors share the positive spectral measure. A
licensed spectral envelope replaces their absolute node sum by a fixed
nonnegative combination of the seven moments of that measure.
Physical errors in \eqref{debit:quotient} always use absolute node
weights. In particular, cancellation in a signed spectral filter does
not refund a physical error with no corresponding common-spectrum
representation.

For a retained row, a fixed near branch $N(s)$ and a global constant
$G_*$ satisfy
\[
 N(R),G_*\le N(0)+\beta MR/2.
\]
Convexity for $0\le s\le R$ and the global cap for $s\ge R$ yield
\[
 N_{\rm actual}(L)\le N(0)+\beta M(L-L_0)/4.
\]
The remaining main slope is precisely $(1-\beta)M$; only that slope
is available to the later arithmetic credit.

\subsection{Refunds preserve both affine branches}

The baseline errors are $e_0=5\cdot10^{-7}$ and
$e_1=5\cdot10^{-5}$. For each actual native height choose a rational
\[
 \kappa\ge\max(E_0/e_0,E_1/e_1),\qquad 0\le\kappa\le1.
\]
Scaling the \emph{same} value and derivative errors in
\eqref{debit:quotient} gives
\[
 \mathcal P(g,B+b;\kappa r_0,\kappa r_1)
 \le\kappa\mathcal P(g,B+b;r_0,r_1).
\]
The refund per unit absolute node weight is therefore
\[
 (1-\kappa)(E_{\rm ph,old}-2\cdot10^{-12}).
\]
For retained rows this uses the anchor physical cap. The near physical
cap increases in $s$ and the global activation cap dominates the anchor,
so this one refund is valid on both branches of the existing affine
proof. Thus $\beta$ and its paid slope are unchanged. For a shared row
the refund uses its global physical cap. Fixed coordinate charges and
the $2\cdot10^{-12}$ allowance remain in the ledger.

\subsection{Joint derivatives and the final 67-row identity}

For a joint row, let $|G-f_x|/|f_x|\le\delta<1$ on its whole
licensed circle. There is one holomorphic logarithm
$\ell=\log(G/f_x)$, with $|\ell|\le\eta:=\delta/(1-\delta)$.
Cauchy's formula gives
\[
 \left|\frac1p\sum_{m=0}^M\frac{b_mp^m}{m!}
           \ell^{(m+1)}(p)\right|
 \le\frac\eta{pr}\sum_{m=0}^M(m+1)|b_m|(p/r)^m.
\]
All derivatives are of this one logarithm. The direct bound is
$\delta\le E_0P/g$. Since $e_1/e_0=100$, the baseline quotient
ledger implies $e_0P/g\le E_{\rm ph,old}/100$; this is the source
of the circle estimate. There is no additional multiplication by $P$.
The normaliser discrepancy is analytic on the same disk and is charged
separately. Row $44$ uses radius $39/20$, so its horizontal translations
remain inside the physical-error domain; radius $49/25$ would require
an additional domain proof.

Let $D_i^{\rm sc}$ be the scalar native debit after its proved refund.
For rows other than $0,16,44$ the final native debit is exactly
$D_i^{\rm sc}$. For those three rows it is
\[
 D_i^{\rm final}=\alpha_iD_i^{\rm sc}+2D_i^{\rm jet}.
\]
The order-eight row-$44$ jet replaces its order-four predecessor; those
jets are not added. The later main-clock gains change only the first
and last finite-head costs. They do not change this native debit.
The final source is consequently
\[
 \underline J_i(t)=M_i-D_i^{\rm final}
                  -(1-\theta)A_i-\theta B_i.
\]
The accompanying machine-readable debit map gives every rational
quantity in these identities for all $67$ rows, including the scalar
baseline, refund, conic multiplier, jet debit, and final operational
native debit. Their equality is checked with exact rational arithmetic.
This ledger does not infer a physical theorem from a hash or a reported
check count.


\clearpage
\section{Derivation and accounting of the arithmetic source payments}
\label{sec:arithmetic-credit-audit}

This section supplies the mathematical implications behind the arithmetic
payments in all $67$ source rows. It distinguishes the $28$ redesigned
scalar functionals from the $39$ retained scalar functionals. The later
conic derivative additions do not create a new arithmetic payment.
All lower capacities, integer allocations, relation vectors and rational
multipliers are data whose defining inequalities must be verified; an
optimizer's status or an earlier acceptance flag is not a premise.

\subsection{One finite polynomial and its coefficients}

Put $K=4096$, $N_0=690988$, $c_0=\log N_0/2$, and $L_0=4c_0$.
For the actual signed probe list $(h_j,v_j,w_j)$ define
\[
 M=\sum_jw_j>0,\qquad
 F(n)=\sum_j w_jn^{-h_j/2}\cos(v_j\log n/2).
\]
The signed sum defining $F(n)$ is formed before taking its modulus.
Let $b_t(n)=\exp(t\log^2n/4)$. Formal Dirichlet convolution defines
$d_t=(b_t\log)*b_t^{-1}$ coefficient by coefficient. Thus
\[
 d_t(n)=(\log n)b_t(n)
       -\sum_{\substack{r\mid n\\2\le r<n}}d_t(r)b_t(n/r).
\]
No infinite positive-time Dirichlet series is used. Differentiation of
this finite coefficient identity gives
\[
 \partial_td_t(n)=\frac{\log^2n}{4}d_t(n)
  +\frac{\log n}{4}\sum_{\substack{ab=n\\a,b\ge2}}d_t(a)d_t(b).
\]
Since $d_0(n)=\Lambda(n)\ge0$, induction and an integrating factor prove
nonnegativity and monotonicity on every closed thermal interval.

Write $s_n=\operatorname{sgn}F(n)$ and
\[
 a_n(t,c)=\tfrac12d_t(n)n^{-1/2-tc}|F(n)|,\qquad
 P_{t,c}(\chi)=\sum_{n=2}^{K}s_na_n(t,c)\chi_n.
\]
Here $\chi_{mn}=\chi_m\chi_n$ is the actual completely multiplicative
unit phase. The common finite-model and remainder calculation in
Section~\ref{src:common} and Step 1 of Theorem~\ref{thm:direct67-reconstruction}
reduces the source to $Mc-\Re P_{t,c}$ minus its fully paid analytic and
physical debit. This section proves the arithmetic improvements of
that finite polynomial bound. The source identification follows from
the common model and remainder formulas, without assuming the completed
uniform source theorem.

If $t\in[t_-,t_+]$ and $c=c_0$, monotonicity gives
\[
 a_n(t,c_0)\le(1-\theta)U_{n,-}+\theta U_{n,+},\quad
 \theta=\frac{t-t_-}{t_+-t_-},
\]
where the stored upward capacities dominate
\[
 \tfrac12d_{t_+}(n)n^{-1/2-t_-c_0}|F(n)|,
 \qquad
 \tfrac12d_{t_+}(n)n^{-1/2-t_+c_0}|F(n)|,
\]
respectively. This follows from convexity of $e^{-tc_0\log n}$ after
freezing the same nonnegative $d_{t_+}(n)$. Both endpoint sums include
every $2\le n\le4096$.

\subsection{The redesigned 28 rows: fixed lower capacities}

For these rows the stored lower capacities obey
\[
 0\le\ell_n\le
 \tfrac12d_{t_-}(n)n^{-1/2-t_+c_0}|F(n)|
 \le a_n(t,c_0).
\]
Every one of the $28$ actual rows selects the fixed-lower-capacity
branch. The earlier endpoint amplitude-box branch is therefore not
needed as a premise for their final support bounds.

For a frustrated edge $(m,n,mn)$, $s_ms_ns_{mn}=-1$. Put
$u=s_m\chi_m$ and $v=s_n\chi_n$, so the third gauged phase is $-uv$.
A template $(A,B,C;p)$ proves
\begin{equation}
 A(1-\Re u)+B(1-\Re v)+C(1+\Re uv)\ge p.
 \label{eq:credit-triad-deficit}
\end{equation}
The actual templates are validated by one of the following exact
identities. For $A,B,C>0$,
\[
 A\Re u+B\Re v-C\Re uv
 \le\tfrac12(AB/C+AC/B+BC/A),
\]
because the difference between the right and left sides is
\[
 \frac C2\left|\sqrt{A/B}\,u+\sqrt{B/A}\,\overline v
                         -\frac{\sqrt{AB}}C\right|^2.
\]
Alternatively, $C(1/A+1/B)\le1$ and the weighted Cauchy inequality
give the stronger deficit $2C$. For repeated labels $m=n$, the same
phase occurs twice; the exact maximum of
$(A+B)\cos u-C\cos(2u)$ is
\[
 \begin{cases}
 C+(A+B)^2/(8C),&A+B\le4C,\\
 A+B-C,&A+B\ge4C.
 \end{cases}
\]
The checker verifies that each stored $p$ is no larger than the
corresponding proved deficit. A dyadic allocation $g$ spends
$Ag,Bg,Cg$ at its three labels and supplies $pg$.

Summing allocations is legal only when their complete cost at each
label is at most $\ell_n$. Both occurrences of a repeated label are
charged. The remaining actual coefficients are nonnegative; their
individual triangle deficits are nonnegative. Consequently the
simultaneous fixed payment is
\[
 \Delta_{\rm fix}=\sum_e p_eg_e,
 \qquad \Re P_{t,c_0}\le\sum_n a_n(t,c_0)-\Delta_{\rm fix}.
\]

\paragraph{The global coordinate reserve.}
For a fixed phase tuple, finite differentiation gives
\[
 \partial_c\{Mc-\Re P_{t,c}(\chi)\}
 \ge M-t_+\sum_{n=2}^{K}(\log n)^+U_{n,-}=:\mu>0
 \quad(c\ge c_0).
\]
This does not differentiate the physical phase as $x$ varies. At each
physical point its actual tuple is inserted into a statement already
proved for every tuple. Since $c=\log(x/(4\pi))/4$,
\[
 c-c_0\ge\tfrac14(1-4\pi N_0^2/x).
\]
For $x\ge X_0=6\cdot10^{12}$ the right side exceeds
$\kappa_0=1/5000000$. The first coordinate payment is precisely
$\mu\kappa_0$; it is subsequently subtracted from every far reserve.

\paragraph{The remaining prime capacities and one Gram bound.}
After the complete fixed allocation, each used prime has a positive
unspent lower capacity $r_p$. For the actual prime gauges let
$\theta_p\in[-\pi,\pi]$ be their principal phases. A relation vector
$k_j$ is checked by multiplying its integer prime powers, checking its
gauge parity, and enclosing its actual logarithmic ratio. On the whole
band $X_0\le x\le X_1=6.6\cdot10^{12}$, with the same thermal cell,
the actual clock $\tau=x/2+\pi t/8$ then gives
\[
 |k_j\cdot\theta|\ge\rho_j>0.
\]
The clock interval is checked continuously at its extreme $x,t$ values,
not on a phase grid. Let $R=\operatorname{diag}(r_p)$, let $s_j>0$ be
the stored rational row scales, and put
\[
 V_{jp}=s_jk_{jp}/\sqrt{r_p},\quad G=VV^{\mathsf T},
 \quad z_p=\sqrt{r_p}\,\theta_p.
\]
All Gram entries are reconstructed exactly from the same unspent
capacities. The certificate rounds $G$ entrywise downward to a symmetric
$\widetilde G$ on grid $2^{-64}$. For $m$ relations,
\[
 U=\widetilde G+m2^{-64}I\succeq G.
\]
This uses the entrywise error bound and the row-sum bound on the
operator norm. Its exact certificate is
\[
 U-UWU=\mathcal L\mathcal D\mathcal L^{\mathsf T},
 \qquad W=\operatorname{diag}(w_j),\quad w_j\ge0,
 \quad\mathcal D\succ0.
\]
Every multiplication identity and every diagonal sign is checked with
fractions. Since $U=\mathcal L\mathcal D\mathcal L^{\mathsf T}+UWU$
is positive definite, this implies
$W^{1/2}UW^{1/2}\preceq I$. Therefore
$W^{1/2}GW^{1/2}\preceq I$ and $V^{\mathsf T}WV\preceq I$.
This proof also covers singular $G$; no inverse of $G$ is taken. It
follows that
\[
 \sum_pr_p\theta_p^2\ge
 E:=\sum_jw_j(s_j\rho_j)^2.
\]
The relations are one combined operator constraint, not separately
funded expenditures from reset prime capacities.

\paragraph{From angle energy to an actual source payment.}
Define $r_*=\min_pr_p$ and $u=\min(E/r_*,6)$. Since $u\leq6<\pi^2$,
the following bounds hold:
\begin{equation}
 \sum_pr_p(1-\cos\theta_p)
 \ge r_*(1-\cos\sqrt u)
 \ge r_*\left(\frac u2-\frac{u^2}{24}
                 +\frac{u^3}{720}-\frac{u^4}{40320}\right)=:C_{
 \rm near}.
 \label{eq:credit-cosine-envelope}
\end{equation}
Indeed, if some $|\theta_p|\ge\sqrt u$, that one term suffices.
Otherwise $(1-\cos v)/v^2$ decreases on $[0,\pi]$, and multiplication
by the angle-energy bound proves the first inequality. The second is
the alternating cosine bound on $0\le u\le6$.

For $x\ge X_1$, the global coordinate derivative instead supplies
\[
 C_{\rm far}=\mu\left\{\tfrac14(1-4\pi N_0^2/X_1)-\kappa_0\right\}>0.
\]
Thus the exact additional orbit payment must obey
$0<C_{\rm orbit}\le\min(C_{\rm near},C_{\rm far})$.
The near and far arguments are alternative domain bounds, so their
payments are not added. The redesigned scalar endpoints are exactly
\[
 H_\pm=\sum_nU_{n,\pm}-\Delta_{\rm fix}
                  -\mu\kappa_0-C_{\rm orbit}.
\]

\subsection{The retained 39 rows: composites and common clock}

For these rows the native proof has already charged an affine growth
cost $\beta M(c-c_0)$, with $0<\beta<1$. The main slope still available
is $A=(1-\beta)M>0$. Lower capacities at thermal endpoint $t_-$ and
attenuation $t_+27/4$ are valid on the whole resource rectangle
$t\in[t_-,t_+]$, $L_0\le L\le27$.

For a composite $n=\prod p^{e_p}$ with $\Omega(n)=\sum e_p$ and a
negative gauged coefficient, allocate $r_n$ no larger than its actual
available magnitude. Writing $z_p$ for the gauged unit prime phases,
\[
 1-\Re\prod_pz_p^{e_p}
 \le\Omega(n)\sum_pe_p(1-\Re z_p)
\]
by telescoping and Cauchy. Hence its full triangle deficit is at least
\[
 r_n\{1+\Re\prod_pz_p^{e_p}\}
 \ge2r_n-r_n\Omega(n)\sum_pe_p(1-\Re z_p).
\]
Prime $p$ is therefore charged $r_n\Omega(n)e_p$. Summing every
allocation proves the composite payment $C_{\rm comp}\le2\sum_nr_n$
while leaving the exact residual prime capacities
$\alpha_p>0$. These capacities are not restored before the angular
duals are evaluated.

\paragraph{Why the composite payment is global.}
At fixed $t$, the complete finite head is the real part of a polynomial
in its finitely many prime variables. The composite allocation gives a
torus-uniform upper bound at $L=27$. For $L>27$ the same polynomial is
evaluated at the radial contraction
\[
 z_p\longmapsto p^{-t(L-27)/4}z_p.
\]
The maximum principle in each variable preserves that upper bound
throughout the polydisc. The baseline endpoint chord at $L_0$ bounds
the triangle mass at $27$ as well. Thus the composite improvement
continues globally without assuming positive constant lower capacities
at arbitrarily large $L$, and without spending any of the main slope $A$.

\paragraph{The final total angular/main payment.}
Let $d$ be the total angular payment that will replace the old angular
payment. Put $\Delta=c-c_0$, $v=e^{4\Delta}-1$ and
$q=t/(16N_0^2)$. A checked prime relation $R_j=\prod p^{k_{jp}}>1$
has $a_j=2N_0^2\log R_j$, a checked gauged odd target $h_j$, and
\[
 \frac{k_j\cdot\theta}{\pi}=2m_j+1+\delta_j,
 \qquad\delta_j=a_j(1+v+q)-h_j.
\]
The exact source clock used here is $\tau=\pi(2e^L+t/8)$.
On the verified alias rectangle $|\delta_j|<1$, the actual sign
$\sigma_j=\operatorname{sgn}(k_j\cdot\theta)$ gives
\[
 \sigma_jk_j\cdot\theta\ge\pi(1+\sigma_j\delta_j).
\]
For each complete sign branch use its fixed nonnegative rational
multipliers $\lambda_j$, and form the combined coefficient
$b_p=\sum_j\lambda_j\sigma_jk_{jp}$. The exact conjugate is
\[
 C^*(\alpha,b)=\max_{|u|\le\pi}\{bu-\alpha(1-\cos u)\}.
\]
For $|b|<\alpha$ its maximum is the larger of
\[
 \pi|b|-2\alpha,\qquad
 \alpha\left[\frac{|b|}{\alpha}\arcsin\frac{|b|}{\alpha}
       +\sqrt{1-\frac{b^2}{\alpha^2}}-1\right];
\]
for $|b|\ge\alpha$ only the endpoint maximum is needed. In particular,
the global principal-angle endpoint is never dropped.

Choose $\omega=4d/A<1$ and $V=\omega/(1-\omega)$. On
$0\le v\le V$ and $q\in[t_-/(16N_0^2),t_+/(16N_0^2)]$, with
$L_0+\log(1+V)<27$, define
\[
 \mathcal F(v,q)=\pi\sum_j\lambda_j(1+\sigma_j\delta_j)
       -\sum_p C^*(\alpha_p,b_p)+\frac A4\log(1+v).
\]
This explicit lower payment is jointly concave, since its only nonlinear
term has second derivative $-A/[4(1+v)^2]$. Its four corner values
therefore bound it from below throughout the rectangle. All $128$
branches of every seven-relation row, and all $256$ branches of the
eight-relation row $8$, are verified. The same $v,q$ is used in all
relations before taking this bound.

If the four corner payments are at least $d$, the near interval
$0\le\Delta\le d/A$ is covered because
$e^{4\Delta}-1\le V$. On the complementary far interval,
$A\Delta\ge d$. These alternatives cover every physical center and
spend the retained main slope once. This argument asserts concavity
of $\mathcal F$, not concavity of the native current itself.

\subsection{Exact operational composition and review data}

The retained scalar endpoints are their complete raw triangle endpoints
minus $C_{\rm comp}+d$. The old angular payment $d_0$ is replaced by
$d$, so the decrement relative to the pinned M8 parent is $d-d_0$,
not another copy of $d$. For row $8$, $d-d_0=29/10000$.
The former row-$55$ local native credit is not used by the direct final
construction.

At rows $0,16,44$ the source is a conic functional
$\alpha_i$ times its complete scalar functional plus twice its separately
charged derivative functional, with $\alpha_i>0$. An arithmetic payment
belonging to the scalar component is multiplied by $\alpha_i$; it is
not assigned to the derivative component. All main, native and head
fields are combined in exactly these same units.

The companion files \path{review/arithmetic_credit/ALL67_ARITHMETIC_CREDIT_EXACT.csv}
and \path{review/arithmetic_credit/ALL67_ARITHMETIC_CREDIT_AUDIT.json} give every scalar
payment, conic factor, effective payment and final operational field as
an exact rational. Their composition script performs $650$ exact
cross-table checks, including that all $28$ redesigned rows actually use
the fixed branch and that the final MAIN decrement changes neither main
nor native debit. This composition audit does not replace the directed
capacity checks or the exact allocation/dual feasibility checks. Those
are provided by the fresh head, phase and source-overlay consumers in
the bundle. Decimal displays in the associated table are explanatory
only; the exact rational files are the review data.


\clearpage
\begingroup\small
\setlength{\tabcolsep}{4pt}
\begin{longtable}{rcrrrr}
\caption{Actual arithmetic payments and operational values. R denotes a retained scalar row; S denotes a shared-resource scalar row. Decimal displays are for orientation; the exact CSV and JSON are authoritative. $C_i$ is the total arithmetic payment in final conic units; $g_i$ is the additional MAIN payment relative to the M8 parent, already included in $C_i$.}\label{tab:credit-audit}\\
\toprule Row & Branch & $C_i$ & $g_i$ & Native debit & First head\\\midrule\endfirsthead
\toprule Row & Branch & $C_i$ & $g_i$ & Native debit & First head\\\midrule\endhead
0 & R & 0.007992754 & 0.000692071 & 0.007352276 & 0.830909465 \\
1 & R & 0.008310455 & 0.000725020 & 0.004473392 & 0.841700097 \\
2 & R & 0.008846060 & 0.000770300 & 0.005000931 & 0.860827810 \\
3 & R & 0.009463726 & 0.000825200 & 0.005536952 & 0.886329286 \\
4 & R & 0.009990737 & 0.000871920 & 0.006038487 & 0.910052483 \\
5 & R & 0.010501562 & 0.000901010 & 0.006562747 & 0.931199801 \\
6 & R & 0.011031045 & 0.000933730 & 0.007162690 & 0.954090179 \\
7 & R & 0.011500752 & 0.001028680 & 0.007488877 & 0.981247488 \\
8 & R & 0.013366970 & 0.002900000 & 0.007973102 & 0.999627216 \\
9 & R & 0.012873033 & 0.001803780 & 0.010538224 & 1.028942841 \\
10 & R & 0.013415637 & 0.001963190 & 0.010856963 & 1.059829839 \\
11 & R & 0.014507521 & 0.002086750 & 0.017702977 & 1.098120775 \\
12 & R & 0.015070784 & 0.002180770 & 0.018940933 & 1.126986096 \\
13 & R & 0.015516128 & 0.002259060 & 0.019991563 & 1.152390263 \\
14 & R & 0.016097677 & 0.002390950 & 0.021443059 & 1.186140203 \\
15 & R & 0.015868172 & 0.002476250 & 0.014065630 & 1.203282544 \\
16 & R & 0.016257247 & 0.002585420 & 0.014594437 & 1.231512530 \\
17 & R & 0.016874786 & 0.002756000 & 0.015690624 & 1.269936367 \\
18 & R & 0.017335611 & 0.002905560 & 0.016811373 & 1.305060028 \\
19 & R & 0.017754214 & 0.003051390 & 0.017389834 & 1.339761469 \\
20 & R & 0.018851274 & 0.003321610 & 0.018241455 & 1.366940774 \\
21 & R & 0.018937267 & 0.003356330 & 0.018384089 & 1.378608746 \\
22 & R & 0.019269513 & 0.003417300 & 0.019359209 & 1.395682437 \\
23 & R & 0.019602610 & 0.003492240 & 0.020547219 & 1.414513496 \\
24 & R & 0.019711483 & 0.003546330 & 0.021992492 & 1.429308430 \\
25 & R & 0.020027084 & 0.003657460 & 0.023867433 & 1.452776191 \\
26 & R & 0.020343147 & 0.003778830 & 0.026095236 & 1.477717729 \\
27 & S & 0.028242099 & 0.000000000 & 0.044129721 & 1.467643936 \\
28 & S & 0.026777735 & 0.000000000 & 0.045465678 & 1.484603418 \\
29 & S & 0.027790520 & 0.000000000 & 0.045820570 & 1.498670787 \\
30 & S & 0.028185726 & 0.000000000 & 0.046303642 & 1.511838457 \\
31 & S & 0.026407072 & 0.000000000 & 0.047195919 & 1.528413131 \\
32 & S & 0.028429880 & 0.000000000 & 0.051003043 & 1.559002153 \\
33 & S & 0.028595615 & 0.000000000 & 0.051558876 & 1.572185886 \\
34 & S & 0.028575250 & 0.000000000 & 0.052568393 & 1.586742435 \\
35 & S & 0.028771523 & 0.000000000 & 0.052868914 & 1.600817782 \\
36 & S & 0.029248344 & 0.000000000 & 0.057294700 & 1.635240981 \\
37 & S & 0.029229774 & 0.000000000 & 0.058289519 & 1.650966508 \\
38 & S & 0.028561564 & 0.000000000 & 0.058648779 & 1.665852839 \\
39 & S & 0.028473963 & 0.000000000 & 0.059021661 & 1.679704793 \\
40 & S & 0.029112812 & 0.000000000 & 0.063734400 & 1.717221416 \\
41 & S & 0.027063314 & 0.000000000 & 0.064156008 & 1.734529499 \\
42 & S & 0.027886168 & 0.000000000 & 0.065097015 & 1.747862654 \\
43 & S & 0.027395467 & 0.000000000 & 0.066042023 & 1.761963696 \\
44 & S & 0.026592799 & 0.000000000 & 0.069865179 & 1.805833736 \\
45 & S & 0.026269794 & 0.000000000 & 0.069789725 & 1.818446555 \\
46 & S & 0.025493330 & 0.000000000 & 0.071248660 & 1.833463484 \\
47 & S & 0.024785208 & 0.000000000 & 0.070803068 & 1.847371066 \\
48 & S & 0.025626963 & 0.000000000 & 0.074972823 & 1.891059020 \\
49 & S & 0.024393778 & 0.000000000 & 0.075222858 & 1.905153435 \\
50 & S & 0.023581134 & 0.000000000 & 0.075283539 & 1.920317573 \\
51 & S & 0.023000336 & 0.000000000 & 0.075384956 & 1.934066267 \\
52 & R & 0.025088829 & 0.005539620 & 0.052801917 & 1.983306678 \\
53 & R & 0.025257545 & 0.005551530 & 0.055196710 & 1.998797936 \\
54 & R & 0.025171524 & 0.005460310 & 0.056277883 & 1.999146894 \\
55 & R & 0.025408491 & 0.005538960 & 0.056820089 & 2.024425013 \\
56 & R & 0.025653990 & 0.005623970 & 0.057063553 & 2.051073524 \\
57 & R & 0.025835110 & 0.005664110 & 0.058453267 & 2.071897773 \\
58 & R & 0.025775362 & 0.005604400 & 0.059338358 & 2.078165784 \\
59 & R & 0.025996625 & 0.005645030 & 0.061728418 & 2.099323832 \\
60 & R & 0.025858250 & 0.005533830 & 0.062200549 & 2.096631518 \\
61 & R & 0.025859724 & 0.005563150 & 0.063602791 & 2.102464680 \\
62 & S & 0.016206921 & 0.000000000 & 0.082599160 & 2.188309328 \\
63 & S & 0.015637810 & 0.000000000 & 0.081251399 & 2.199655740 \\
64 & S & 0.014936423 & 0.000000000 & 0.079801036 & 2.210346408 \\
65 & R & 0.026341046 & 0.005533790 & 0.060396113 & 2.172074803 \\
66 & R & 0.026560964 & 0.005592840 & 0.058847507 & 2.199420129 \\
\bottomrule
\end{longtable}

\endgroup

\clearpage
\begin{thebibliography}{99}
\bibitem{deBruijn}
N. G. de Bruijn, \emph{The roots of trigonometric integrals},
Duke Math. J. \textbf{17} (1950), no.~3, 197--226.
\href{https://doi.org/10.1215/S0012-7094-50-01720-0}{doi:10.1215/S0012-7094-50-01720-0}.

\bibitem{Newman}
C. M. Newman, \emph{Fourier transforms with only real zeros},
Proc. Amer. Math. Soc. \textbf{61} (1976), no.~2, 245--251.
\href{https://doi.org/10.1090/S0002-9939-1976-0434982-5}{doi:10.1090/S0002-9939-1976-0434982-5}.

\bibitem{Polymath}
D. H. J. Polymath, \emph{Effective approximation of heat flow evolution
of the Riemann $\xi$ function, and a new upper bound for the
de Bruijn--Newman constant}, Research in the Mathematical Sciences
\textbf{6} (2019), article 31.
\href{https://doi.org/10.1007/s40687-019-0193-1}{doi:10.1007/s40687-019-0193-1}.
Numbered references use
\href{https://arxiv.org/abs/1904.12438v2}{arXiv:1904.12438v2}.

\bibitem{PlattTrudgian}
D. J. Platt and T. S. Trudgian,
\emph{The Riemann hypothesis is true up to $3\cdot10^{12}$},
Bull. Lond. Math. Soc. \textbf{53} (2021), 792--797.
\href{https://doi.org/10.1112/blms.12460}{doi:10.1112/blms.12460};
\href{https://arxiv.org/abs/2004.09765}{arXiv:2004.09765}.
Theorem~1 states the exact verification height $3000175332800$.

\bibitem{RodgersTao}
B. Rodgers and T. Tao,
\emph{The de Bruijn--Newman constant is non-negative},
Forum of Mathematics, Pi \textbf{8} (2020), e6, 62 pp.
\href{https://doi.org/10.1017/fmp.2020.6}{doi:10.1017/fmp.2020.6};
\href{https://arxiv.org/abs/1801.05914}{arXiv:1801.05914}.

\bibitem{SevenEighths}
OpenAI, \nolinkurl{lean/OAI/NumberTheory/DirichletL/Nonvanishing.lean},
Lean source for the ordinary-zeta declaration
\nolinkurl{OAI.riemannZeta_ne_zero_of_seven_eighths_lt_re},
in the official \nolinkurl{openai/math} repository, revision
\nolinkurl{fd4aeeb2ee4fc729c18d98444fed42fd0529eeeb}.
\url{https://github.com/openai/math/blob/fd4aeeb2ee4fc729c18d98444fed42fd0529eeeb/lean/OAI/NumberTheory/DirichletL/Nonvanishing.lean}.
The independent rebuild in \cite{ExternalVerifier} uses revision
\nolinkurl{adc7f1241b42e322a6451854ab7e4b4c146bf78a}; the relevant proof
subtree and build configuration match the retained source pin.

\bibitem{ExternalVerifier}
tomoto0, \emph{Quasi-Riemann hypothesis $7/8$ verification},
source/build report and explicit ordinary-zeta challenge,
\url{https://github.com/tomoto0/quasi-riemann-hypothesis-7-8-verification},
commit \nolinkurl{b83a73c6cb626a391e551baafdb7d0cceafaa55a}.
The external report is supplied as verification evidence; the present
work did not repeat the full Lean build locally.

\bibitem{Arb}
F. Johansson, \emph{Arb: efficient arbitrary-precision midpoint-radius
interval arithmetic}, IEEE Transactions on Computers \textbf{66}
(2017), 1281--1292.
\href{https://doi.org/10.1109/TC.2017.2690633}{doi:10.1109/TC.2017.2690633}.

\bibitem{RaymondPrevious}
D. J. Raymond, \emph{Retuned translated probes and the bound for the
de Bruijn--Newman constant}, public manuscript and proof materials,
\url{https://github.com/Dale-Dug/de-bruijn-newman-0116398},
commit \nolinkurl{4694f9ab2e8bc88aeeae1949d2e690c2360d2eba}.
The pinned repository states $\Lambda\le0.116397726408291$; its
October 26 update adds the ordinary $7/8$ formal-input evidence.

\bibitem{ProofPacket}
D. J. Raymond, \emph{Proof and replay materials for
$\Lambda\le0.107964910079$}, exact-version research and publication
packet, version 1.0, 29 October 2026.
DOI: \href{https://doi.org/10.5281/zenodo.23752103}{10.5281/zenodo.23752103}.
The accompanying manifest specifies each
source file, certificate, checker, size and SHA-256 digest. The proof
does not depend on the historical names of the research rounds.

\bibitem{RaymondPhase}
D. J. Raymond, \emph{Phase-correlated certificates and the bound
$\Lambda<0.163297$ for the de Bruijn--Newman constant},
Zenodo preprint, 8 September 2026.
\href{https://doi.org/10.5281/zenodo.22626161}{doi:10.5281/zenodo.22626161}.

\bibitem{RaymondCorrelated}
D. J. Raymond, \emph{Correlated translated probes and the bound
$\Lambda<0.149396$ for the de Bruijn--Newman constant},
Zenodo preprint, 21 September 2026.
\href{https://doi.org/10.5281/zenodo.22867547}{doi:10.5281/zenodo.22867547}.

\bibitem{RaymondRetuned}
D. J. Raymond, \emph{Retuned translated probes and the conditional bound
$\Lambda<0.116398$ for the de Bruijn--Newman constant},
Zenodo preprint, 10 October 2026.
\href{https://doi.org/10.5281/zenodo.23272725}{doi:10.5281/zenodo.23272725}.

\bibitem{HakimogluBrown2026}
Cetin Hakimoglu-Brown. \emph{A conditional bound $\Lambda\leq0.13963$ for the de Bruijn--Newman constant}.
8 October 2026.
DOI: \href{https://doi.org/10.48550/arXiv.2610.12803}{10.48550/arXiv.2610.12803}.
arXiv:2610.12803v1.

\bibitem{BellottiTrudgianYang2026}
Chiara Bellotti, Tim Trudgian, Andrew Yang. \emph{Zero-free regions inspired by work of Heath-Brown}.
23 March 2026.
DOI: \href{https://doi.org/10.48550/arXiv.2603.21490}{10.48550/arXiv.2603.21490}.
arXiv:2603.21490v1.

\bibitem{Kornaus2026}
Daniel Martin Kornaus. \emph{Collision Exclusion in the de Bruijn--Newman Heat Flow}.
8 September 2026.
DOI: \href{https://doi.org/10.5281/zenodo.22658444}{10.5281/zenodo.22658444}.
Zenodo preprint, record 22658444.

\bibitem{KiKimLee2009}
H. Ki, Y.-O. Kim and J. Lee,
\emph{On the de Bruijn--Newman constant},
Advances in Mathematics \textbf{222} (2009), no.~1, 281--306.
\href{https://doi.org/10.1016/j.aim.2009.04.003}{doi:10.1016/j.aim.2009.04.003}.

\bibitem{PlanatCollective2026}
M. Planat, \emph{Collective Contraction in the de Bruijn--Newman Heat Flow:
Off-Zero Logarithmic Derivatives and Certified Barrier Refinements},
preprint, version~2, 7 October 2026, 54 pp.; first version 29 September 2026.
\href{https://arxiv.org/abs/2609.37164v2}{arXiv:2609.37164v2};
\href{https://doi.org/10.48550/arXiv.2609.37164}{doi:10.48550/arXiv.2609.37164}.

\bibitem{PlanatClassical2026}
M. Planat, \emph{A certified classical-strip comparison for the
de Bruijn--Newman constant}, preprint, 13 October 2026, 19 pp.
\href{https://arxiv.org/abs/2610.18214v1}{arXiv:2610.18214v1};
\href{https://doi.org/10.48550/arXiv.2610.18214}{doi:10.48550/arXiv.2610.18214}.

\bibitem{RaymondAnalytic}
D. J. Raymond, \emph{Analytic foundations for the time-zero upper-bound
construction}, version~1.0, 21 September 2026, 19 pp.
Companion manuscript in the deposit \cite{RaymondCorrelated},
\href{https://doi.org/10.5281/zenodo.22867547}{doi:10.5281/zenodo.22867547};
supplied as \nolinkurl{Prior_Analytic_Foundations.pdf}.
The original public companion is also retained in
\url{https://github.com/Dale-Dug/de-bruijn-newman-0149396}.

\bibitem{XuSang2026}
Y. Xu and H. Sang, \emph{Amplified moments and a refinement of the
quasi-Riemann hypothesis}, preprint, version~2, 15 October 2026;
first version 11 October 2026.
\href{https://arxiv.org/abs/2610.15291v2}{arXiv:2610.15291v2};
\href{https://doi.org/10.48550/arXiv.2610.15291}{doi:10.48550/arXiv.2610.15291}.

\bibitem{LiuRefinement2026}
B. Liu, \emph{Improved zero-free half-planes for the quasi-Riemann
hypothesis}, preprint, version~2, 12 October 2026;
first version 8 October 2026.
\href{https://arxiv.org/abs/2610.12234v2}{arXiv:2610.12234v2};
\href{https://doi.org/10.48550/arXiv.2610.12234}{doi:10.48550/arXiv.2610.12234}.

\bibitem{LiuHecke2026}
B. Liu, \emph{Zero-free half-planes for Hecke $L$-functions over number
fields}, preprint, submitted 19 October 2026.
\href{https://arxiv.org/abs/2610.26835v1}{arXiv:2610.26835v1};
\href{https://doi.org/10.48550/arXiv.2610.26835}{doi:10.48550/arXiv.2610.26835}.

\bibitem{OAI:The-Quasi-Riemann-Hypothesis-September-30-2026}
OpenAI. \emph{The Quasi-Riemann Hypothesis: A Zero-Free Half-Plane $\mathrm{Re}(s)>7/8$}.
30 September 2026.
\url{https://github.com/openai/math/blob/main/preprints/The-Quasi-Riemann-Hypothesis-September-30-2026/paper.pdf}.

\bibitem{MI2026}
Mosaic Intelligence. \emph{A certified unconditional upper bound $\Lambda\leq0.1965$ for the de Bruijn--Newman constant}.
2026.
DOI: \href{https://doi.org/10.5281/zenodo.20724170}{10.5281/zenodo.20724170}.

\bibitem{MI2026b}
Mosaic Intelligence. \emph{A certified unconditional upper bound $\Lambda\leq0.1875$ for the de Bruijn--Newman constant}.
2026.
DOI: \href{https://doi.org/10.5281/zenodo.21175533}{10.5281/zenodo.21175533}.

\bibitem{MI2026c}
Mosaic Intelligence. \emph{A certified unconditional upper bound $\Lambda<0.16$ for the de Bruijn--Newman constant}.
2026.
\url{https://github.com/463464q435q43/optimizationproblems/tree/6baae4c15dccec1e94ac011de399ec382fcea7b6/certificates/certified0160}.
Repository revision \texttt{6baae4c15dccec1e94ac011de399ec382fcea7b6}.

\bibitem{MI2026d}
Mosaic Intelligence. \emph{A certified unconditional upper bound $\Lambda<0.1595$ for the de Bruijn--Newman constant}.
2026.
\url{https://github.com/463464q435q43/optimizationproblems/tree/cc8107854fedd5a96d7eb695178f405568712f35/certificates/certified01595}.
Repository revision \texttt{cc8107854fedd5a96d7eb695178f405568712f35}.

\bibitem{Gomila2026}
Jude Gomila. \emph{A computer-assisted proof of the bound $\Lambda\leq0.1787854$ for the de Bruijn--Newman constant}.
20 August 2026.
\url{https://github.com/judegomila/dbn-lambda-01787854-candidate-audit}.
Author's primary release manuscript; not yet peer reviewed.

\bibitem{Gordon2026}
Stefan Gordon. \emph{A computer-assisted upper bound for the de Bruijn--Newman constant}.
12 September 2026.
DOI: \href{https://doi.org/10.2139/ssrn.7608638}{10.2139/ssrn.7608638}.
Preprint; author-hosted PDF dated September 2026. SSRN posting 14 October 2026, revised 15 October 2026.

\bibitem{TaoCatalogue2026}
T. Tao and contributors, \emph{de Bruijn--Newman constant},
\emph{Optimization Problems and Constants}, maintained catalogue,
accessed 29 October 2026.
\url{https://teorth.github.io/optimizationproblems/constants/21a.html}.
\end{thebibliography}
\end{document}
